\documentclass{article}


\usepackage[table,dvipsnames]{xcolor}

\usepackage{amsmath}
\usepackage{amssymb}
\usepackage{amsthm}

\definecolor{darkgreen}{RGB}{0,100,0}
\newcommand{\amir}[1]{\textcolor{darkgreen}{Amir: #1}}

\usepackage{graphicx}
\usepackage{svg}
\usepackage{wrapfig}
\usepackage{subcaption}
\usepackage{booktabs}
\usepackage{siunitx}
\usepackage{tabularx}
\usepackage{array}

\usepackage{float}
\usepackage{placeins}

\usepackage{tikz}
\usepackage{tikz-cd}
\usetikzlibrary{positioning,calc,arrows.meta,fit}

\usepackage{multicol}

\usepackage[most]{tcolorbox}

\usepackage[numbers,sort&compress]{natbib}


\usepackage{aliascnt}

\usepackage{amsmath,amsfonts,bm,amssymb}

\newcommand{\figleft}{{\em (Left)}}
\newcommand{\figcenter}{{\em (Center)}}
\newcommand{\figright}{{\em (Right)}}
\newcommand{\figtop}{{\em (Top)}}
\newcommand{\figbottom}{{\em (Bottom)}}
\newcommand{\captiona}{{\em (a)}}
\newcommand{\captionb}{{\em (b)}}
\newcommand{\captionc}{{\em (c)}}
\newcommand{\captiond}{{\em (d)}}

\newcommand{\newterm}[1]{{\bf #1}}


\def\figref#1{figure~\ref{#1}}
\def\Figref#1{Figure~\ref{#1}}
\def\twofigref#1#2{figures \ref{#1} and \ref{#2}}
\def\quadfigref#1#2#3#4{figures \ref{#1}, \ref{#2}, \ref{#3} and \ref{#4}}
\def\secref#1{section~\ref{#1}}
\def\Secref#1{Section~\ref{#1}}
\def\twosecrefs#1#2{sections \ref{#1} and \ref{#2}}
\def\secrefs#1#2#3{sections \ref{#1}, \ref{#2} and \ref{#3}}
\def\eqref#1{equation~\ref{#1}}
\def\Eqref#1{Equation~\ref{#1}}
\def\plaineqref#1{\ref{#1}}
\def\chapref#1{chapter~\ref{#1}}
\def\Chapref#1{Chapter~\ref{#1}}
\def\rangechapref#1#2{chapters\ref{#1}--\ref{#2}}
\def\algref#1{algorithm~\ref{#1}}
\def\Algref#1{Algorithm~\ref{#1}}
\def\twoalgref#1#2{algorithms \ref{#1} and \ref{#2}}
\def\Twoalgref#1#2{Algorithms \ref{#1} and \ref{#2}}
\def\partref#1{part~\ref{#1}}
\def\Partref#1{Part~\ref{#1}}
\def\twopartref#1#2{parts \ref{#1} and \ref{#2}}

\def\ceil#1{\lceil #1 \rceil}
\def\floor#1{\lfloor #1 \rfloor}
\def\1{\bm{1}}
\newcommand{\train}{\mathcal{D}}
\newcommand{\valid}{\mathcal{D_{\mathrm{valid}}}}
\newcommand{\test}{\mathcal{D_{\mathrm{test}}}}

\def\eps{{\epsilon}}


\def\reta{{\textnormal{$\eta$}}}
\def\ra{{\textnormal{a}}}
\def\rb{{\textnormal{b}}}
\def\rc{{\textnormal{c}}}
\def\rd{{\textnormal{d}}}
\def\re{{\textnormal{e}}}
\def\rf{{\textnormal{f}}}
\def\rg{{\textnormal{g}}}
\def\rh{{\textnormal{h}}}
\def\ri{{\textnormal{i}}}
\def\rj{{\textnormal{j}}}
\def\rk{{\textnormal{k}}}
\def\rl{{\textnormal{l}}}
\def\rn{{\textnormal{n}}}
\def\ro{{\textnormal{o}}}
\def\rp{{\textnormal{p}}}
\def\rq{{\textnormal{q}}}
\def\rr{{\textnormal{r}}}
\def\rs{{\textnormal{s}}}
\def\rt{{\textnormal{t}}}
\def\ru{{\textnormal{u}}}
\def\rv{{\textnormal{v}}}
\def\rw{{\textnormal{w}}}
\def\rx{{\textnormal{x}}}
\def\ry{{\textnormal{y}}}
\def\rz{{\textnormal{z}}}

\def\rvepsilon{{\mathbf{\epsilon}}}
\def\rvtheta{{\mathbf{\theta}}}
\def\rva{{\mathbf{a}}}
\def\rvb{{\mathbf{b}}}
\def\rvc{{\mathbf{c}}}
\def\rvd{{\mathbf{d}}}
\def\rve{{\mathbf{e}}}
\def\rvf{{\mathbf{f}}}
\def\rvg{{\mathbf{g}}}
\def\rvh{{\mathbf{h}}}
\def\rvu{{\mathbf{i}}}
\def\rvj{{\mathbf{j}}}
\def\rvk{{\mathbf{k}}}
\def\rvl{{\mathbf{l}}}
\def\rvm{{\mathbf{m}}}
\def\rvn{{\mathbf{n}}}
\def\rvo{{\mathbf{o}}}
\def\rvp{{\mathbf{p}}}
\def\rvq{{\mathbf{q}}}
\def\rvr{{\mathbf{r}}}
\def\rvs{{\mathbf{s}}}
\def\rvt{{\mathbf{t}}}
\def\rvu{{\mathbf{u}}}
\def\rvv{{\mathbf{v}}}
\def\rvw{{\mathbf{w}}}
\def\rvx{{\mathbf{x}}}
\def\rvy{{\mathbf{y}}}
\def\rvz{{\mathbf{z}}}

\def\erva{{\textnormal{a}}}
\def\ervb{{\textnormal{b}}}
\def\ervc{{\textnormal{c}}}
\def\ervd{{\textnormal{d}}}
\def\erve{{\textnormal{e}}}
\def\ervf{{\textnormal{f}}}
\def\ervg{{\textnormal{g}}}
\def\ervh{{\textnormal{h}}}
\def\ervi{{\textnormal{i}}}
\def\ervj{{\textnormal{j}}}
\def\ervk{{\textnormal{k}}}
\def\ervl{{\textnormal{l}}}
\def\ervm{{\textnormal{m}}}
\def\ervn{{\textnormal{n}}}
\def\ervo{{\textnormal{o}}}
\def\ervp{{\textnormal{p}}}
\def\ervq{{\textnormal{q}}}
\def\ervr{{\textnormal{r}}}
\def\ervs{{\textnormal{s}}}
\def\ervt{{\textnormal{t}}}
\def\ervu{{\textnormal{u}}}
\def\ervv{{\textnormal{v}}}
\def\ervw{{\textnormal{w}}}
\def\ervx{{\textnormal{x}}}
\def\ervy{{\textnormal{y}}}
\def\ervz{{\textnormal{z}}}

\def\rmA{{\mathbf{A}}}
\def\rmB{{\mathbf{B}}}
\def\rmC{{\mathbf{C}}}
\def\rmD{{\mathbf{D}}}
\def\rmE{{\mathbf{E}}}
\def\rmF{{\mathbf{F}}}
\def\rmG{{\mathbf{G}}}
\def\rmH{{\mathbf{H}}}
\def\rmI{{\mathbf{I}}}
\def\rmJ{{\mathbf{J}}}
\def\rmK{{\mathbf{K}}}
\def\rmL{{\mathbf{L}}}
\def\rmM{{\mathbf{M}}}
\def\rmN{{\mathbf{N}}}
\def\rmO{{\mathbf{O}}}
\def\rmP{{\mathbf{P}}}
\def\rmQ{{\mathbf{Q}}}
\def\rmR{{\mathbf{R}}}
\def\rmS{{\mathbf{S}}}
\def\rmT{{\mathbf{T}}}
\def\rmU{{\mathbf{U}}}
\def\rmV{{\mathbf{V}}}
\def\rmW{{\mathbf{W}}}
\def\rmX{{\mathbf{X}}}
\def\rmY{{\mathbf{Y}}}
\def\rmZ{{\mathbf{Z}}}

\def\ermA{{\textnormal{A}}}
\def\ermB{{\textnormal{B}}}
\def\ermC{{\textnormal{C}}}
\def\ermD{{\textnormal{D}}}
\def\ermE{{\textnormal{E}}}
\def\ermF{{\textnormal{F}}}
\def\ermG{{\textnormal{G}}}
\def\ermH{{\textnormal{H}}}
\def\ermI{{\textnormal{I}}}
\def\ermJ{{\textnormal{J}}}
\def\ermK{{\textnormal{K}}}
\def\ermL{{\textnormal{L}}}
\def\ermM{{\textnormal{M}}}
\def\ermN{{\textnormal{N}}}
\def\ermO{{\textnormal{O}}}
\def\ermP{{\textnormal{P}}}
\def\ermQ{{\textnormal{Q}}}
\def\ermR{{\textnormal{R}}}
\def\ermS{{\textnormal{S}}}
\def\ermT{{\textnormal{T}}}
\def\ermU{{\textnormal{U}}}
\def\ermV{{\textnormal{V}}}
\def\ermW{{\textnormal{W}}}
\def\ermX{{\textnormal{X}}}
\def\ermY{{\textnormal{Y}}}
\def\ermZ{{\textnormal{Z}}}

\def\vzero{{\bm{0}}}
\def\vone{{\bm{1}}}
\def\vmu{{\bm{\mu}}}
\def\vtheta{{\bm{\theta}}}
\def\va{{\bm{a}}}
\def\vb{{\bm{b}}}
\def\vc{{\bm{c}}}
\def\vd{{\bm{d}}}
\def\ve{{\bm{e}}}
\def\vf{{\bm{f}}}
\def\vg{{\bm{g}}}
\def\vh{{\bm{h}}}
\def\vi{{\bm{i}}}
\def\vj{{\bm{j}}}
\def\vk{{\bm{k}}}
\def\vl{{\bm{l}}}
\def\vm{{\bm{m}}}
\def\vn{{\bm{n}}}
\def\vo{{\bm{o}}}
\def\vp{{\bm{p}}}
\def\vq{{\bm{q}}}
\def\vr{{\bm{r}}}
\def\vs{{\bm{s}}}
\def\vt{{\bm{t}}}
\def\vu{{\bm{u}}}
\def\vv{{\bm{v}}}
\def\vw{{\bm{w}}}
\def\vx{{\bm{x}}}
\def\vy{{\bm{y}}}
\def\vz{{\bm{z}}}

\def\evalpha{{\alpha}}
\def\evbeta{{\beta}}
\def\evepsilon{{\epsilon}}
\def\evlambda{{\lambda}}
\def\evomega{{\omega}}
\def\evmu{{\mu}}
\def\evpsi{{\psi}}
\def\evsigma{{\sigma}}
\def\evtheta{{\theta}}
\def\eva{{a}}
\def\evb{{b}}
\def\evc{{c}}
\def\evd{{d}}
\def\eve{{e}}
\def\evf{{f}}
\def\evg{{g}}
\def\evh{{h}}
\def\evi{{i}}
\def\evj{{j}}
\def\evk{{k}}
\def\evl{{l}}
\def\evm{{m}}
\def\evn{{n}}
\def\evo{{o}}
\def\evp{{p}}
\def\evq{{q}}
\def\evr{{r}}
\def\evs{{s}}
\def\evt{{t}}
\def\evu{{u}}
\def\evv{{v}}
\def\evw{{w}}
\def\evx{{x}}
\def\evy{{y}}
\def\evz{{z}}

\def\mA{{\bm{A}}}
\def\mB{{\bm{B}}}
\def\mC{{\bm{C}}}
\def\mD{{\bm{D}}}
\def\mE{{\bm{E}}}
\def\mF{{\bm{F}}}
\def\mG{{\bm{G}}}
\def\mH{{\bm{H}}}
\def\mI{{\bm{I}}}
\def\mJ{{\bm{J}}}
\def\mK{{\bm{K}}}
\def\mL{{\bm{L}}}
\def\mM{{\bm{M}}}
\def\mN{{\bm{N}}}
\def\mO{{\bm{O}}}
\def\mP{{\bm{P}}}
\def\mQ{{\bm{Q}}}
\def\mR{{\bm{R}}}
\def\mS{{\bm{S}}}
\def\mT{{\bm{T}}}
\def\mU{{\bm{U}}}
\def\mV{{\bm{V}}}
\def\mW{{\bm{W}}}
\def\mX{{\bm{X}}}
\def\mY{{\bm{Y}}}
\def\mZ{{\bm{Z}}}
\def\mBeta{{\bm{\beta}}}
\def\mPhi{{\bm{\Phi}}}
\def\mLambda{{\bm{\Lambda}}}
\def\mSigma{{\bm{\Sigma}}}

\DeclareMathAlphabet{\mathsfit}{\encodingdefault}{\sfdefault}{m}{sl}
\SetMathAlphabet{\mathsfit}{bold}{\encodingdefault}{\sfdefault}{bx}{n}
\newcommand{\tens}[1]{\bm{\mathsfit{#1}}}
\def\tA{{\tens{A}}}
\def\tB{{\tens{B}}}
\def\tC{{\tens{C}}}
\def\tD{{\tens{D}}}
\def\tE{{\tens{E}}}
\def\tF{{\tens{F}}}
\def\tG{{\tens{G}}}
\def\tH{{\tens{H}}}
\def\tI{{\tens{I}}}
\def\tJ{{\tens{J}}}
\def\tK{{\tens{K}}}
\def\tL{{\tens{L}}}
\def\tM{{\tens{M}}}
\def\tN{{\tens{N}}}
\def\tO{{\tens{O}}}
\def\tP{{\tens{P}}}
\def\tQ{{\tens{Q}}}
\def\tR{{\tens{R}}}
\def\tS{{\tens{S}}}
\def\tT{{\tens{T}}}
\def\tU{{\tens{U}}}
\def\tV{{\tens{V}}}
\def\tW{{\tens{W}}}
\def\tX{{\tens{X}}}
\def\tY{{\tens{Y}}}
\def\tZ{{\tens{Z}}}


\def\gA{{\mathcal{A}}}
\def\gB{{\mathcal{B}}}
\def\gC{{\mathcal{C}}}
\def\gD{{\mathcal{D}}}
\def\gE{{\mathcal{E}}}
\def\gF{{\mathcal{F}}}
\def\gG{{\mathcal{G}}}
\def\gH{{\mathcal{H}}}
\def\gI{{\mathcal{I}}}
\def\gJ{{\mathcal{J}}}
\def\gK{{\mathcal{K}}}
\def\gL{{\mathcal{L}}}
\def\gM{{\mathcal{M}}}
\def\gN{{\mathcal{N}}}
\def\gO{{\mathcal{O}}}
\def\gP{{\mathcal{P}}}
\def\gQ{{\mathcal{Q}}}
\def\gR{{\mathcal{R}}}
\def\gS{{\mathcal{S}}}
\def\gT{{\mathcal{T}}}
\def\gU{{\mathcal{U}}}
\def\gV{{\mathcal{V}}}
\def\gW{{\mathcal{W}}}
\def\gX{{\mathcal{X}}}
\def\gY{{\mathcal{Y}}}
\def\gZ{{\mathcal{Z}}}

\def\sA{{\mathbb{A}}}
\def\sB{{\mathbb{B}}}
\def\sC{{\mathbb{C}}}
\def\sD{{\mathbb{D}}}
\def\sF{{\mathbb{F}}}
\def\sG{{\mathbb{G}}}
\def\sH{{\mathbb{H}}}
\def\sI{{\mathbb{I}}}
\def\sJ{{\mathbb{J}}}
\def\sK{{\mathbb{K}}}
\def\sL{{\mathbb{L}}}
\def\sM{{\mathbb{M}}}
\def\sN{{\mathbb{N}}}
\def\sO{{\mathbb{O}}}
\def\sP{{\mathbb{P}}}
\def\sQ{{\mathbb{Q}}}
\def\sR{{\mathbb{R}}}
\def\sS{{\mathbb{S}}}
\def\sT{{\mathbb{T}}}
\def\sU{{\mathbb{U}}}
\def\sV{{\mathbb{V}}}
\def\sW{{\mathbb{W}}}
\def\sX{{\mathbb{X}}}
\def\sY{{\mathbb{Y}}}
\def\sZ{{\mathbb{Z}}}

\def\emLambda{{\Lambda}}
\def\emA{{A}}
\def\emB{{B}}
\def\emC{{C}}
\def\emD{{D}}
\def\emE{{E}}
\def\emF{{F}}
\def\emG{{G}}
\def\emH{{H}}
\def\emI{{I}}
\def\emJ{{J}}
\def\emK{{K}}
\def\emL{{L}}
\def\emM{{M}}
\def\emN{{N}}
\def\emO{{O}}
\def\emP{{P}}
\def\emQ{{Q}}
\def\emR{{R}}
\def\emS{{S}}
\def\emT{{T}}
\def\emU{{U}}
\def\emV{{V}}
\def\emW{{W}}
\def\emX{{X}}
\def\emY{{Y}}
\def\emZ{{Z}}
\def\emSigma{{\Sigma}}

\newcommand{\etens}[1]{\mathsfit{#1}}
\def\etLambda{{\etens{\Lambda}}}
\def\etA{{\etens{A}}}
\def\etB{{\etens{B}}}
\def\etC{{\etens{C}}}
\def\etD{{\etens{D}}}
\def\etE{{\etens{E}}}
\def\etF{{\etens{F}}}
\def\etG{{\etens{G}}}
\def\etH{{\etens{H}}}
\def\etI{{\etens{I}}}
\def\etJ{{\etens{J}}}
\def\etK{{\etens{K}}}
\def\etL{{\etens{L}}}
\def\etM{{\etens{M}}}
\def\etN{{\etens{N}}}
\def\etO{{\etens{O}}}
\def\etP{{\etens{P}}}
\def\etQ{{\etens{Q}}}
\def\etR{{\etens{R}}}
\def\etS{{\etens{S}}}
\def\etT{{\etens{T}}}
\def\etU{{\etens{U}}}
\def\etV{{\etens{V}}}
\def\etW{{\etens{W}}}
\def\etX{{\etens{X}}}
\def\etY{{\etens{Y}}}
\def\etZ{{\etens{Z}}}

\newcommand{\pdata}{p_{\rm{data}}}
\newcommand{\ptrain}{\hat{p}_{\rm{data}}}
\newcommand{\Ptrain}{\hat{P}_{\rm{data}}}
\newcommand{\pmodel}{p_{\rm{model}}}
\newcommand{\Pmodel}{P_{\rm{model}}}
\newcommand{\ptildemodel}{\tilde{p}_{\rm{model}}}
\newcommand{\pencode}{p_{\rm{encoder}}}
\newcommand{\pdecode}{p_{\rm{decoder}}}
\newcommand{\precons}{p_{\rm{reconstruct}}}

\newcommand{\laplace}{\mathrm{Laplace}} %

\newcommand{\E}{\mathbb{E}}
\newcommand{\Ls}{\mathcal{L}}
\newcommand{\R}{\mathbb{R}}
\newcommand{\emp}{\tilde{p}}
\newcommand{\lr}{\alpha}
\newcommand{\reg}{\lambda}
\newcommand{\rect}{\mathrm{rectifier}}
\newcommand{\softmax}{\mathrm{softmax}}
\newcommand{\sigmoid}{\sigma}
\newcommand{\softplus}{\zeta}
\newcommand{\KL}{D_{\mathrm{KL}}}
\newcommand{\Var}{\mathrm{Var}}
\newcommand{\standarderror}{\mathrm{SE}}
\newcommand{\Cov}{\mathrm{Cov}}
\newcommand{\normlzero}{L^0}
\newcommand{\normlone}{L^1}
\newcommand{\normltwo}{L^2}
\newcommand{\normlp}{L^p}
\newcommand{\normmax}{L^\infty}

\newcommand{\parents}{Pa} %

\DeclareMathOperator*{\argmax}{arg\,max}
\DeclareMathOperator*{\argmin}{arg\,min}

\DeclareMathOperator{\sign}{sign}
\DeclareMathOperator{\Tr}{Tr}
\let\ab\allowbreak

\newcommand{\fk}[1]{ {\textcolor{OliveGreen}{\textbf{FK:}~#1}} }



\newcounter{setting}[section]
\renewcommand{\thesetting}{\thesection.\arabic{setting}}

\newcounter{model}[section]
\renewcommand{\themodel}{\thesection.\arabic{model}}

\newcounter{initialization}[section]
\renewcommand{\theinitialization}{\thesection.\arabic{initialization}}

\newenvironment{setting}
{%
    \refstepcounter{setting}%
    \par\medskip
    \noindent\textbf{Setting~\thesetting.}\enspace
}
{%
    \par\medskip
}

\newenvironment{model}
{%
    \refstepcounter{model}%
    \par\medskip
    \noindent\textbf{Model~\themodel.}\enspace
}
{%
    \par\medskip
}

\newenvironment{initialization}
{%
    \refstepcounter{initialization}%
    \par\medskip
    \noindent\textbf{Initialization~\theinitialization.}\enspace
}
{%
    \par\medskip
}


\usepackage[
    colorlinks=true,
    citecolor=MidnightBlue,
    linkcolor=MidnightBlue,
    urlcolor=MidnightBlue
]{hyperref}

\usepackage{cleveref}

\setlength{\parindent}{0pt}
\setlength{\parskip}{0.6\baselineskip}

\newcommand{\experimentfigure}[2][]{%
  \IfFileExists{#2}{%
    \includegraphics[#1]{#2}%
  }{%
    \fbox{%
      \parbox[c][0.22\textheight][c]{0.88\linewidth}{%
        \centering
        Figure file not found:\\[0.5em]
        \texttt{\detokenize{#2}}%
      }%
    }%
  }%
}


\newcounter{datasetting}
\renewcommand{\thedatasetting}{\arabic{datasetting}}
\crefname{datasetting}{Data Setting}{Data Settings}
\Crefname{datasetting}{Data Setting}{Data Settings}

\newcounter{trainingsetup}
\renewcommand{\thetrainingsetup}{\arabic{trainingsetup}}
\crefname{trainingsetup}{Training Setup}{Training Setups}
\Crefname{trainingsetup}{Training Setup}{Training Setups}

\newcounter{initializationsetting}
\renewcommand{\theinitializationsetting}{\arabic{initializationsetting}}
\crefname{initializationsetting}{Initialization}{Initializations}
\Crefname{initializationsetting}{Initialization}{Initializations}



\tcbset{
  theorysetting/.style={
    enhanced,
    breakable,
    boxrule=0.8pt,
    arc=2pt,
    left=7pt,
    right=7pt,
    top=6pt,
    bottom=6pt,
    before skip=8pt,
    after skip=8pt,
    fonttitle=\bfseries,
  }
}

\NewDocumentEnvironment{datasetting}{O{}}
{
  \refstepcounter{datasetting}
  \begin{tcolorbox}[
    theorysetting,
    colback=blue!3,
    colframe=blue!45!black,
    title={
      Data Setting~\thedatasetting
      \if\relax\detokenize{#1}\relax\else\ --- #1\fi
    }
  ]
}
{
  \end{tcolorbox}
}

\NewDocumentEnvironment{trainingsetup}{O{}}
{
  \refstepcounter{trainingsetup}
  \begin{tcolorbox}[
    theorysetting,
    colback=green!3,
    colframe=green!40!black,
    title={
      Training Setup~\thetrainingsetup
      \if\relax\detokenize{#1}\relax\else\ --- #1\fi
    }
  ]
}
{
  \end{tcolorbox}
}

\NewDocumentEnvironment{initializationsetting}{O{}}
{
  \refstepcounter{initializationsetting}
  \begin{tcolorbox}[
    theorysetting,
    colback=orange!4,
    colframe=orange!55!black,
    title={
      Initialization~\theinitializationsetting
      \if\relax\detokenize{#1}\relax\else\ --- #1\fi
    }
  ]
}
{
  \end{tcolorbox}
}


\newtheorem{theorem}{Theorem}

\newaliascnt{lemma}{theorem}
\newtheorem{lemma}[lemma]{Lemma}
\aliascntresetthe{lemma}

\newaliascnt{corollary}{theorem}
\newtheorem{corollary}[corollary]{Corollary}
\aliascntresetthe{corollary}

\crefname{theorem}{Theorem}{Theorems}
\Crefname{theorem}{Theorem}{Theorems}

\crefname{lemma}{Lemma}{Lemmas}
\Crefname{lemma}{Lemma}{Lemmas}

\crefname{corollary}{Corollary}{Corollaries}
\Crefname{corollary}{Corollary}{Corollaries}

\newcommand{\objective}[1]{%
\begin{tcolorbox}[
    colback=blue!4,
    colframe=blue!45!black,
    boxrule=0.6pt,
    arc=2mm,
    left=1.2mm,
    right=1.2mm,
    top=1mm,
    bottom=1mm,
    before skip=0.8em,
    after skip=1em,
    title={Objective for Section~\thesubsection},
    fonttitle=\bfseries
]
#1
\end{tcolorbox}
}

\crefname{lemma}{Lemma}{Lemmas}
\newcommand\misha[1]{\textcolor{blue}{(Misha: #1)}}

\begin{document}

\begin{center}

\makebox[\textwidth][c]{%
  \parbox[t]{0.84\paperwidth}{%
    \centering
    \Large\bfseries
    Revenge of Monosemanticity: Neuron Specialization\\[-0.05em]
    as a New Form of Feature Learning in MLPs
    \par
  }%
}

\vspace{1.1em}

{\large
\begin{tabular}{cc}
Amirhesam Abedsoltan$^{1}$
&
Enric Boix-Adsera$^{2}$ \\[0.35em]
Fivos Kalogiannis$^{1}$
&
Mikhail Belkin$^{3,1}$
\end{tabular}
}

\vspace{0.9em}

{\small
$^{1}$Department of Computer Science and Engineering, UC San Diego
\par
\makebox[\linewidth][c]{%
$^{2}$Department of Statistics and Data Science, The Wharton School, University of Pennsylvania}
\par
$^{3}$Halıcıoğlu Data Science Institute, UC San Diego
\par}

\end{center}


\vspace{1.0em}

\begin{abstract}
Understanding how neural networks learn and organize features is central to understanding their behavior. Much existing theory of feature learning has focused on the emergence of a global low-dimensional representation. We show that this picture is incomplete. In regression problems with clustered data, we demonstrate that multilayer perceptrons (MLPs) naturally develop monosemantic specialized neurons: individual neurons become strongly aligned with a specific predictive feature relevant to a particular region of the input space. Rather than learning a single global low-dimensional representation, MLPs learn a collection of local low-dimensional representations. We show that this ability to specialize gives MLPs a provable data-efficiency advantage over feature-learning methods based on a global low-dimensional representation.
\end{abstract}

\vspace{1.0em}




\section{Introduction}


A major advance in understanding neural networks was the recognition that neural networks can perform feature learning by exploiting low-dimensional structure in the target, where prediction depends on only a small number of directions in the input space. In such settings, neural networks can adapt their representations to identify these predictive directions~\citep{bach2017breaking,soltanolkotabi2017learning,yehudai2019power,wei2019regularization,ghorbani2019limitations,ghorbani2020when,malach2021quantifying,abbe2022merged,abbe2023sgd,ba2022highdimensional,damian2022neural,dandi2024giant,bruna2025survey}.

In this setting, the response takes the form
\[
f(x)
=
g\!\left(U^{\top}x\right),
\qquad
U\in\mathbb{R}^{d\times r},
\qquad
r\ll d,
\]
where \(g:\mathbb{R}^r\to\mathbb{R}\) is the \emph{link function}. Although \(x\in\mathbb{R}^d\), the response
depends only on the \(r\)-dimensional predictive subspace spanned by the
columns of $U$.
During training, a neural network can learn features aligned with that predictive subspace, effectively reducing the dimensionality of the problem from the ambient dimension $d$ to the intrinsic dimension $r$.


This perspective views feature learning primarily as a mechanism for discovering low-dimensional predictive structure: supervision reshapes the representation so that learning is governed by the intrinsic dimension of the task rather than the ambient dimension of the input.



In our work, we identify a qualitatively different mechanism of feature learning
in multilayer perceptrons (MLPs) that is not tied to recovering a
single global low-dimensional predictive subspace.

\ifdefined\iclrformat
\begin{wrapfigure}{r}{0.37\textwidth}
\else
\begin{wrapfigure}{r}{0.45\textwidth}
\fi
    \centering
    \vspace{-15pt}
    \includegraphics[width=\linewidth]
    {figs/clusters_equation_box_refined_cropped.png}
\end{wrapfigure}

In this setting, the data are drawn from multiple clusters, and each cluster can
have its own predictive directions and link function. Thus, prediction may be
low-dimensional within each cluster even when there is no low-dimensional
global structure. The diagram on the right illustrates the \emph{single-index}
case, where each cluster has one predictive direction.

As we will see, MLPs can simultaneously discover the cluster structure and learn the predictive functions relevant to each cluster. This uncovers a form of feature learning that goes beyond recovering a single global low-dimensional predictive subspace.





\noindent\textbf{Main contributions.}



We show that, in trained MLPs, a substantial fraction of individual neurons become \emph{monosemantic}, 
specializing by aligning predominantly with a single cluster-specific predictive direction. 
This specialization allows the MLP to learn both the relevant local low-dimensional features and an implicit 
clustering that determines where each feature is useful. The resulting behavior resembles a mixture of experts,
 but emerges within a standard MLP without an explicit routing module or expert decomposition.


We show empirically that this specialization leads to increasingly favorable sample-complexity 
scaling as the number of clusters grows, allowing MLPs to maintain strong performance even when 
no single global low-dimensional structure exists. 
While specialization occurs for standard ReLU and GELU~\citep{hendrycks2016gaussian}, 
modern architectures with multiplicative gating, such as ReGLU and SwiGLU~\citep{shazeer2020glu}, 
can substantially improve sample complexity
in this setting.


We also establish a provable advantage of MLPs over methods based on a global low-dimensional representation. In particular, in our theoretical setting, as the number of clusters $K$ increases, the MLP has the ability to achieve optimality (consistency) with only polynomially many training samples, $n=\mathrm{poly}(K)$. In contrast, methods based on a single global low-dimensional representation have a constant error when the number of samples is any polynomial in $K$.










\subsection{Discussion}
The form of feature learning identified in this work differs from classical feature learning, where the main goal
is to recover a single global low-dimensional predictive subspace. In such
settings, methods designed for learning low-dimensional features can perform as well as or
better than MLPs. A prominent example is the Recursive Feature Machine (RFM), a supervised kernel-based method that uses input-response pairs to learn a global low-dimensional predictive subspace~\citep{radhakrishnan2024mechanism,radhakrishnan2025linear}.
RFM is therefore a particularly informative comparison: like an MLP, it learns
features from the data rather than operating with a fixed representation.

\begin{figure}[h]
    \centering
    \includegraphics[width=\linewidth]{figs/main_idea_final.png}
\end{figure}


To study this phenomenon systematically, we consider a mixture of single-index or
multi-index models. Suppose the data are drawn from \(K\) clusters. For an input
\(x\in\mathbb{R}^d\) belonging to cluster \(c\in[K]\), the target is
\[
f(x)=g_c\!\left(U_c^\top x\right),
\qquad
U_c\in\mathbb{R}^{d\times r}.
\]
Thus, prediction is $r$-dimensional within each cluster, but both the
predictive subspace spanned by the columns of \(U_c\) and the link function
\(g_c\) may vary across clusters. The case \(r=1\) is a single-index model,
whereas \(r>1\) gives a multi-index model.

Note that the standard single- or multi-index model corresponds to \(K=1\). When
\(K>1\), however, the column spaces of \(U_1,\ldots,U_K\) may not generally admit a common
low-dimensional subspace. Hence, prediction can be low-dimensional locally within
each cluster even though there is no useful global low-dimensional predictive
subspace to recover.

\subsection{Prior Work}
\label{sec:prior_work}




\paragraph{Feature learning of global low-dimensional structure.}
A large body of work studies feature learning for multi-index
targets. These works have shown that neural-network models can adapt to an unknown
low-dimensional subspace of the inputs, succeeding where fixed-kernel methods are sample-inefficient
\citep{bach2017breaking,soltanolkotabi2017learning,
yehudai2019power,ghorbani2020when}. Recently, the literature has characterized how gradient-based training succeeds in recovering this subspace for isotropic data
\citep{abbe2022merged,abbe2023sgd,damian2022neural,ba2022highdimensional,dandi2024giant,damian2025generative,bruna2025survey}, and in the related setting of classifying Gaussian-mixture data with a constant number of clusters
\citep{pmlr-v139-refinetti21b,arous2024highdimensional}.


However, in our work, we prove that neuron specialization is a different form of feature learning, allowing MLPs to efficiently tackle settings where a single global low-dimensional representation is insufficient.

\paragraph{Neuron specialization.} The emergence of
specialized neurons has been observed in several settings. For instance, training neural networks to do modular arithmetic leads to specialized neurons, each representing different Fourier components
\citep{nanda2023progress,gromov2023grokking,morwani2024feature,
he2026mechanism,he2026neural}. When learning XOR-type targets, it has also been shown that neurons specialize to four clusters rather than being distributed evenly across a predictive low-dimensional subspace \citep{frei2023random,glasgow2024sgd}. In teacher--student settings, student neurons have also been shown to align frequently with teacher neurons \citep{tian2020student,oostwal2021hidden,zhu2025gradient}, although this depends on initialization \citep{jarvis2025theory}.

These works mainly study specialization to globally relevant features or to cluster directions that themselves determine the target. In our work, neurons specialize to locally predictive features. Thus, an MLP must not only learn the predictive directions, but also preserve their association with the clusters in which they are relevant. Moreover, these works do not isolate a sample-complexity advantage of specialization over an adaptive feature-learning method based on a single global representation, such as the Recursive Feature Machine \citep{radhakrishnan2024mechanism}. We establish such an advantage: when cluster-specific predictive directions collectively span a high-dimensional space, neuron specialization allows an MLP to preserve these local feature-cluster associations and achieve better sample complexity.

\paragraph{Monosemanticity.}
Interpretability works on language models have found that some individual neurons are monosemantic, responding primarily to a single interpretable concept or pattern, while others are \emph{polysemantic}, responding to multiple distinct concepts or patterns \citep{bills2023language,elhage2022solu}. More recent work
has argued that monosemantic features need not align with individual neurons,
and has used sparse dictionary learning to recover more monosemantic feature
directions from neural representations
\citep{bricken2023monosemanticity,templeton2024scaling}. Our work shows that,
at least in simple MLPs, monosemantic features can emerge directly at the level
of individual neurons and that this neuron-level specialization can improve
sample efficiency compared with methods based on a global low-dimensional representation.
\paragraph{Mixture-of-experts models.} A complementary line of work studies specialization and latent-cluster recovery in mixture-of-experts (MoE) architectures. An MoE router can learn cluster-center features that partition a clustered classification problem into simpler subproblems handled by different experts \citep{chen2022towards}. Related theoretical and empirical results show that a learned router can route inputs according to latent clusters \citep{dikkala-etal-2023-benefits}. Particularly close to our statistical setting, recent work studies nonlinear regression with an underlying cluster structure of single-index models and shows that an MoE trained by SGD can detect the latent organization and divide the task into cluster-specific subproblems; under its assumptions, a vanilla neural network does not detect this organization within the same polynomial complexity regime \citep{pmlr-v267-kawata25a}.

These works build specialization into the architecture through an explicit router and separate experts. In contrast, we study a standard MLP with no explicit routing or expert decomposition, and show that both the implicit clustering and the corresponding local predictive-feature specialization can emerge through the specialization of individual neurons.


\subsection{Paper Structure}
The remainder of the paper studies neuron specialization in MLPs empirically and theoretically. Section~\ref{sec:prelim} introduces the setting. Section~\ref{sec:Experiments} shows empirically that MLPs jointly learn cluster structure and cluster-specific predictive functions, achieving increasingly favorable sample complexity over global low-rank methods as the number of clusters grows. We also show that gated activations such as ReGLU and SwiGLU improve sample complexity over ReLU and GELU. Finally, Section~\ref{sec:specialization-theory} theoretically analyzes the emergence of neuron specialization in simplified settings.




\section{Preliminaries}\label{sec:prelim}

Throughout our synthetic experiments and theoretical analysis, we consider
data drawn from a Gaussian mixture model with \(K\) clusters. The cluster
index \(c\) is drawn uniformly from \(\{1,\ldots,K\}\). For each cluster \(c\),
let \(\mu_c\in\mathbb{R}^d\) denote its mean and
\(\Sigma_c\in\mathbb{R}^{d\times d}\) its covariance matrix. Conditioned on
cluster \(c\), the covariates \(x\in\mathbb{R}^d\) are distributed as
\[
x \mid c \sim \mathcal{N}(\mu_c, \Sigma_c).
\]
In all settings considered in this paper, we choose the cluster means so that
distinct clusters are well separated. The majority of our experiments, as well as the
theoretical analysis, use a cluster-specific single-index model with additive
Gaussian noise:
\[
y = g_c\!\left(\langle x, u_c\rangle\right) + \varepsilon,
\qquad
\varepsilon \sim \mathcal{N}(0,\sigma^2),
\]
where \(\varepsilon\) is independent of \(x\) and \(c\),
\(u_c\in\mathbb{R}^d\), with \(\|u_c\|_2=1\), is the cluster-specific
predictive direction, and \(g_c:\mathbb{R}\to\mathbb{R}\) is a
nonlinear link function that may vary across clusters.

In the last experiment, we also consider a multi-index
setting in which each cluster \(c\) has three orthonormal predictive
directions \(u_{c,1},u_{c,2},u_{c,3}\). The same cluster-specific link
function is applied along each direction and the resulting components are
aggregated:
\[
y
=
\frac{1}{\sqrt 3}
\sum_{r=1}^{3}
g_c\!\left(\langle x,u_{c,r}\rangle\right)
+\varepsilon.
\]

Throughout, we use mean squared error (MSE) as the regression loss,
\[
\mathcal{L}(f)
=
\mathbb{E}\!\left[(f(x)-y)^2\right],
\]
where \(f:\mathbb{R}^d\to\mathbb{R}\) denotes the learned predictor.


\section{Experiments}
\label{sec:Experiments}
In this section, we empirically show that MLPs can learn both the cluster structure and the predictive function within each cluster through monosemantic specialized neurons. This gives rise to a form of feature learning that goes beyond recovering a single low-dimensional predictive subspace. We demonstrate these phenomena through four experiments, each presented in a separate subsection.




\subsection{MLPs Develop Specialized First-Layer Neurons}
\label{sec:coordinate-specialization}

We first show that individual first-layer neurons specialize to the
cluster-specific predictive directions. We use the clustered single-index
model introduced in Section~\ref{sec:prelim}, and for this experiment choose
the predictive directions to coincide with distinct coordinate axes. This
choice makes specialization directly visible in the learned first-layer
weights.

\paragraph{Main findings.}

Table~\ref{tab:coordinate-specialization} shows that a significant number of
first-layer neurons in trained MLPs become \emph{monosemantic}, specializing to
individual cluster-specific predictive directions. We illustrate this specialization in
Fig.~\ref{fig:coordinate-specialization-hermite3-neurons} by showing
the first-layer weights of selected neurons with high cosine similarity to their corresponding predictive directions.



\begin{figure}[H]
    \centering
    \includegraphics[width=1.0\textwidth]
    {figs/relu_high_cosine_neurons.pdf}
    \caption{\textbf{Representative specialized neurons.}
     These specialized neurons concentrate
    nearly all of their weights on one cluster-specific predictive coordinate,
    with little weight on the cluster-position and nuisance coordinates.
    Because every predictive direction \(v_c\) is a standard basis vector, concentration on a
    predictive coordinate is equivalent to alignment with the corresponding
    predictive direction.}
    \label{fig:coordinate-specialization-hermite3-neurons}
\end{figure}







\paragraph{Data setting.}
Following the notation of Section~\ref{sec:prelim}, we use \(K=10\) clusters
in \(d=40\). We choose the cluster centers and predictive directions to lie
along separate coordinate axes:
\[
[\mu_c]_j =
\begin{cases}
20, & j=c,\\
0, & \text{otherwise},
\end{cases}
\qquad
[v_c]_j =
\begin{cases}
1, & j=K+c,\\
0, & \text{otherwise},
\end{cases}
\qquad
j\in\{1,\ldots,d\}.
\]
We further set
\[
\Sigma_c = I_d/d,
\qquad
\varepsilon \sim \mathcal{N}(0,0.005^2).
\]
Every cluster uses the normalized third-Hermite polynomial as its link
function. We make this choice to match the theoretical setting analyzed in
Section~\ref{sec:early-specialization}.

Thus, the first \(K\) coordinates identify the clusters, the next \(K\)
coordinates contain the predictive directions, and the remaining coordinates
are nuisance dimensions. Because each \(v_c\) is a standard basis vector,
concentration on a predictive coordinate is equivalent to alignment with the
corresponding predictive direction.

\paragraph{Specialization measures.}
For a neuron's first-layer weight \(w_j\), we measure its maximum
predictive-direction cosine by
\[
    \max_{1\leq c\leq K}
    \frac{|\langle w_j,v_c\rangle|}
    {\|w_j\|_2\,\|v_c\|_2}.
\]
Table~\ref{tab:coordinate-specialization} reports the mean and standard
deviation of the percentage of active neurons for which this cosine is at
least \(0.71\) or \(0.90\). The first threshold means that at
least approximately half of the squared weight norm lies along a single predictive
direction, while \(0.90\) is a stricter measure of alignment. We also report
whether the largest predictive coordinate exceeds every other coordinate and
whether it is more than twice as large as every other coordinate. These
coordinate-dominance measures capture specialization when many small
coefficients reduce cosine similarity; see
Appendix~\ref{app:coordinate-specialization-details} for more detail.


\begin{table}[H]
    \centering
    \small
    \renewcommand{\arraystretch}{1.15}
    \setlength{\tabcolsep}{5pt}

    \resizebox{\textwidth}{!}{%
    \begin{tabular}{@{}lccccc@{}}
        \toprule
        & & \multicolumn{4}{c}{\textbf{Specialization measures
        (\% of active neurons)}} \\
        \cmidrule(lr){3-6}
        \textbf{Model/weight}
        & \shortstack{\textbf{Active}\\\textbf{neurons}\\
          \textbf{(out of 2048)}}
        & \shortstack{\textbf{Max. predictive}\\
          \textbf{cosine} \(\mathbf{\geq .71}\,\uparrow\)}
        & \shortstack{\textbf{Max. predictive}\\
          \textbf{cosine} \(\mathbf{\geq .90}\,\uparrow\)}
        & \shortstack{\textbf{Predictive coord.}\\
          \textbf{is largest} \(\mathbf{\uparrow}\)}
        & \shortstack{\textbf{Predictive coord.}\\
          \(\mathbf{>2\times}\) \textbf{every other} \(\mathbf{\uparrow}\)} \\
        \midrule
        ReLU & \(469\pm57\) & \(56.75\pm3.91\) & \(33.22\pm1.92\) & \(92.6\pm0.9\) & \(58.8\pm3.1\) \\
        GELU & \(1{,}747\pm83\) & \(30.15\pm2.48\) & \(18.40\pm1.20\) & \(71.2\pm0.7\) & \(34.6\pm2.2\) \\
        ReGLU gate & \(200\pm56\) & \(51.89\pm1.22\) & \(33.82\pm1.27\) & \(92.7\pm0.6\) & \(54.2\pm1.3\) \\
        ReGLU value & \(200\pm56\) & \(24.98\pm0.52\) & \(3.94\pm2.24\) & \(69\pm1\) & \(22.0\pm1.4\) \\
        SwiGLU gate & \(432\pm5\) & \(25.06\pm2.29\) & \(12.26\pm1.27\) & \(72.3\pm0.9\) & \(23.3\pm2.0\) \\
        SwiGLU value & \(432\pm5\) & \(17.88\pm2.88\) & \(4.62\pm0.96\) & \(86.5\pm2.6\) & \(17.9\pm2.1\) \\
        \bottomrule
    \end{tabular}%
    }

    \caption{
    \textbf{First-layer neurons are monosemantic.}
    Across activations, many neurons align with a single cluster-specific predictive direction. The cosine similarity columns report the percentage of active neurons exceeding the indicated alignment threshold. Appendix~\ref{app:coordinate-specialization-details} describes the activeness criterion, predictive-coordinate dominance, and additional experiments. Entries show mean $\pm$ standard deviation over three runs, resampling the training, validation, and test data, label noise, network initialization, and optimization minibatches.
}
    \label{tab:coordinate-specialization}
\end{table}





\paragraph{Details and additional experiments.}
Appendix~\ref{app:specialization-measures-examples} motivates
predictive-coordinate dominance as a complementary specialization measure
and illustrates it with ReLU and GELU neurons.
Appendix~\ref{app:shared-specialization-analysis} defines the active-neuron
criterion, while Appendix~\ref{app:specialization-null} compares alignment
with the true predictive directions against a random-direction baseline.
Training details are provided in
Appendix~\ref{app:hermite3-specialization-experiment}, and
Appendix~\ref{app:hermite2-specialization-experiment} presents an additional
link-function experiment showing specialization beyond the setting of our
theoretical analysis.




\subsection{MLPs Jointly Learn Cluster Structure and Cluster-Specific Predictive Functions}
\label{sec:local-routing}

Following the clustered single-index model introduced in
Section~\ref{sec:prelim}, we consider regression problems in
\(d=20\) in which both the predictive direction \(v_c\) and link function
\(g_c\) may differ across clusters. We vary the number of clusters over
\(K\in\{1,2,10,50\}\).

\paragraph{Main findings.}
Figure~\ref{fig:local-sample-complexity} and
Table~\ref{tab:singleindex-mixed-k10-specialization-null} highlight three main
conclusions:
\begin{enumerate}
    \item \textbf{Sample complexity as the number of clusters grows.}
    MLPs outperform other methods as the number of clusters increases.
    Remarkably, MLPs achieve performance close to that of the local predictors,
    which are given the true cluster identities and fit a separate RFM or
    Laplace predictor within each cluster. Thus, MLPs recover much of the
    benefit of knowing the cluster structure without ever observing the
    cluster identities.

    \item \textbf{Gated activations enhance cluster-dependent feature learning.}
    ReLU already learns useful cluster-dependent structure and substantially
    outperforms global RFM when many clusters are present. ReGLU improves
    further in this regime. Its separate gate and value branches allow the
    network to select cluster-specific predictive features more directly than
    an ordinary ReLU hidden layer.

    \item \textbf{Specialization beyond axis-aligned orthogonal
    directions.}
     Table~\ref{tab:singleindex-mixed-k10-specialization-null} shows that a substantial fraction of active first-layer neurons specialize to the cluster-specific predictive directions. In this experiment, the predictive
    directions are neither axis-aligned nor constrained to be mutually orthogonal. As a baseline, we show that the alignment of the
    same active neurons with random directions in the same predictive
    subspace is negligible compared with their alignment with the true
    predictive directions.
    Appendix~\ref{app:singleindex-specialization-null} provides the full
    construction and analysis.
\end{enumerate}

\begin{figure}[H]
    \centering

    \includegraphics[width=\textwidth]
    {figs/singleindex_final_n500k_3seed_original_relu_reglu_plot_only.pdf}

    \par\vspace{-0.2em}

    \includegraphics[width=0.84\textwidth]
    {figs/singleindex_final_n500k_3seed_original_relu_reglu_compact_legend_only.pdf}

    \vspace{-0.3em}

    \caption{
    \textbf{MLPs jointly learn clustering and local predictive
    functions.}
    As the number of clusters grows, MLPs, especially ReGLU, maintain strong performance close to local Laplace and RFM, while global RFM becomes increasingly less sample efficient. The local methods are given the true cluster identities and fit a separate predictor within each cluster. Curves show loss averaged over three runs with different link-function assignments, data, label noise, model initialization, and optimization minibatches.
}
    \label{fig:local-sample-complexity}
\end{figure}

\begin{table}[H]
    \centering
    \small
    \renewcommand{\arraystretch}{1.15}
    \setlength{\tabcolsep}{4pt}
    \resizebox{\textwidth}{!}{%
    \begin{tabular}{@{}lccccc@{}}
        \toprule
        & & \multicolumn{4}{c}{\textbf{Specialization measures
        (\% of active neurons)}} \\
        \cmidrule(lr){3-6}
        \textbf{Model/weight}
        & \shortstack{\textbf{Active}\\\textbf{neurons}\\\textbf{(out of 4096)}}
        & \shortstack{\textbf{True}\\\textbf{directions}\\
          \textbf{cosine} \(\mathbf{\geq .71}\,\uparrow\)}
        & \shortstack{\textbf{Random}\\\textbf{directions}\\
          \textbf{cosine} \(\mathbf{\geq .71}\)}
        & \shortstack{\textbf{True}\\\textbf{directions}\\
          \textbf{cosine} \(\mathbf{\geq .90}\,\uparrow\)}
        & \shortstack{\textbf{Random}\\\textbf{directions}\\
          \textbf{cosine} \(\mathbf{\geq .90}\)} \\
        \midrule
        ReLU & \(558\pm36\) & \(12.72\pm3.54\) & \(0.29\pm0.49\) & \(3.02\pm0.79\) & \(0.00\pm0.02\) \\
        GELU & \(287\pm15\) & \(16.87\pm0.79\) & \(0.83\pm0.85\) & \(8.01\pm0.66\) & \(0.00\pm0.04\) \\
        ReGLU gate & \(133\pm44\) & \(16.87\pm4.31\) & \(0.56\pm1.08\) & \(5.02\pm1.31\) & \(0.00\pm0.04\) \\
        ReGLU value & \(133\pm44\) & \(79.65\pm4.15\) & \(6.76\pm7.44\) & \(60.12\pm6.44\) & \(0.05\pm0.46\) \\
        SwiGLU gate & \(248\pm6\) & \(5.39\pm1.15\) & \(0.17\pm0.29\) & \(2.00\pm0.77\) & \(0.00\pm0.02\) \\
        SwiGLU value & \(248\pm6\) & \(16.14\pm0.79\) & \(0.67\pm0.85\) & \(7.16\pm2.43\) & \(0.00\pm0.04\) \\
        \bottomrule
    \end{tabular}%
    }

    \caption{\textbf{First-layer neurons align specifically with the true
    cluster-dependent predictive directions.}
    In the single-index experiment of Fig.~\ref{fig:local-sample-complexity},
    the fraction of active neurons aligned with a true predictive direction is
    substantially larger than the fraction aligned with
    random directions in the same predictive subspace. Results use \(K=10\) clusters and \(n=500{,}000\) training samples.
True-direction entries and active-neuron counts report mean \(\pm\) sample standard deviation over the
same three runs as Fig.~\ref{fig:local-sample-complexity}.
Random-direction entries report mean \(\pm\) sample standard deviation
over \(15{,}000\) pooled comparisons (\(5{,}000\) random direction sets per
trained model).
Appendix~\ref{app:singleindex-specialization-null}
    defines the active-neuron criterion and randomization protocol.}
    \label{tab:singleindex-mixed-k10-specialization-null}
\end{table}

\paragraph{Data setting.}
Following Section~\ref{sec:prelim}, we consider
\(K\in\{1,2,10,50\}\) clusters in \(d=20\). We divide the input into
ten cluster-identifying coordinates and ten predictive coordinates. The cluster centers are chosen as well-separated unit vectors $s_c\in\mathbb S^9:=\{s\in\mathbb R^{10}:\|s\|_2=1\}$. Details of their construction are provided in Appendix~\ref{appendix:main_exp_detail}. They are embedded as
\[
    \mu_c=
    \begin{bmatrix}
        s_c\\
        0
    \end{bmatrix}.
\]
The cluster-specific predictive directions are sampled independently and
uniformly from the unit sphere in the predictive subspace:
\[
    \widetilde v_c\sim\operatorname{Unif}(\mathbb S^9),
    \qquad
    v_c=
    \begin{bmatrix}
        0\\
        \widetilde v_c
    \end{bmatrix},
    \qquad
    \Sigma_c=
    \begin{pmatrix}
        \sigma_K^2 I_{10} & 0\\
        0 & I_{10}
    \end{pmatrix},
    \qquad
    \varepsilon\sim\mathcal N(0,0.02^2).
\]
Here, \(\sigma_K\) is adjusted with \(K\) so that the clusters remain well
separated. Full details of the center construction and the choice of
\(\sigma_K\) are provided in
Appendix~\ref{appendix:main_exp_detail}.

Each cluster is also assigned a link function $g_c$, sampled uniformly with
replacement from four choices: the normalized second-order Hermite polynomial,
the normalized third-order Hermite polynomial, the sine function, and the
hyperbolic tangent function. Thus, both the predictive direction $v_c$ and the
link function $g_c$ may vary across clusters.

\paragraph{Details and additional experiments.}
Appendix~\ref{appendix:main_exp_detail} provides the experimental protocol
and extends the comparison in Figure~\ref{fig:local-sample-complexity}
to include GELU and SwiGLU.
Appendix~\ref{app:singleindex-specialization-null} defines the
specialization measure and the random-direction baseline used to assess
alignment with the true predictive directions.
Appendix~\ref{app:importance-pruning} examines whether active neurons
retain predictive performance after pruning, through full and
readout-only fine-tuning.
Appendix~\ref{app:fixed-h3-single-index-details} considers a common
third-Hermite link function to match the setting of our theoretical analysis.



\subsection{Trained MLPs Encode Cluster Structure}
\label{sec:learned-gating}


\begin{figure}[H]
  \centering
  \experimentfigure[width=1.0\linewidth]{figs/kmeans_final_5.png}
  \caption{\textbf{Overview of the MLP-gated local-prediction procedure.} A ReLU
  MLP is trained without cluster labels, K-means is applied to its first-layer
  representations, and a local predictor is fit within each learned cluster.
  Test inputs are routed using the nearest learned centroid.}
  \label{fig:mlp-gating-procedure}
\end{figure}


In this experiment, we show that an MLP's first-layer representation captures both cluster identity and cluster-specific predictive structure. We train a ReLU MLP without cluster labels and construct representations using
only its active first-layer neurons. We apply K-means to these
representations. For each learned cluster, we use the corresponding active
neurons' first-layer weights to construct features for an independent
local Laplace or RFM predictor. At test time, an input is assigned to its
nearest learned centroid and evaluated by the corresponding local predictor.
Figure~\ref{fig:mlp-gating-procedure} summarizes this procedure.
Appendix~\ref{app:learned-gating-details} describes the active-neuron
selection criterion used in this experiment and provides the data,
training, representation-clustering, and evaluation protocols.



\paragraph{Main findings.}
Figure~\ref{fig:local-sample-complexity_kmeans} shows that using clusters
extracted from the MLP representation substantially improves both Laplace and
RFM over their global counterparts. In particular, MLP-gated RFM consistently
outperforms global RFM and approaches local RFM, which is fitted using the
true clusters. Thus, the
first-layer representation both separates the clusters and preserves their
local predictive features.

\begin{figure}[H]
    \centering
    \experimentfigure[width=\linewidth]
    {figs/mixed_h3_sin_tanh_active999_test10k_3seed_mean_plot_only.pdf}

    \vspace{0.3em}

    \experimentfigure[width=0.82\linewidth]
    {figs/mlp_gated_h3_sin_tanh_test10k_3seed_plain_legend_only.pdf}

    \caption{\textbf{Trained MLP representations encode cluster
    structure.}
    Clustering representations extracted from a trained ReLU MLP and fitting
    separate local predictors within the resulting groups improve both
    Laplace and RFM. In particular, MLP-gated RFM consistently outperforms
    global RFM and narrows the gap to local RFM, which is fitted using the
    true clusters. Curves show averaged loss over three runs with different cluster-specific link assignments, training and test examples, label noise, model and \(K\)-means initializations, and optimization minibatch orders.}
    \label{fig:local-sample-complexity_kmeans}
\end{figure}

\paragraph{Data setting.}
We use the clustered-data construction of
Section~\ref{sec:local-routing} with
\(K\in\{2,5,10,50\}\). Each cluster's link function is sampled uniformly from
the normalized third-Hermite polynomial, sine, and hyperbolic tangent. We
average results over three independent seeds and evaluate each run on
\(10{,}000\) test examples.






\subsection{MLPs Retain Their Advantage for Clustered Multi-Index Models}
\label{sec:clustered-multi-index}

We extend the experiment of Section~\ref{sec:local-routing} to clustered
multi-index models and show that MLPs retain their advantage over other methods.

\paragraph{Main findings.}
Figure~\ref{fig:clustered-multi-index-results} shows that MLPs retain their
advantage over global Laplace and RFM as the number of clusters increases.
The gated architectures generally perform best and approach local RFM at
larger sample sizes. Thus, the
benefit of cluster-dependent feature learning is not limited to
single-index targets. Table~\ref{tab:multiindex-mixed-k10-span-specialization} further shows that
active neurons align with cluster-specific predictive spans. Across all activations,
alignment with the true cluster-specific
predictive spans exceeds alignment with random three-dimensional
subspaces drawn within the entire predictive subspace.

\begin{figure}[H]
    \centering
    \includegraphics[width=\linewidth]
    {figs/multiindex_d50_final_n500k_3seed_original_all_activations_plot_only.pdf}
    \includegraphics[width=\linewidth]
    {figs/multiindex_d50_final_n500k_3seed_original_all_activations_legend_only.pdf}

    \caption{\textbf{MLPs retain their advantage for clustered multi-index
    targets.}
    As the number of clusters grows, MLPs, particularly ReGLU and SwiGLU,
    remain more sample efficient than global Laplace and RFM and approach
    local RFM at larger sample sizes. The local methods are given the true
    cluster identities and fit a separate predictor within each cluster. Curves show mean
    test MSE over three runs with different link-function assignments,
    data, label noise, model initialization, and optimization minibatches.
    Cluster centers and predictive spans are held fixed across these runs.
    The dotted line marks the observation-noise variance.}
    \label{fig:clustered-multi-index-results}
\end{figure}

\begin{table}[H]
    \centering
    \small
    \renewcommand{\arraystretch}{1.15}
    \setlength{\tabcolsep}{4pt}
    \resizebox{\linewidth}{!}{%
    \begin{tabular}{@{}lccccc@{}}
        \toprule
        & &
        \multicolumn{2}{c}{\textbf{Span cosine} \(\boldsymbol{\geq0.71}\)}
        & \multicolumn{2}{c}{\textbf{Span cosine} \(\boldsymbol{\geq0.90}\)}
        \\
        \cmidrule(lr){3-4}
        \cmidrule(lr){5-6}
        \textbf{Model/weight}
        & \shortstack{\textbf{Active neurons}\\\textbf{(out of 2048)}}
        & \shortstack{\textbf{True subspace}\\\textbf{(\%)}}
        & \shortstack{\textbf{Random subspaces}\\\textbf{(\%)}}
        & \shortstack{\textbf{True subspaces}\\\textbf{(\%)}}
        & \shortstack{\textbf{Random subspaces}\\\textbf{(\%)}}
        \\
        \midrule
        ReLU & \(1875\pm126\) & \(13.07\pm6.69\) & \(0.05\pm0.08\) & \(1.80\pm1.33\) & \(0.00\pm0.00\) \\
        GELU & \(661\pm332\) & \(11.35\pm5.62\) & \(0.12\pm0.20\) & \(4.68\pm3.47\) & \(0.00\pm0.00\) \\
        ReGLU gate & \(274\pm10\) & \(8.59\pm2.34\) & \(0.06\pm0.15\) & \(1.59\pm0.87\) & \(0.00\pm0.00\) \\
        ReGLU value & \(274\pm10\) & \(47.46\pm3.14\) & \(1.32\pm0.97\) & \(27.13\pm7.45\) & \(0.00\pm0.01\) \\
        SwiGLU gate & \(487\pm21\) & \(7.29\pm0.88\) & \(0.10\pm0.16\) & \(3.17\pm0.73\) & \(0.00\pm0.00\) \\
        SwiGLU value & \(487\pm21\) & \(12.70\pm1.67\) & \(0.12\pm0.19\) & \(4.75\pm1.09\) & \(0.00\pm0.00\) \\
        \bottomrule
    \end{tabular}%
    }
    \caption{\textbf{First-layer neurons specialize to cluster-specific
    predictive spans in the multi-index targets.}
    A substantially larger fraction of active neurons aligns with the true
    predictive spans than with random subspaces. Results use \(d=50\),
    \(K=10\) clusters, and \(n=500{,}000\) training samples.
    True-span entries and active-neuron counts report mean \(\pm\) sample standard
    deviation over the same three runs as
    Fig.~\ref{fig:clustered-multi-index-results}. Random-subspace entries
    report mean \(\pm\) sample standard deviation over \(15{,}000\) pooled
    comparisons (\(5{,}000\) random subspace collections per trained model).
    Appendix~\ref{app:multiindex-specialization-null} defines the
    active-neuron criterion, alignment measure, and random-subspace control.}
    \label{tab:multiindex-mixed-k10-span-specialization}
\end{table}

\paragraph{Data setting.}
We use the multi-index model defined in Section~\ref{sec:prelim}, with
\(d=50\), \(K\in\{1,2,10,25\}\), three predictive directions per cluster,
and noise standard deviation \(\sigma=0.02\). The input comprises twenty
cluster-identifying coordinates, twenty predictive coordinates, and ten
nuisance coordinates.
Each cluster independently receives one link function, sampled uniformly
from the normalized second-order Hermite polynomial, sine, and hyperbolic
tangent.
Only the local Laplace and RFM methods receive the true cluster identities.

Appendix~\ref{app:clustered-multi-index-details} provides the data,
training, evaluation, and specialization-analysis protocols.



\clearpage
\section{Theory}
\label{sec:specialization-theory}

In this section, we prove two main results describing how
MLPs learn data with cluster structure. In
Theorem~\ref{thm:neurons-specialize}, we prove that neurons in the MLP
specialize during training. In Theorem~\ref{thm:rfm-mlp-gap}, we prove that,
as the number of clusters tends to infinity, MLPs outperform standard kernel
methods and Recursive Feature Machines (RFM) in terms of sample complexity. This comes as a consequence of the fact that the latter two methods cannot compute specialized features for each
cluster.

\subsection{MLP Neurons Specialize When Learning on Gaussian Mixture Model Data}
\label{sec:early-specialization}

In this subsection, we show that, under small initialization,
neurons trained on well-separated Gaussian mixture data specialize to
cluster-specific predictive directions.


We consider $K$-cluster Gaussian mixture data in $d=2K$ dimensions. Let
$e_1,\ldots,e_K$ denote the standard basis of $\mathbb R^K$. We use the first
$K$ coordinates to separate the cluster means and the last $K$ coordinates
for the cluster-specific predictive directions.\footnote{Note that independent orthogonal changes of basis in the
two blocks preserve the isotropic Gaussian noise and all conclusions below.}

\begin{datasetting}[Symmetric Gaussian mixture]
\label{data:neuron-specialization}

Let $R > 0$ be a cluster separation parameter and, for $c\in[K]$, let
\[
    \mu_c=
    \begin{bmatrix}
        Re_c\\
        0
    \end{bmatrix},
    \qquad
    v_c=
    \begin{bmatrix}
        0\\
        e_c
    \end{bmatrix}.
\]
Set each cluster to be isotropic with $\Sigma_c=I_d$ and use the common link function $g_c=h_3$, where $h_3(t)=(t^3-3t)/\sqrt6$ is the third Hermite polynomial.
\\

In the notation of Section~\ref{sec:prelim}, the data distribution is
\[
    c\sim\operatorname{Unif}([K]),
    \qquad
    x\mid c\sim\mathcal N(\mu_c,\Sigma_c),
    \qquad
    y=g_c\!\left(\langle x,v_c\rangle\right)
      =h_3\!\left(\langle x,v_c\rangle\right).
\]
\end{datasetting}

Thus the routing coordinates encode cluster identity through the means
$Re_c$, where $R$ should be thought of as a large cluster-separation
parameter, so that the clusters are well separated, while within cluster $c$ the response depends only on the
cluster-specific predictive direction $v_c$. The cubic Hermite link function is chosen because there is an explicit expression for the expected product $\mathbb{E}[y\,\phi(\omega^\top x)]$ between the response and the ReLU neuron's activation, allowing us to characterize the directions to which neurons converge.
Next, we consider training a neural network to learn this data distribution.

\begin{trainingsetup}[Two-layer ReLU population gradient flow]
\label{train:neuron-specialization}
We train the two-layer ReLU network
\[
    f_\theta(x)=\sum_{j=1}^m a_j\phi(w_j^\top x),
    \qquad
    \phi(t)=t_+,
\]
by population gradient flow on the squared loss
\[
    \mathcal L(\theta)
    =
    \frac12\mathbb E\left[
        \bigl(f_\theta(x)-y\bigr)^2
    \right].
\]
\end{trainingsetup}


\begin{initializationsetting}[Small random hidden weights and zero output layer]
\label{init:neuron-specialization}

The hidden directions are initialized independently and uniformly at random, scaled by a parameter $\varepsilon > 0$,
while the output layer is initialized at zero:
\[
    \omega_j^0
    \stackrel{\mathrm{iid}}{\sim}
    \operatorname{Unif}(\mathbb S^{d-1}),
    \qquad
    w_j(0)=\varepsilon\omega_j^0,
    \qquad
    a_j(0)=0.
\]

\end{initializationsetting}

Under this initialization and data distribution, we are able to prove that the neurons in the MLP specialize to the clusters of the Gaussian mixture.
This result is consistent with, and provides theoretical support for, our empirical observations of neuron specialization in \Cref{sec:coordinate-specialization}.

\begin{theorem}[Randomly initialized neurons specialize]
\label{thm:neurons-specialize}
Under \cref{data:neuron-specialization}, \cref{train:neuron-specialization},
and \cref{init:neuron-specialization}, there are universal constants
$R_0,C<\infty$ such that the following holds
for every $R\geq R_0$. For almost every draw of
$\omega_1^0,\ldots,\omega_m^0$, there are cluster labels
$J_1,\ldots,J_m\in[K]$ and orientations
$\tau_1,\ldots,\tau_m\in\{\pm1\}$ such that, for every $\delta>0$, there
exist a finite time $T_\delta$ and $\varepsilon_0>0$ for which
\[
    0<\varepsilon\leq\varepsilon_0
    \quad\Longrightarrow\quad
    \left\|
        \frac{w_j(T_\delta)}{\|w_j(T_\delta)\|_2}
        -\tau_j v_{J_j}
    \right\|_2
    \leq
    \delta+\frac{C}{R}
    \qquad\text{for every }j\in[m].
\]
Thus every neuron specializes, up to orientation and a vanishing
$O(R^{-1})$ routing component, to the predictive direction of one cluster.

The selected labels $J_1,\ldots,J_m$ are independent and uniformly
distributed on $[K]$. Consequently,
\[
    \mathbb P\left(
        \text{every cluster is covered by a specialized neuron}
    \right)
    \geq
    1-Ke^{-m/K}.
\]
In particular, if
\[
    m\geq K\log\left(\frac{K}{\eta}\right),
\]
then all $K$ clusters are covered with probability at least $1-\eta$.
\end{theorem}

\paragraph{Proof sketch.}
We employ a proof strategy of studying feature learning
through effectively independent neuron dynamics, which has been used in prior
work on early-time and small-initialization regimes
\citep{abbe2022merged,min2024early,glasgow2024sgd}.
At initialization, the network output is zero and the hidden weights are
stationary. The initial derivative of each output weight is proportional to
the population correlation
\[
    \Phi(\omega)
    =\mathbb E\left[y\,
      \phi(\omega^\top x)\right]
\]
of its own randomly initialized hidden direction. Except on a measure-zero
set, this correlation is nonzero, so the output weight immediately acquires
the appropriate sign. ReLU homogeneity then turns the neuron's subsequent
directional dynamics into a positive time reparameterization of gradient
ascent on the corresponding signed correlation objective.

For the cubic Hermite link function, the Hermite calculation makes this objective a sum of
clusterwise cubic terms. Its positive local maxima are one-cluster solutions:
a neuron aligns with $\pm v_c$ and uses only an $O(R^{-1})$ component
in the associated routing direction to place its ReLU threshold within
cluster $c$. Indeed, we prove that any positive mixed-cluster critical point has an unstable
direction. Analytic-gradient-flow convergence and strict-saddle avoidance
therefore imply that a random neuron specializes almost surely.

Permutation symmetry makes the selected cluster uniform on $[K]$, and
independence of the initial hidden directions makes the selected labels
independent. The coverage estimate is then the usual coupon-collector union
bound. Finally, on every fixed early-time interval the full network output is
$O(m\varepsilon^2)$, so the coupled network dynamics are a vanishing
perturbation of the isolated-neuron dynamics. A Gr\"onwall argument transfers
the specialization result to population gradient flow. The complete proof is
given in Appendix~\ref{app:proof-neurons-specialize}.



\subsection{Specialization Yields a Sample-Complexity Gap Between MLPs and Adaptable Kernel Methods}
\label{sec:mlp-rfm-separation}

In this subsection, we establish a sample-complexity
separation between two-layer MLPs and the Recursive Feature Machines (RFM) adaptive kernel method. Our clustered data distribution is similar to that of the previous subsection, but we let the number of clusters $K$ tend to infinity.


We show that fitting a norm-constrained MLP with a
polynomial number of training samples yields a consistent estimator. Our proof
of this fact relies on showing that the MLP can have specialized neurons that
approximate the predictive function separately within each cluster. In
contrast, we show that RFM is not consistent under any fixed polynomial scaling of
the sample size in $K$.


Concretely, the data distribution we use for this result is similar to that in the previous subsection, but we suppose full separation between clusters for
ease of analysis.

\begingroup
\tcbset{theorysetting/.append style={breakable=false}}
\begin{datasetting}[Gaussian mixture with noiseless routing]
\label{data:rfm-mlp-gap}
We consider the same data distribution as in \cref{data:neuron-specialization}, except that the clusters can now be separated noiselessly. 
In other words, we set the means to $\mu_c=[e_c;0]$ and the predictive directions to
$v_c=[0;e_c]$ in $d=2K$ dimensions, and make the first $K$ coordinates
noiseless by setting
\[
    \Sigma_c=
    \begin{bmatrix}
        0_{K\times K} & 0\\
        0 & I_K
    \end{bmatrix}.
\]


Additionally, we consider a non-constant link function $g:\mathbb R\to\mathbb R$ which is independent of $K$, and satisfies boundedness and Lipschitzness
\[
    \|g\|_\infty\leq G_0,
    \qquad
    |g(s)-g(t)|\leq L|s-t|
    \quad\text{for all }s,t\in\mathbb R,
\]
for constants $G_0<\infty$ and $L>0$. We use the same link in every
cluster, $g_c=g$. Thus, in the notation of Section~\ref{sec:prelim},
\[
    c\sim\operatorname{Unif}([K]),
    \qquad
    x\mid c\sim\mathcal N(\mu_c,\Sigma_c),
    \qquad
    y=g\!\left(\langle x,v_c\rangle\right).
\]

\end{datasetting}
\endgroup

The routing block now reveals the cluster exactly, while the response within
cluster $c$ depends only on its cluster-specific predictive direction.

Instead of studying a network trained through gradient-based
updates, we study a two-layer MLP fit by empirical risk minimization subject
to a Frobenius-norm constraint on its weights.

\begingroup
\tcbset{theorysetting/.append style={breakable=false}}

\begin{trainingsetup}[Frobenius-constrained MLP ERM]
\label{train:mlp-gap}
We fit the two-layer ReLU network
\[
    f_\theta(x)=\sum_{j=1}^m a_j\phi(w_j^\top x),
    \qquad
    \phi(t)=t_+,
\]
using empirical risk minimization. For a width $m$ and budget $B>0$, define
\[
    \mathcal F_{B,m}
    :=
    \left\{
        f_\theta:
        \|\theta\|_{\mathrm F}^2
        =\frac12\sum_{j=1}^m
        \left(a_j^2+\|w_j\|_2^2\right)
        \leq B
    \right\}.
\]
Given $n$ samples, squared-loss empirical risk minimization (ERM) gives
\[
    \widehat f_{\mathrm{MLP}}
    \in
    \arg\min_{f\in\mathcal F_{B,m}}
    \frac1n\sum_{i=1}^n\bigl(f(x_i)-y_i\bigr)^2.
\]
\end{trainingsetup}

\endgroup

We compare this MLP ERM with standard rotationally invariant kernel ridge
regression and with RFM, a supervised kernel method
\citep{radhakrishnan2024mechanism,radhakrishnan2025linear}.\footnote{
We give the RFM comparator the ground-truth population average gradient outer
product (AGOP), thereby removing metric-estimation error and isolating the
limitation of using one global feature metric. An analysis of RFM using an
empirically estimated AGOP would require controlling empirical fluctuations
and is not implied by our theorem.}

\begingroup
\tcbset{theorysetting/.append style={breakable=false}}


\begin{trainingsetup}[Kernel methods and RFM with ground-truth AGOP]
\label{train:rfm-gap}

Let $\mathcal K$ be a rotationally invariant kernel. Given $n$ samples,
$\widehat f_{\mathrm{Kernel}}$ is kernel ridge regression with kernel
$\mathcal K(x,x')$. We give the RFM comparator the ground-truth population
AGOP\footnotemark
\[
    M
    :=\mathbb E\left[
        \nabla_x f^\star(x)\nabla_x f^\star(x)^\top
      \right].
\]
For $\rho\geq0$, define the regularized AGOP and corresponding kernel
\[
    M_\rho:=M+\rho I_d,
    \qquad
    \mathcal K_{M_\rho}(x,x')
    :=\mathcal K\bigl(\sqrt{M_\rho}x,\sqrt{M_\rho}x'\bigr).
\]
$\widehat f_{\mathrm{RFM}}$ is kernel ridge regression with kernel
$\mathcal K_{M_\rho}$. Both estimators may use any kernel ridge parameter
$\lambda_n\geq0$.

\end{trainingsetup}


\footnotetext{Because $P_x$ is supported on a union of affine subspaces,
the regression function $f^\star$ is intrinsically defined only on that
support, and its derivatives normal to the support are not determined.
Throughout, $\nabla_x f^\star(x)$ denotes the gradient along the affine
support component containing $x$, embedded in $\mathbb R^d$ with zero normal
component. Since $g$ is Lipschitz, it is differentiable almost everywhere,
and for an input $x$ from cluster $c$,
\[
    \nabla_x f^\star(x)
    =g'\!\left(\langle x,v_c\rangle\right)v_c
\]
almost everywhere.}

\endgroup


In the theorem below, we show that a two-layer MLP learns the
cluster-structured target with polynomial sample complexity for every fixed,
bounded, nonconstant Lipschitz link function $g$. In contrast, the standard rotationally
invariant kernel estimator and RFM with the ground-truth AGOP remain
inconsistent under every fixed polynomial sample-size scaling. The appendix
also proves the MLP upper bound in a more general fixed-dimensional
multi-index setting; the single-index result stated here is a
corollary for $r=1$, which is simpler to state.



\begin{theorem}[Specialization enables data-efficient learning]
\label{thm:rfm-mlp-gap}
Suppose \cref{data:rfm-mlp-gap} holds. Then the target function $f^\star$ admits approximation by cluster-specialized two-layer
ReLU networks with polynomially growing width and Frobenius norm,
and with approximation error tending to zero as $K\to\infty$.

Consequently, there are polynomially growing width
and norm-budget sequences $(m_K)$ and $(B_K)$ such that, with polynomially many samples $n_K=\Omega(K^{7})$,
the MLP estimator in \cref{train:mlp-gap} is consistent:
\[
    \lim_{K\to\infty}
    \mathbb E\left[
        \|\widehat f_{\mathrm{MLP}}-f^\star\|_{L^2(P_x)}^2
    \right]
    =0.
\]

On the other hand, for every fixed $A<\infty$, every sequence
$n_K=O(K^A)$, and every sequence $\rho_K\geq0$,
\[
    \liminf_{K\to\infty}
    \mathbb E\left[
        \|\widehat f-f^\star\|_{L^2(P_x)}^2
    \right]
    >0,
    \qquad
    \widehat f\in
    \left\{
        \widehat f_{\mathrm{Kernel}},
        \widehat f_{\mathrm{RFM}}
    \right\}.
\]
\end{theorem}






\paragraph{Proof sketch.}
The MLP claim follows from a more general multi-index result proved in
Appendix~\ref{app:proof-mlp-rfm-separation}. We first show that the link function can be approximated arbitrarily well
by a finite shallow network. By allowing the approximation accuracy to
improve sufficiently slowly with $K$, we obtain a sequence of
cluster-specialized MLPs whose approximation error tends to zero while
their width and Frobenius norm grow only polynomially in $K$. This construction follows by using neurons specialized to each cluster as follows. Since
$g$ is Lipschitz and bounded, it can be approximated in Gaussian $L^2$ by a
finite shallow network of ridge-ramp functions, each of which is a
difference of two ReLUs. We can thus approximate $g$ on each cluster by using routing
coordinates to convert every ridge-ramp unit into a pair of ReLU neurons whose
contributions cancel on all nonselected clusters.

The construction that we provide implies that there is a Frobenius-norm-bounded network with good approximation error. Thus, standard Barron- and path-norm Rademacher-complexity bounds control the
estimation error
\citep{e2022barron,bach2017breaking,golowich2018size}.
Because the Gaussian covariates are unbounded, we derive the required
sample-dependent version carefully, using Gaussian concentration to control
the MLP predictions on typical inputs.

For the kernel lower bound, expand the link function in terms of normalized Hermite
polynomials,
\[
    g(t)=\sum_{q=0}^\infty\widehat g_qh_q(t),
\]
where $h_q$ denotes the normalized probabilists' Hermite polynomial of degree
$q$, and $\widehat g_q$ is the corresponding scalar Hermite coefficient of $g$.
A bounded, continuous, nonconstant function cannot have only finitely many
nonzero Hermite coefficients: otherwise it would equal a polynomial, and a
bounded polynomial on $\mathbb R$ is constant. Thus every finite Hermite tail
of $g$ has positive energy.


The ground-truth population AGOP is a rescaling of
$P_V:=\sum_{c=1}^K v_cv_c^\top$, the orthogonal projector onto the span of
the cluster-specific predictive directions:
\[
    M=\frac{\alpha}{K}P_V,
    \qquad
    \alpha=\mathbb E[g'(G)^2]>0.
\]
Hence the standard kernel metric and every regularized AGOP metric remain
rotationally invariant within the predictive coordinates. At Hermite order
$q$, the clusterwise target contains a component in an irreducible harmonic
subspace of dimension $\Theta_q(K^q)$. The representer theorem restricts a
kernel predictor to the span of at most $n$ kernel sections, and the projected
span is an invariant random subspace. Therefore, when $n=O(K^A)$, every fixed
order $q>A$ leaves asymptotically all of its target energy unrecovered. Since
$g$ has nonzero Hermite energy beyond every finite order, both kernel
estimators retain nonzero error under every fixed polynomial sample scaling. This degree-by-degree harmonic obstruction is closely related to prior
high-dimensional analyses of rotationally invariant kernel methods and
invariant kernel models
\citep{ghorbani2021linearized,ghorbani2020outperform,mei2021invariances}.
The complete proof is given in
Appendix~\ref{app:proof-mlp-rfm-separation}.

\newpage




\section*{Acknowledgements}

We gratefully acknowledge support from the National Science Foundation (NSF)
under grants CCF-2112665 and MFAI 2502258, the Office of Naval Research
(ONR N000142412631), and the Defense Advanced Research Projects Agency
(DARPA) under Contract No. HR001125CE020.
This work used the Delta system at the National Center for Supercomputing Applications through allocation TG-CIS220009 from the Advanced Cyberinfrastructure Coordination Ecosystem: Services \& Support (ACCESS) program, which is supported by National Science Foundation grants \#2138259, \#2138286, \#2138307, \#2137603, and \#2138296.


AI tools were used to assist with aspects of the experiments and theoretical analysis.



\bibliographystyle{plainnat}
\bibliography{references_feature_learning}

\newpage
\appendix

\section{Additional Experimental Details}
\label{app:experimental-details}

\subsection{Additional Details for First-Layer Neuron Specialization}
\label{app:coordinate-specialization-details}

This appendix supplements Section~\ref{sec:coordinate-specialization}.
Section~\ref{app:specialization-measures-examples} motivates the complementary
specialization measures and illustrates them with representative ReLU and GELU
neurons. Section~\ref{app:shared-specialization-analysis} introduces the
active-neuron criterion and the shared analysis protocol.
Section~\ref{app:specialization-null} introduces controls based on random
directions and subspaces to assess specialization, and applies this
comparison to the axis-aligned experiment.
Finally, Section~\ref{app:hermite3-specialization-experiment} gives the
third-Hermite training protocol, and
Section~\ref{app:hermite2-specialization-experiment} presents the additional
second-Hermite experiment.

\subsubsection{Specialization Measures and Representative Neurons}
\label{app:specialization-measures-examples}

\paragraph{Why predictive-coordinate dominance?}
For an incoming weight vector \(w_j\), define its maximum
predictive-direction cosine as
\[
    A_j
    =
    \max_{1\leq c\leq K}
    \frac{|\langle w_j,v_c\rangle|}
    {\|w_j\|_2\,\|v_c\|_2}.
\]
Cosine alignment measures how much of a neuron's total weight norm lies in a
single predictive direction. It can therefore be conservative when one
predictive coordinate is large but many individually small coefficients
collectively contribute appreciably to the norm. Let \(\widehat c_j\) index
the predictive direction with the largest absolute inner product with
\(w_j\). Because \(v_c\) is supported on coordinate \(K+c\), define the
predictive-coordinate dominance ratio as
\[
    D_j
    =
    \frac{|[w_j]_{K+\widehat c_j}|}
    {\displaystyle
     \max_{\substack{1\leq \ell\leq d\\
                     \ell\neq K+\widehat c_j}}
     |[w_j]_\ell|}.
\]
Thus, \(D_j>1\) means that the largest coordinate is predictive, while
\(D_j>2\) means that this coordinate is more than twice as large as every
competing coordinate. Unlike cosine alignment, this ratio asks whether a
predictive coordinate dominates each alternative coordinate individually.
The two measures therefore provide complementary views of specialization.

\paragraph{ReLU examples.}
Figure~\ref{fig:low-cosine-dominance} illustrates this distinction for ReLU
neurons trained with the normalized third-Hermite link function. Although the
displayed
neurons have only moderate cosine alignment because their norms include many
smaller coefficients, each has a pronounced spike on one cluster-specific
predictive coordinate. Predictive-coordinate dominance therefore captures
clear specialization that a strict cosine threshold can miss.


\begin{figure}[H]
    \centering
    \includegraphics[width=\textwidth]
    {figs/relu_high_dominance_moderate_cosine_neurons.pdf}
    \caption{\textbf{Predictive-coordinate dominance complements cosine
    alignment.}
    This plot illustrates why predictive-coordinate dominance can be a more informative measure than cosine similarity with a strict threshold. A neuron may have only moderate cosine alignment because its norm includes many smaller coefficients, while still exhibiting a pronounced spike along a single cluster-specific predictive direction $v_c$. Recall that here, each $v_c$ is chosen as a standard basis vector, so concentration on a predictive coordinate is equivalent to alignment with the corresponding predictive direction.
}
    \label{fig:low-cosine-dominance}
\end{figure}

\paragraph{GELU examples.}
Figure~\ref{fig:gelu-hermite3-specialization-examples} shows the corresponding
behavior for GELU. The top row contains neurons with high cosine alignment,
whereas the bottom row contains neurons with moderate cosine alignment but a
single strongly dominant predictive coordinate. This confirms that the
distinction between the two specialization measures is not specific to ReLU.


\begin{figure}[H]
    \centering
    \includegraphics[width=1.0\textwidth]
    {figs/gelu_high_cosine_neurons.pdf}

    \vspace{0.25em}

    \includegraphics[width=1.0\textwidth]
    {figs/gelu_high_dominance_moderate_cosine_neurons.pdf}
    \caption{\textbf{Normalized third-Hermite specialization in GELU
    neurons.}
    Top: examples whose first-layer weights are nearly entirely concentrated
    on one predictive coordinate. Bottom: examples for which many smaller
    coefficients reduce cosine alignment, while one predictive coordinate
    remains more than twice as large as every competing coordinate.}
    \label{fig:gelu-hermite3-specialization-examples}
\end{figure}

\clearpage

\subsubsection{Active Neurons and Specialization}
\label{app:shared-specialization-analysis}

\paragraph{Active-neuron criterion.}
We restrict the specialization analysis to neurons that contribute materially
to the network output. For an ordinary ReLU or GELU neuron, we define its
importance as
\[
    I_j=|a_j|\,\|w_j\|_2,
\]
where \(a_j\) is its output weight and \(w_j\) is its incoming weight. For a
gated ReGLU or SwiGLU neuron, we use
\[
    I_j=|a_j|\,\|w_{g,j}\|_2\,\|w_{v,j}\|_2,
\]
where \(w_{g,j}\) and \(w_{v,j}\) denote its gate and value weights.

We sort neurons in decreasing order of importance \(I_j\), then add their
importance scores in that order until the cumulative sum reaches at least
\(99.9\%\) of the total importance, \(\sum_j I_j\), summed over all neurons.
We retain all neurons up to and including the one that reaches this threshold
and refer to them as \emph{active}. For gated architectures, the same active-neuron set is used
when analyzing the gate and value weights separately.

The \(99.9\%\) threshold is a descriptive sparsification rule rather than a
statistical significance threshold. Its purpose is to prevent the
specialization statistics from being dominated by neurons with negligible
influence on the network output.

\paragraph{Specialization statistics and reporting.}
For each active neuron, we compute the cosine-alignment score \(A_j\) and the
predictive-coordinate dominance score \(D_j\) defined above. Absolute inner
products are used because alignment with \(v_c\) and \(-v_c\) represents the
same predictive direction. Table~\ref{tab:coordinate-specialization} and the
top panel of Fig.~\ref{fig:coordinate-specialization-hermite2} report the
fractions of active neurons satisfying
\[
    A_j\geq0.71,\qquad
    A_j\geq0.90,\qquad
    D_j>1,\qquad
    D_j>2.
\]
For ReGLU and SwiGLU, all statistics are computed and reported separately for
the gate and value weights.

The individual neurons displayed in the bottom panel of
Fig.~\ref{fig:coordinate-specialization-hermite2} are the
highest-importance active GELU neurons satisfying \(A_j\geq0.90\). Their
coordinate-wise absolute incoming weights are shown so that specialization is
visible independently of sign. For the normalized third-Hermite link function,
Fig.~\ref{fig:coordinate-specialization-hermite3-neurons} and the top panel of
Fig.~\ref{fig:gelu-hermite3-specialization-examples} show high-alignment ReLU
and GELU examples, respectively. Figure~\ref{fig:low-cosine-dominance} and the
bottom panel of Fig.~\ref{fig:gelu-hermite3-specialization-examples} instead
show neurons with only moderate \(A_j\) but large \(D_j\), making the
complementary roles of the two measures explicit for both activations.

\subsubsection{Comparing Alignment with True and Random Directions}
\label{app:specialization-null}

In this section, we introduce a control experiment to determine whether the
observed neuron alignment reflects specialization to the true predictive
directions. If we replace these directions with random directions in the same
predictive subspace, would a similar fraction of active neurons still exceed
the same alignment threshold?

We answer this question by comparing alignment
with the true directions against alignment with randomly generated
alternatives. For multi-index models, we make the analogous comparison using
predictive subspaces. The random directions or subspaces match their true
counterparts in number, dimension, and coordinate support, while respecting
the geometric constraints of each experiment. A substantially larger fraction
of neurons aligned with the true directions or subspaces provides evidence
that the learned weights capture the specific predictive structure of the
problem.

We use this idea in the specialization experiments of
Sections~\ref{sec:coordinate-specialization}, \ref{sec:local-routing}, and
\ref{sec:clustered-multi-index}. The appropriate random object depends on the
experiment: it is a set of directions for the single-index models and a set
of three-dimensional subspaces for the multi-index model.

\paragraph{Shared post-hoc protocol.}
This is a post-hoc analysis: we first train each network and select its active
neurons using the criterion reported for that experiment. We then keep both
the learned weights and this active set fixed. For each trained seed and
alignment threshold, we compare the observed percentage of active neurons
exceeding the threshold with the corresponding percentages under
\(B=5{,}000\) random sets of directions or subspaces. Thus, the control changes
only the directions or subspaces against which the learned neurons are
evaluated; it does not retrain
or otherwise modify the network. For each trained seed \(s\in\{1,2,3\}\)
and threshold \(\tau\in\{0.71,0.90\}\), we compute the one-sided Monte Carlo
\(p\)-value
\[
    p_{\mathrm{MC},s}(\tau)
    =
    \frac{1+\sum_{b=1}^{B}
    \mathbf 1\!\left\{T_{s,\mathrm{null}}^{(b)}(\tau)\geq T_{s,\mathrm{obs}}(\tau)\right\}}
    {B+1}.
\]
Here, \(T_{s,\mathrm{obs}}(\tau)\) is the observed percentage of active
neurons meeting the threshold, and \(T_{s,\mathrm{null}}^{(b)}(\tau)\) is
the corresponding percentage for random set \(b\).
We report the largest per-seed value,
\[
    p_{\max}(\tau)=\max_{s\in\{1,2,3\}}p_{\mathrm{MC},s}(\tau).
\]
Thus, p-values are computed within each seed, not from the pooled random
percentages. If none of the random sets matches or exceeds the observed
percentage, the per-seed value is \(1/(B+1)=1/5001\approx2.0\times10^{-4}\),
the finite Monte Carlo resolution limit, not zero.
These dimensionless values describe conditional random-geometry comparisons
for fixed learned weights; they do not test the complete training procedure
or an improvement due to fine-tuning. We apply no multiplicity adjustment
across rows or thresholds.

Within each training seed, the same random directions or subspaces are used
for every architecture, so model comparisons are not affected by different
Monte Carlo draws. The coordinate-model construction is given next; the
non-axis-aligned single-index and multi-index constructions are given in
Sections~\ref{app:singleindex-specialization-null} and
\ref{app:multiindex-specialization-null}, respectively.

True-direction or true-span entries report means and sample standard
deviations over three trained models. Random-baseline entries report the
mean and sample standard deviation of all \(15{,}000\) specialization
percentages pooled across the three models and their \(5{,}000\) random
draws each. These are neither standard errors nor standard deviations of
seed-level means. Specialization entries are percentages, with standard
deviations in percentage points. Displayed values are rounded to two
decimal places; a displayed zero may therefore represent a small nonzero
value.

\paragraph{Application to the axis-aligned experiment.}
For the experiment in Section~\ref{sec:coordinate-specialization}, we use the
active-neuron criterion defined in
Section~\ref{app:shared-specialization-analysis}. The true predictive
directions form an orthonormal coordinate basis of the ten-dimensional
predictive subspace, so the randomized directions must preserve this joint
geometry rather than being sampled independently.
For each null draw, we sample a Haar-random orthogonal matrix
\(Q\in\mathbb R^{10\times10}\) and embed its columns in the same ten
predictive coordinates as the true directions. The randomized directions are
therefore mutually orthonormal and span exactly the same predictive subspace
as the true directions. This construction preserves each learned weight's
norm and its total mass in the predictive subspace, while breaking its
alignment with the particular direction assigned to each cluster.

\begin{table*}[t]
    \centering
    \small
    \renewcommand{\arraystretch}{1.12}
    \setlength{\tabcolsep}{5pt}
    \resizebox{\textwidth}{!}{%
    \begin{tabular}{@{}lccccc@{}}
        \toprule
        \textbf{Model/weight}
        & \shortstack{\textbf{Active neurons}\\\textbf{(out of 2048)}}
        & \multicolumn{2}{c}{\textbf{Cosine \(\mathbf{\geq .71}\) (\%)}}
        & \multicolumn{2}{c}{\textbf{Cosine \(\mathbf{\geq .90}\) (\%)}} \\
        \cmidrule(lr){3-4}\cmidrule(lr){5-6}
        & & \textbf{True directions} \(\mathbf{\uparrow}\)
          & \textbf{Random directions}
          & \textbf{True directions} \(\mathbf{\uparrow}\)
          & \textbf{Random directions} \\
        \midrule
        ReLU & \(469\pm57\) & \(56.75\pm3.91\) & \(5.80\pm3.96\) & \(33.22\pm1.92\) & \(0.03\pm0.24\) \\
        GELU & \(1{,}747\pm83\) & \(30.15\pm2.48\) & \(2.56\pm2.05\) & \(18.40\pm1.20\) & \(0.02\pm0.12\) \\
        ReGLU gate & \(200\pm56\) & \(51.89\pm1.22\) & \(5.58\pm4.33\) & \(33.82\pm1.27\) & \(0.05\pm0.37\) \\
        ReGLU value & \(200\pm56\) & \(24.98\pm0.52\) & \(1.08\pm1.07\) & \(3.94\pm2.24\) & \(0.00\pm0.03\) \\
        SwiGLU gate & \(432\pm5\) & \(25.06\pm2.29\) & \(3.30\pm1.80\) & \(12.26\pm1.27\) & \(0.02\pm0.11\) \\
        SwiGLU value & \(432\pm5\) & \(17.88\pm2.88\) & \(2.83\pm1.19\) & \(4.62\pm0.96\) & \(0.01\pm0.04\) \\
        \bottomrule
    \end{tabular}%
    }
    \caption{\textbf{Predictive-direction specialization exceeds a
    geometry-matched random-direction baseline.}
    Observed alignment with the true cluster-specific directions is
    substantially stronger than alignment with randomized orthonormal
    directions in the same predictive subspace. True-direction entries are
    means and sample standard deviations over three trained models.
    Random-direction entries summarize \(15{,}000\) pooled comparisons
    (\(5{,}000\) random bases for each trained model).
    Every row and threshold has \(p_{\max}=1/5001\), as defined in the shared
    protocol in Section~\ref{app:specialization-null}.}
    \label{tab:coordinate-specialization-null}
\end{table*}

We apply this randomization only to cosine-based specialization. The
coordinate-dominance statistic depends on the chosen coordinate system and
therefore does not have a rotation-invariant interpretation under this null.

\clearpage

\subsubsection{Normalized Third-Hermite Experiment}
\label{app:hermite3-specialization-experiment}

Let \(\operatorname{He}_3\) denote the third probabilists' Hermite
polynomial. We use its unit-variance normalization,
\[
    \operatorname{He}_3(t)=t^3-3t,
    \qquad
    h_3(t)
    =
    \frac{\operatorname{He}_3(t)}{\sqrt{3!}}
    =
    \frac{t^3-3t}{\sqrt6}.
\]
For an example from cluster \(c\), the response is
\[
    y
    =
    h_3\!\left(\sqrt d\,\langle v_c,x\rangle\right)
    +\varepsilon,
    \qquad
    \varepsilon\sim\mathcal N(0,0.005^2).
\]
Under the data distribution in Section~\ref{sec:coordinate-specialization},
\(\langle v_c,x\rangle\) has conditional variance \(1/d\). Thus,
\(\sqrt d\,\langle v_c,x\rangle\) has unit variance, and for
\(Z\sim\mathcal N(0,1)\), the normalization above gives
\(\mathbb E[h_3(Z)^2]=1\).

\paragraph{Training protocol.}
All four architectures use width \(2{,}048\), a batch size of \(2{,}048\),
trainable biases, the width-aware parameterization, and
\(\mu\mathrm{P}\)-scaled Adam. Each run uses \(100{,}000\) training examples,
\(4{,}096\) validation examples, and \(2{,}048\) test examples. The
activation-specific optimization configurations are given in
Table~\ref{tab:hermite3-training-protocol}.

\begin{table}[H]
    \centering
    \small
    \caption{Training configurations for the normalized third-Hermite
    specialization experiment. Each learning rate follows a cosine schedule
    from the first value to the second.}
    \label{tab:hermite3-training-protocol}
    \begin{tabularx}{\textwidth}{@{}lXXX@{}}
        \toprule
        Activation & Cosine LR & Weight decay & Maximum steps \\
        \midrule
        ReLU   & \(2\times10^{-3}\to2\times10^{-4}\)     & \(10^{-6}\) & \(50{,}000\) \\
        GELU   & \(1.6\times10^{-2}\to1.6\times10^{-3}\) & \(0\)       & \(20{,}000\) \\
        ReGLU  & \(1.6\times10^{-3}\to1.6\times10^{-4}\) & \(0\)       & \(20{,}000\) \\
        SwiGLU & \(2\times10^{-3}\to2\times10^{-4}\)     & \(0\)       & \(50{,}000\) \\
        \bottomrule
    \end{tabularx}
\end{table}

For every run, we restore the
checkpoint with the lowest validation MSE and use the test set only after
model selection. The specialization statistics in
Table~\ref{tab:coordinate-specialization} are reported as the mean and
standard deviation over three independent seeds.

\clearpage

\subsubsection{Normalized Second-Hermite Experiment}
\label{app:hermite2-specialization-experiment}

Figure~\ref{fig:coordinate-specialization-hermite2} provides an additional
experiment using a normalized second-Hermite link function. First-layer
neurons remain strongly specialized across all four activation families, with
the ReGLU gate weights exhibiting especially strong alignment. Thus,
specialization also occurs beyond the third-Hermite setting analyzed in
Section~\ref{sec:early-specialization}.




\paragraph{Experimental setting and link function.}
This experiment uses exactly the same data geometry, input distribution, noise level, model architectures, and width as Section~\ref{sec:coordinate-specialization}. Within the data-generating setup, the only change is the link function: in every cluster, we replace the normalized third-Hermite link function with the unit-variance normalized second-Hermite link function. Optimization hyperparameters are specified separately below.
\[
    h_2(t)=\frac{\operatorname{He}_2(t)}{\sqrt{2!}}
          =\frac{t^2-1}{\sqrt{2}},
    \qquad
    \operatorname{He}_2(t)=t^2-1.
\]
Thus, for an example from cluster \(c\),
\[
    y=h_2\!\left(\sqrt d\,\langle v_c,x\rangle\right)+\varepsilon,
    \qquad
    \varepsilon\sim\mathcal N(0,0.005^2).
\]
For \(Z\sim\mathcal N(0,1)\), this normalization gives
\(\mathbb E[h_2(Z)^2]=1\), matching the scale of the normalized
third-Hermite link function. We generate the training, validation, and test
sets independently; they contain \(100{,}000\), \(4{,}096\), and \(2{,}048\)
examples, respectively.


\begin{figure}[H]
    \centering

    \small
    \renewcommand{\arraystretch}{1.12}
    \setlength{\tabcolsep}{5pt}

    \resizebox{\textwidth}{!}{%
    \begin{tabular}{@{}lccccc@{}}
        \toprule
        & & \multicolumn{4}{c}{\textbf{Specialization measures
        (\% of active neurons)}} \\
        \cmidrule(lr){3-6}
        \textbf{Model/weight}
        & \shortstack{\textbf{Active}\\\textbf{neurons}\\\textbf{(out of 2048)}}
        & \(\mathbf{A_j\geq .71}\,\uparrow\)
        & \(\mathbf{A_j\geq .90}\,\uparrow\)
        & \(\mathbf{D_j>1}\,\uparrow\)
        & \(\mathbf{D_j>2}\,\uparrow\) \\
        \midrule
        ReLU         & 613   & 45.5 & 33.4 & 71.1 & 49.1 \\
        GELU         & 1,710 & 49.9 & 28.2 & 68.2 & 43.1 \\
        ReGLU gate   & 453   & 73.1 & 56.1 & 96.7 & 73.5 \\
        ReGLU value  & 453   & 28.5 &  3.8 & 69.1 & 24.7 \\
        SwiGLU gate  & 551   & 43.6 & 19.1 & 91.7 & 38.1 \\
        SwiGLU value & 551   & 46.3 & 14.9 & 87.5 & 41.4 \\
        \bottomrule
    \end{tabular}%
    }

    \vspace{1em}

    \includegraphics[width=\textwidth]{figs/top3_neurons_gelu.png}

    \caption{\textbf{Strong neuron specialization also occurs for the normalized
    second-Hermite link function.}
    Top: percentage of specialized neurons among active neurons. Across
    standard and gated activations, many neurons align strongly with a single
    cluster-specific predictive direction. Bottom: absolute first-layer
    weights of selected GELU neurons. Each neuron concentrates on one
    predictive coordinate, with comparatively little weight on the
    cluster-position and nuisance coordinates.}
    \label{fig:coordinate-specialization-hermite2}
\end{figure}


\paragraph{Training protocol.}
All models use Adam, width \(2048\), batch size \(2048\), zero weight decay,
and the width-aware parameterization of
Section~\ref{app:hermite3-specialization-experiment}. The cosine
learning-rate schedules are
\begin{center}
\begin{tabular}{@{}r@{\;}l@{}}
    ReLU:   & \(4\times10^{-3}\to4\times10^{-4}\), \\
    GELU:   & \(1.6\times10^{-2}\to1.6\times10^{-3}\), \\
    ReGLU:  & \(1.6\times10^{-3}\to1.6\times10^{-4}\), \\
    SwiGLU: & \(2.56\times10^{-2}\to2.56\times10^{-3}\).
\end{tabular}
\end{center}
ReLU is trained for at most \(40{,}000\) steps; the remaining models use
at most \(20{,}000\) steps. Each reported model uses its lowest-validation-MSE
checkpoint, at step \(37{,}000\) for ReLU. Test data are used only for
final evaluation.

\FloatBarrier

\subsection{Additional Details for MLPs Jointly Learning Cluster Structure
and Cluster-Specific Predictive Functions}
\label{app:local-routing-details}

This appendix provides the experimental details and additional results
supporting Section~\ref{sec:local-routing}.
Section~\ref{app:singleindex-specialization-null} details the specialization
measure and random-direction control used in
Table~\ref{tab:singleindex-mixed-k10-specialization-null}.
Section~\ref{app:importance-pruning} evaluates whether active neurons
retain predictive performance after pruning, comparing fine-tuning the
active neurons with updating only the readout while keeping the
hidden parameters fixed.
Section~\ref{appendix:main_exp_detail} gives the complete protocol
for the single-index experiment with mixed link functions and extends the
comparison in Fig.~\ref{fig:local-sample-complexity} to all four MLP
activations---ReLU, GELU, ReGLU, and SwiGLU---by adding GELU and SwiGLU
(Fig.~\ref{fig:mixed_link_functions_all_activations}).
Finally, Section~\ref{app:fixed-h3-single-index-details} presents the
experiment in Fig.~\ref{fig:fixed-h3-single-index-results}, which uses only
the normalized third-order Hermite link function, as in our theoretical
analysis in Section~\ref{sec:early-specialization}.

\subsubsection{Neuron Specialization and Random-Direction Control}
\label{app:singleindex-specialization-null}

Table~\ref{tab:singleindex-mixed-k10-specialization-null} complements the
predictive-performance results in Fig.~\ref{fig:local-sample-complexity} by
testing whether first-layer neurons align with the true cluster-specific
predictive directions. The analysis uses the \(K=10\), \(n=500{,}000\)
models and the active-neuron criterion defined in
Section~\ref{app:shared-specialization-analysis}. For gated architectures, the
same active set is used to analyze gate and value weights separately.

\paragraph{Specialization measure.}
For each active first-layer weight \(w_j\), we compute its maximum absolute
cosine with the ten true predictive directions,
\[
    \max_{1\leq c\leq K}
    \frac{|\langle w_j,v_c\rangle|}
    {\|w_j\|_2\,\|v_c\|_2}.
\]
The observed columns in
Table~\ref{tab:singleindex-mixed-k10-specialization-null} report the percentage
of active neurons whose maximum cosine is at least \(0.71\) or \(0.90\).
The denominator uses the full twenty-dimensional weight norm.

\paragraph{Random-direction control.}
We use the shared post-hoc protocol of
Section~\ref{app:specialization-null}. Unlike the coordinate experiment in
Section~\ref{sec:coordinate-specialization}, the true directions here are
independently sampled, are not axis-aligned, and are not constrained to be
mutually orthogonal. We therefore sample random directions from the same
distribution as the true directions. For each trained model, we generate
\(B=5{,}000\) independent sets of random directions, indexed by
\(b\in\{1,\ldots,B\}\). Within each set \(b\), we independently sample one
direction for each cluster \(c\):
\[
    \widetilde q_c^{(b)}\sim\operatorname{Unif}(\mathbb S^9),
    \qquad
    q_c^{(b)}
    =
    \begin{bmatrix}
        0\\
        \widetilde q_c^{(b)}
    \end{bmatrix}
    \in\mathbb R^{20}.
\]
For every active weight, we replace the maximum cosine with the true
directions by the maximum cosine with
\(q_1^{(b)},\ldots,q_K^{(b)}\), and then recompute the percentage exceeding
each threshold. This experiment-specific construction preserves the sampling
distribution of the predictive directions in addition to the properties
preserved by the shared protocol.

We aggregate the results as specified in
Section~\ref{app:specialization-null}. At both thresholds, alignment with the true directions is
substantially stronger than alignment with the randomized directions for
every architecture and weight type.
For every row and threshold in
Table~\ref{tab:singleindex-mixed-k10-specialization-null}, the maximum
per-seed Monte Carlo p-value is \(p_{\max}=1/5001\), using the shared
definition in Section~\ref{app:specialization-null}.


\subsubsection{Importance-Based Pruning and Fine-Tuning}
\label{app:importance-pruning}

The active-neuron criterion in
Section~\ref{app:shared-specialization-analysis} is based on weight
importance. We evaluate the selected neurons' predictive role through two pruning
experiments: fine-tuning only the active neurons recovers approximately
the full network's test MSE with similar neuron specialization. In the
second experiment, we freeze both the first-layer weights and biases and
fine-tune only the readout. This recovers much of the performance lost
through pruning without changing the retained neurons' directional alignment.
In both experiments, fine-tuning uses exactly the same training data as
the original model, and evaluation uses the same held-out test data.

\paragraph{Fine-tuning the retained neurons.}
We retain and fine-tune only the active neurons.
Table~\ref{tab:importance-pruning-performance} shows that fine-tuning
recovers similar or better test MSE than the full model, with similar
neuron specialization (Table~\ref{tab:importance-pruning-alignment}).

\begin{table}[H]
    \centering
    \small
    \renewcommand{\arraystretch}{1.15}
    \setlength{\tabcolsep}{5pt}
    \resizebox{\linewidth}{!}{%
    \begin{tabular}{@{}lcccc@{}}
        \toprule
        & & \multicolumn{3}{c}{\textbf{Test MSE} \(\boldsymbol{(\times10^{-3})}\,\downarrow\)} \\
        \cmidrule(lr){3-5}
        \textbf{Model}
        & \shortstack{\textbf{Active neurons}\\\textbf{(out of 4096)}}
        & \textbf{Full model} & \textbf{Pruned} & \textbf{Fine-tuned} \\
        \midrule
        ReLU   & \(558\pm36\) & \(1.97\pm0.30\) & \(2.09\pm0.29\) & \(1.94\pm0.29\) \\
        GELU   & \(287\pm15\) & \(1.42\pm0.30\) & \(3.93\pm0.75\) & \(1.38\pm0.32\) \\
        ReGLU  & \(133\pm44\) & \(1.54\pm0.31\) & \(2.02\pm0.19\) & \(1.49\pm0.38\) \\
        SwiGLU & \(248\pm6\)  & \(0.74\pm0.13\) & \(0.92\pm0.09\) & \(0.73\pm0.14\) \\
        \bottomrule
    \end{tabular}%
    }
    \caption{\textbf{Fine-tuning active neurons recovers performance after pruning.}
    Removing the remaining neurons increases test error, but fine-tuning
    all retained parameters restores approximately the full model's performance.
    Results use the same \(K=10\), \(n=500{,}000\) models as
    Table~\ref{tab:singleindex-mixed-k10-specialization-null}. Entries report
    mean \(\pm\) sample standard deviation over the same three runs.}
    \label{tab:importance-pruning-performance}
\end{table}

\begin{table}[H]
    \centering
    \small
    \renewcommand{\arraystretch}{1.15}
    \setlength{\tabcolsep}{5pt}
    \resizebox{\linewidth}{!}{%
    \begin{tabular}{@{}lcccc@{}}
        \toprule
        & \multicolumn{2}{c}{\textbf{Cosine} \(\boldsymbol{\geq0.71}\) \textbf{(\%)}}
        & \multicolumn{2}{c}{\textbf{Cosine} \(\boldsymbol{\geq0.90}\) \textbf{(\%)}} \\
        \cmidrule(lr){2-3}\cmidrule(lr){4-5}
        \textbf{Model/weight}
        & \textbf{Before} & \textbf{After} & \textbf{Before} & \textbf{After} \\
        \midrule
        ReLU         & \(12.72\pm3.54\) & \(12.18\pm3.57\) & \(3.02\pm0.79\)  & \(3.14\pm0.73\) \\
        GELU         & \(16.87\pm0.79\) & \(17.21\pm0.79\) & \(8.01\pm0.66\)  & \(8.01\pm0.66\) \\
        ReGLU gate   & \(16.87\pm4.31\) & \(16.87\pm4.31\) & \(5.02\pm1.31\)  & \(4.99\pm1.25\) \\
        ReGLU value  & \(79.65\pm4.15\) & \(79.65\pm4.15\) & \(60.12\pm6.44\) & \(59.92\pm5.63\) \\
        SwiGLU gate  & \(5.39\pm1.15\)  & \(5.39\pm0.93\)  & \(2.00\pm0.77\)  & \(1.87\pm0.98\) \\
        SwiGLU value & \(16.14\pm0.79\) & \(16.14\pm0.79\) & \(7.16\pm2.43\)  & \(7.30\pm2.67\) \\
        \bottomrule
    \end{tabular}%
    }
    \caption{\textbf{Active neurons retain similar specialization after fine-tuning.}
    Percentages are measured among the same active neurons before and after
    fine-tuning all retained parameters, using alignment with the original predictive directions.
    Entries report mean \(\pm\) sample standard deviation over three runs.
    The before-fine-tuning values and retained-neuron counts match
    Table~\ref{tab:singleindex-mixed-k10-specialization-null}.
    Every row, stage, and threshold has \(p_{\max}=1/5001\) against its
    fixed-weight random-direction baseline; these are not tests of the
    before-to-after change.}
    \label{tab:importance-pruning-alignment}
\end{table}

\paragraph{Readout-only fine-tuning after pruning.}
To isolate whether the active neurons already provide useful predictive
features, we again retain only the active neurons, freeze both their first-layer
weights and biases, and update only the output weights and output bias.
For ReGLU and SwiGLU, both the gate and value branches are frozen.
Table~\ref{tab:active-neuron-readout-only} shows that readout adjustment
recovers much of the performance lost through pruning.

\begin{table}[H]
    \centering
    \small
    \renewcommand{\arraystretch}{1.15}
    \setlength{\tabcolsep}{5pt}
    \resizebox{\linewidth}{!}{%
    \begin{tabular}{@{}lcccc@{}}
        \toprule
        & & \multicolumn{3}{c}{\textbf{Test MSE} \(\boldsymbol{(\times10^{-3})}\,\downarrow\)} \\
        \cmidrule(lr){3-5}
        \textbf{Model}
        & \shortstack{\textbf{Active neurons}\\\textbf{(out of 4096)}}
        & \textbf{Full model} & \textbf{Pruned} & \shortstack{\textbf{Readout-only}\\\textbf{fine-tuning}} \\
        \midrule
        ReLU   & \(558\pm36\) & \(1.97\pm0.30\) & \(2.09\pm0.29\) & \(1.94\pm0.27\) \\
        GELU   & \(287\pm15\) & \(1.42\pm0.30\) & \(3.93\pm0.75\) & \(1.53\pm0.31\) \\
        ReGLU  & \(133\pm44\) & \(1.54\pm0.31\) & \(2.02\pm0.19\) & \(1.57\pm0.28\) \\
        SwiGLU & \(248\pm6\)  & \(0.74\pm0.13\) & \(0.92\pm0.09\) & \(0.79\pm0.11\) \\
        \bottomrule
    \end{tabular}%
    }
    \caption{\textbf{Readout adjustment recovers much of the performance lost through pruning.}
    Updating only the output weights and output bias reduces test error
    while keeping all retained hidden parameters fixed.
    Results use the same \(K=10\), \(n=500{,}000\) models and active-neuron
    sets as Table~\ref{tab:importance-pruning-performance}.
    Entries report mean \(\pm\) sample standard deviation over the same
    three data and model-initialization runs.}
    \label{tab:active-neuron-readout-only}
\end{table}

\paragraph{Shared experimental setting and pruning.}
Both experiments use the same trained models as
Table~\ref{tab:singleindex-mixed-k10-specialization-null}: ReLU, GELU,
ReGLU, and SwiGLU networks of width \(4096\), with \(d=20\), \(K=10\),
and \(n=500{,}000\). The data and original training protocols are given
in Section~\ref{appendix:main_exp_detail}.

In both experiments, we retain the active neurons defined in
Section~\ref{app:shared-specialization-analysis}, using the same set as
the specialization analysis in Section~\ref{app:singleindex-specialization-null}.
This is the smallest importance-ranked set carrying at least \(99.9\%\)
of total importance.
For gated architectures, each retained unit includes both its gate and
value branches. Pruning removes all other neurons while preserving the
retained parameters and the original width-dependent output scaling.
The selected neuron identities remain fixed during fine-tuning; we do not
reapply the importance criterion afterward.
Each fine-tuning run starts from the original trained model immediately
after pruning, not from a checkpoint fine-tuned in the other experiment.

\paragraph{Shared fine-tuning protocol.}
Both experiments use the original \(500{,}000\) training examples.
A fresh Adam optimizer runs
for \(20{,}000\) steps with minibatches of \(4096\) examples sampled with
replacement. Cosine decay reduces the learning rate from \(10^{-3}\) to
\(10^{-5}\) for ReLU, GELU, and SwiGLU, and from \(10^{-4}\) to \(10^{-6}\)
for ReGLU. Weight decay is unchanged from the original models:
\(10^{-6}\) for ReLU and GELU, \(10^{-5}\) for ReGLU, and zero for SwiGLU,
applied only to trainable parameters.
For the readout-only control, Adam uses
\(\beta_1=0.9\), \(\beta_2=0.999\), and \(\epsilon=10^{-8}\).
We save checkpoints and evaluate full-training MSE every \(1000\) steps,
including the pruned model at step zero. We select the checkpoint with
the lowest training MSE; when all retained parameters are fine-tuned,
this is the final checkpoint in all twelve runs.
Neither validation nor test data are used for fine-tuning or checkpoint
selection. Evaluation uses the original \(4096\) test examples per seed;
results summarize the same three runs.

\paragraph{Specialization evaluation.}
We use the alignment measure in
Section~\ref{app:singleindex-specialization-null}, with the original
predictive directions and the full \(20\)-dimensional weight norm.
The thresholds are \(0.71\) and \(0.90\), as in
Table~\ref{tab:singleindex-mixed-k10-specialization-null}.
Before and after fine-tuning, we evaluate the same retained neurons,
with gate and value weights evaluated separately. Pruning itself leaves
their alignment unchanged because it preserves their incoming weights.

For the all-parameter fine-tuning experiment, we apply the random-direction
construction of Section~\ref{app:singleindex-specialization-null} separately
before and after fine-tuning, holding each stage's weights fixed. Both stages
use the same \(5{,}000\) random direction sets per seed and the same retained
neuron identities, without reselection. We aggregate the per-seed p-values
using the shared definition of \(p_{\max}\) in
Section~\ref{app:specialization-null}; the results appear in
Table~\ref{tab:importance-pruning-alignment}.

\FloatBarrier
\subsubsection{Experimental Details for the Single-Index Experiment with Mixed Link Functions}
\label{appendix:main_exp_detail}

\begin{figure*}[t]
    \centering
    \includegraphics[width=\textwidth]
    {figs/singleindex_final_n500k_3seed_original_all_activations_plot_only.pdf}

    \par\vspace{+0.2em}

    \includegraphics[width=1.0\textwidth]
    {figs/singleindex_final_n500k_3seed_original_all_activations_legend_only.pdf}

    \vspace{-0.3em}
    \caption{\textbf{Additional activations for the single-index model with
    mixed link functions.}
    Same plot as Figure~\ref{fig:local-sample-complexity}, with GELU and SwiGLU added. As the number of clusters grows, MLPs, especially those with gated activations such as ReGLU and SwiGLU, maintain strong performance close to local Laplace and RFM, while global RFM becomes increasingly less sample efficient. Curves show loss averaged over three runs with different link-function assignments, data, label noise, model initialization, and optimization minibatches.
    }
    \label{fig:mixed_link_functions_all_activations}
\end{figure*}


\paragraph{Cluster and predictive geometry.}
We use \(d=20\), with the first ten coordinates containing cluster-identifying
information and the final ten coordinates containing the predictive signal.
We consider
\[
    K\in\{1,2,10,50\}.
\]
The finite datasets contain equal numbers of examples from each cluster
whenever the sample size is divisible by \(K\), and otherwise differ by at
most one example per cluster.

For each \(K\), cluster-center directions are obtained by farthest-point
sampling from \(\max\{20{,}000,500K\}\) independent Gaussian vectors
normalized to unit norm, giving
\[
    s_1,\ldots,s_K\in\mathbb S^9,
    \qquad
    \mathbb S^9
    =
    \{s\in\mathbb R^{10}:\|s\|_2=1\}.
\]

For \(K>1\), define
\[
    \rho_K
    =
    \max_{c\neq c'}s_c^\top s_{c'}
\]
and set
\[
    \mu_c
    =
    \begin{bmatrix}
        s_c\\
        0
    \end{bmatrix}
    \in\mathbb R^{20}.
\]
For \(K=1\), we set \(\mu_1=0\), since no cluster discrimination is required.

The standard deviation of the noise in the first ten coordinates is
\begin{equation}
    \sigma_K
    =
    \frac{1-\rho_K}{2\gamma},
    \qquad
    \gamma=4,
    \label{eq:local-routing-noise}
\end{equation}
where we take \(\rho_1=0\). Thus, the cluster noise adapts to the
least-separated pair of centers and keeps the clusters well separated as
\(K\) changes.

Each cluster independently receives a unit predictive direction
\[
    \widetilde v_c
    \sim
    \operatorname{Unif}(\mathbb S^9),
    \qquad
    v_c
    =
    \begin{bmatrix}
        0\\
        \widetilde v_c
    \end{bmatrix}
    \in\mathbb R^{20}.
\]
The cluster-identifying and predictive coordinates are therefore disjoint.


\paragraph{Input distribution.}
Given cluster \(c\), an input is sampled as
\begin{equation}
    x
    =
    \mu_c+
    \begin{bmatrix}
        \sigma_K z_{\mathrm{cluster}}\\
        z_{\mathrm{pred}}
    \end{bmatrix},
    \qquad
    z_{\mathrm{cluster}},z_{\mathrm{pred}}
    \stackrel{\mathrm{iid}}{\sim}
    \mathcal N(0,I_{10}).
    \label{eq:local-routing-data}
\end{equation}
Equivalently,
\[
    x\mid C=c
    \sim
    \mathcal N(\mu_c,\Sigma_c),
    \qquad
    \Sigma_c
    =
    \begin{pmatrix}
        \sigma_K^2 I_{10} & 0\\
        0 & I_{10}
    \end{pmatrix}.
\]


\paragraph{Cluster-specific link functions.}
Each cluster independently receives a link function sampled uniformly with
replacement from
\begin{equation}
    \mathcal G
    =
    \left\{
        t\mapsto\frac{t^2-1}{\sqrt{2}},
        \quad
        t\mapsto\frac{t^3-3t}{\sqrt{6}},
        \quad
        t\mapsto\sin(t),
        \quad
        t\mapsto\tanh(t)
    \right\}.
    \label{eq:local-link-function-family}
\end{equation}
The first two functions are the normalized probabilists' Hermite polynomials
of degrees two and three. Sampling with replacement allows several clusters
to receive the same link function.

For an example from cluster \(c\), the response is
\begin{equation}
    y
    =
    g_c\!\left(\langle x,v_c\rangle\right)+\varepsilon,
    \qquad
    \varepsilon\sim\mathcal N(0,0.02^2).
    \label{eq:local-routing-mixed-target}
\end{equation}
Because \(v_c\) is supported only on the predictive coordinates,
\[
    \langle x,v_c\rangle
    =
    \widetilde v_c^\top z_{\mathrm{pred}}.
\]
Thus, both the predictive direction and the nonlinear link can change across
clusters.

The realized link assignments for the three displayed seeds are reported in
Table~\ref{tab:link-function-counts}. For a fixed \((K,\text{seed})\), these
assignments are shared across all methods and training sizes.

\begin{table}[t]
\centering
\small
\caption{Number of clusters assigned to each link function for the three
runs shown in Fig.~\ref{fig:mixed_link_functions_all_activations}.}
\label{tab:link-function-counts}
\begin{tabular}{@{}rrrrrr@{}}
\toprule
\(K\) & Seed & Hermite-2 & Hermite-3 & Sine & Tanh\\
\midrule
1  & 1000 & 0  & 0  & 0  & 1\\
1  & 1001 & 0  & 1  & 0  & 0\\
1  & 1002 & 0  & 0  & 0  & 1\\
\midrule
2  & 1000 & 0  & 0  & 0  & 2\\
2  & 1001 & 0  & 2  & 0  & 0\\
2  & 1002 & 0  & 1  & 0  & 1\\
\midrule
10 & 1000 & 3  & 2  & 1  & 4\\
10 & 1001 & 2  & 3  & 3  & 2\\
10 & 1002 & 5  & 1  & 2  & 2\\
\midrule
50 & 1000 & 17 & 13 & 12 & 8\\
50 & 1001 & 14 & 17 & 12 & 7\\
50 & 1002 & 11 & 14 & 14 & 11\\
\bottomrule
\end{tabular}
\end{table}


\paragraph{Training, validation, and test data.}
For each \((K,\text{seed})\), the training sets of sizes
\(20{,}000\), \(50{,}000\), \(100{,}000\), and \(200{,}000\) are
cluster-balanced nested subsets of one \(200{,}000\)-example pool.
The \(500{,}000\)-example set is generated separately with the same
geometry and link assignments and is not a nested extension of this pool.
All methods share the same training data and independent \(4096\)-example
test set for each condition. The runs listed in
Table~\ref{tab:validation-protocol} additionally use an independent
\(20{,}000\)-example validation set for checkpoint selection.
Test data are used only for final evaluation.

\paragraph{Random seeds and averaging.}
Results are means over three runs, with seeds \(1000,1001,1002\), without
standard-deviation bands. Cluster centers and predictive directions are
held fixed using geometry seed \(271828\); link assignments, sampled data,
label noise, network initialization, and optimization minibatches vary
across runs. For run \(r\in\{0,1,2\}\), the training, test, and validation
seeds are \(12345+r\), \(54321+r\), and \(84321+r\), respectively.

\paragraph{MLP architectures.}
We evaluate one-hidden-layer ReLU, GELU, ReGLU, and SwiGLU networks of width
\(4{,}096\). ReLU and GELU use
\[
    h(x)=\phi(Wx+b),
\]
while the gated models use
\[
    h(x)
    =
    \phi(W_gx+b_g)\odot(W_vx+b_v).
\]
For ReGLU, \(\phi=\operatorname{ReLU}\); for SwiGLU,
\(\phi=\operatorname{SiLU}\). All hidden weights, readout weights, and biases
are trainable.

Hidden-weight entries are initialized independently with standard deviation
\(10^{-3}/\sqrt d\), and all biases are initialized to zero. For ordinary
MLPs, the initial readout magnitude of neuron \(j\) is proportional to
\(\|w_j\|_2\), with balanced random signs. For gated MLPs, it is proportional
to
\[
    \sqrt{\|w_{g,j}\|_2\|w_{v,j}\|_2}.
\]
The output multiplier is \(1024/4096\).


\paragraph{MLP optimization and evaluation.}
All MLPs use Adam, batch size \(4096\), and cosine learning-rate decay
with horizon \(200{,}000\) updates and final rate \(\eta_0/100\).
Table~\ref{tab:mlp-hparams} gives the learning rates and weight decay
used in each condition. Except for the runs listed in
Table~\ref{tab:validation-protocol}, models are trained on all \(n\)
examples for the following fixed durations, in thousands of updates,
and evaluated at the final checkpoint:
\[
\begin{array}{c|cccc}
K & \text{ReLU} & \text{GELU} & \text{ReGLU} & \text{SwiGLU}\\
\hline
1  & 18    & 35    & 39   & 189.5\\
2  & 104   & 198.5 & 71.5 & 194\\
10 & 129.5 & 198.5 & 135  & 198\\
50 & 17.5  & 113.5 & 54   & 196
\end{array}
\]

\begin{table*}[t]
\centering
\small
\renewcommand{\arraystretch}{1.12}
\setlength{\tabcolsep}{4pt}
\caption{Single-index runs using validation-based checkpoint selection.
Each run uses at most \(200{,}000\) updates and restores the checkpoint
with the lowest MSE on an independent \(20{,}000\)-example validation set.
All other runs use the fixed durations reported in this subsection.}
\label{tab:validation-protocol}
\begin{tabularx}{\textwidth}{@{}lXr@{}}
\toprule
Model & Conditions & Validation interval \\
\midrule
ReLU & \(K=10\), \(n\in\{200{,}000,500{,}000\}\) & \(500\) updates \\
SwiGLU & \(K\in\{10,50\}\), \(n\geq50{,}000\) & \(1000\) updates \\
ReLU, GELU & \(K=50\), \(n\geq100{,}000\) & \(1000\) updates \\
\bottomrule
\end{tabularx}
\end{table*}

For the \(K=10\) ReLU runs in Table~\ref{tab:validation-protocol},
early stopping uses a training-MSE target of \(4\times10^{-4}\) or a
relative-improvement threshold of \(5\times10^{-4}\) over 60 evaluations
after the minimum training duration.
The other listed runs use validation-based early stopping: after training
MSE reaches \(10^{-4}\), the relative-improvement threshold is
\(5\times10^{-3}\) over 15 evaluations; after \(175{,}000\) updates,
a second criterion uses \(10^{-3}\) over 20 evaluations.

\paragraph{MLP hyperparameters used in the displayed figure.}
Table~\ref{tab:mlp-hparams} reports the exact initial learning rate and weight
decay used for every displayed MLP curve. Each entry is
\((\eta_0,\lambda)\).

\begin{table*}[t]
\centering
\small
\caption{MLP hyperparameters used in
Fig.~\ref{fig:mixed_link_functions_all_activations}. Each entry is
\((\eta_0,\lambda)\), denoting initial learning rate and weight decay.}
\label{tab:mlp-hparams}
\resizebox{\textwidth}{!}{%
\begin{tabular}{@{}lcccc@{}}
\toprule
Condition & ReLU & GELU & ReGLU & SwiGLU\\
\midrule
\(K=1\), all \(n\)
& \((10^{-2},10^{-5})\)
& \((10^{-2},10^{-6})\)
& \((10^{-3},10^{-5})\)
& \((3{\times}10^{-2},10^{-6})\)
\\

\(K=2\), all \(n\)
& \((3.33{\times}10^{-3},10^{-5})\)
& \((10^{-2},10^{-6})\)
& \((10^{-3},10^{-5})\)
& \((3.33{\times}10^{-3},10^{-6})\)
\\

\(K=10,\ n=20\mathrm{k}\)
& \((3{\times}10^{-2},10^{-5})\)
& \((10^{-2},10^{-6})\)
& \((10^{-3},10^{-5})\)
& \((10^{-2},10^{-6})\)
\\

\(K=10,\ n\in\{50\mathrm{k},100\mathrm{k}\}\)
& \((3{\times}10^{-2},10^{-5})\)
& \((10^{-2},10^{-6})\)
& \((10^{-3},10^{-5})\)
& \((10^{-2},0)\)
\\

\(K=10,\ n\in\{200\mathrm{k},500\mathrm{k}\}\)
& \((10^{-2},10^{-6})\)
& \((10^{-2},10^{-6})\)
& \((10^{-3},10^{-5})\)
& \((10^{-2},0)\)
\\

\(K=50,\ n=20\mathrm{k}\)
& \((3{\times}10^{-2},0)\)
& \((3{\times}10^{-2},10^{-6})\)
& \((9{\times}10^{-3},10^{-6})\)
& \((3{\times}10^{-2},10^{-6})\)
\\

\(K=50,\ n=50\mathrm{k}\)
& \((3{\times}10^{-2},0)\)
& \((3{\times}10^{-2},10^{-6})\)
& \((9{\times}10^{-3},10^{-6})\)
& \((10^{-2},0)\)
\\

\(K=50,\ n=100\mathrm{k}\)
& \((3{\times}10^{-2},10^{-6})\)
& \((3{\times}10^{-2},10^{-6})\)
& \((9{\times}10^{-3},10^{-6})\)
& \((10^{-2},0)\)
\\

\(K=50,\ n\in\{200\mathrm{k},500\mathrm{k}\}\)
& \((10^{-2},0)\)
& \((3{\times}10^{-2},10^{-6})\)
& \((9{\times}10^{-3},10^{-6})\)
& \((10^{-2},0)\)
\\
\bottomrule
\end{tabular}%
}
\end{table*}


\paragraph{Global Laplace and RFM.}
RFM uses the metric-dependent Laplace kernel
\begin{equation}
    k_t(x,x')
    =
    \exp\!\left(
        -\frac{\|M_t^{1/2}(x-x')\|_2}{h_t}
    \right),
    \label{eq:local-routing-rfm-kernel}
\end{equation}
where \(M_t\) is the metric at iteration \(t\) and \(h_t\) is the bandwidth.
We initialize \(M_0=I_{20}\), so iteration zero is the isotropic Laplace
kernel.

After fitting the kernel predictor, RFM estimates
\[
    \widehat G_t
    =
    \frac{1}{m}
    \sum_{i=1}^{m}
    \nabla\widehat f_t(x_i)
    \nabla\widehat f_t(x_i)^\top
\]
using at most \(m=20{,}000\) training examples. A normalized version of
\(\widehat G_t\) defines the next metric. We perform at most three metric
updates, producing iterations \(t\in\{0,1,2,3\}\).

The base bandwidth is computed from within-cluster pairwise distances in the
current transformed space: we compute one median distance per cluster and
take the median across clusters. This calculation uses cluster identities
only to calibrate one scalar bandwidth. Global Laplace and global RFM still
fit a single predictor to the complete training set and do not receive the
cluster identity at prediction time.

Table~\ref{tab:kernel-hparams} gives the configurations used across
sample sizes and seeds, including \(n=500{,}000\). Each entry is
\((b,t,\lambda)\), where \(b\) is the bandwidth multiplier, \(t\) is the RFM
iteration, and \(\lambda\) is the ridge parameter.

\begin{table*}[t]
\centering
\small
\caption{Kernel configurations used in
Fig.~\ref{fig:mixed_link_functions_all_activations}.}
\label{tab:kernel-hparams}
\resizebox{\textwidth}{!}{%
\begin{tabular}{@{}lcccc@{}}
\toprule
Condition
& Global Laplace
& Global RFM
& Local Laplace
& Local RFM\\
\midrule
\(K=1\), all \(n\)
& \((2,0,10^{-6})\)
& \((2,3,10^{-6})\)
& same as global
& same as global
\\

\(K=2\), all \(n\)
& \((2,0,10^{-6})\)
& \((1,1,10^{-6})\)
& \((2,0,10^{-8})\)
& \((2,3,10^{-8})\)
\\

\(K=10\), all \(n\)
& \((2,0,10^{-6})\)
& \((2,1,10^{-6})\)
& \((2,0,10^{-8})\)
& \((2,3,10^{-6})\)
\\

\(K=50\), all \(n\)
& \((2,0,10^{-6})\)
& \((2,2,10^{-6})\)
& \((2,0,10^{-8})\)
& \((2,3,10^{-8})\)
\\
\bottomrule
\end{tabular}%
}
\end{table*}

Exact kernel solves are used when an individual kernel problem contains fewer
than \(11{,}000\) training examples. Larger kernel problems are solved using
EigenPro~\citep{ma2018power,ma2019kernel,abedsoltan2023toward,
abedsoltan2025fast}. EigenPro uses a training-MSE target of
\(4\times10^{-4}\), equal to the observation-noise variance. The explicit
ridge parameter applies directly to exact solves; EigenPro is regularized
primarily through early stopping.


\paragraph{Local methods.}
The local Laplace and RFM methods receive the true cluster identity
during both training and testing. They partition the training set by cluster,
fit an independent predictor within each partition, and evaluate each test
example using the predictor associated with its true cluster.

Each local RFM  performs at most three metric updates. For \(K=1\), the global
and local methods coincide.

These local methods are diagnostic references rather than fair predictive
baselines because the cluster identity is supplied directly. Their purpose is to
measure the difficulty of learning the local predictive functions once the
cluster assignment is known.


\paragraph{Bayes reference.}
The dotted line is the empirical test MSE of the Bayes predictor under the
known data-generating distribution. It computes the posterior probability of
each cluster given \(x\) and averages the corresponding noiseless local
predictions. The irreducible observation-noise variance is
\[
    0.02^2=4\times10^{-4}.
\]
Across the test sets used in the updated figure, the empirical Bayes MSE lies
between approximately \(3.97\times10^{-4}\) and
\(4.14\times10^{-4}\), consistent with finite-test-set variation around the
noise variance.

\FloatBarrier
\subsubsection{A Common Third-Hermite Link Function}
\label{app:fixed-h3-single-index-details}

\paragraph{Fixed third-Hermite single-index model.}
We use the same clustered single-index geometry with \(d=20\) and
\[
    K\in\{1,2,10,50\}.
\]
Each cluster has an independently sampled predictive direction, but every
cluster uses the common normalized third-Hermite link
\[
    g_c(t)
    =
    h_3(t)
    =
    \frac{t^3-3t}{\sqrt{6}},
    \qquad c\in[K].
\]
A common link function isolates the role of cluster-dependent predictive
directions and supports the third-Hermite analysis in
Section~\ref{sec:early-specialization}.

\begin{figure}[H]
    \centering
    \includegraphics[width=\textwidth]
    {figs/singleindex_fixed_h3_3seed_original_all_activations_plot_only.pdf}
    \includegraphics[width=\textwidth]
    {figs/singleindex_final_n500k_3seed_original_all_activations_legend_only.pdf}
    \caption{
    \textbf{Fixed third-Hermite single-index model.}
    This uses the same clustered single-index data construction as Figure Figure~\ref{fig:local-sample-complexity}, but with a common normalized third-Hermite link, providing a setting closer to the theory in Section~\ref{sec:specialization-theory}.
    }
    \label{fig:fixed-h3-single-index-results}
\end{figure}

\paragraph{Results.}
Figure~\ref{fig:fixed-h3-single-index-results} shows that the separation
between the global methods and the MLPs grows with the number of clusters.
Global Laplace and RFM become progressively less effective as the
cluster-specific directions occupy more of the predictive subspace. The MLPs
continue to exploit the cluster-dependent directions, and the gated
architectures perform especially well in the many-cluster regime.


\subsection{Additional Details for Trained MLPs Encoding Cluster Structure}
\label{app:learned-gating-details}

This appendix supplements the experiment in
Section~\ref{sec:learned-gating}, which builds on the clustered single-index
setting of Section~\ref{sec:local-routing}. We describe how cluster gates are
extracted from a trained ReLU MLP and reused to construct independent local
Laplace and RFM predictors.

\paragraph{Clustered data with mixed link functions.}
We use the \(d=20\) input distribution in
Eq.~\ref{eq:local-routing-data}, with ten cluster-identifying and ten
predictive coordinates, cluster-center radius \(R=1\), and
\(K\in\{2,5,10,50\}\) equally represented clusters. Cluster-center and
predictive directions follow Section~\ref{appendix:main_exp_detail},
with the cluster-noise rule in Eq.~\ref{eq:local-routing-noise} and
\(\gamma=4\). Each cluster's link function is sampled independently and
uniformly from normalized Hermite-3, sine, and hyperbolic tangent.
Responses follow Eq.~\ref{eq:local-routing-mixed-target}, with
observation-noise standard deviation \(0.02\).

\paragraph{Datasets and repetitions.}
For every \(K\), we generate a balanced training pool containing
\(200{,}000\) examples and use nested prefixes of sizes
\[
    n\in\{20{,}000,50{,}000,100{,}000,200{,}000\}.
\]
Evaluation uses an independently generated test set containing \(10{,}000\)
examples. Within each run, exactly the same serialized training and test
tensors are used by the MLP, global kernel methods, ground-truth local
methods, and MLP-gated methods.

Results are averaged over three seeds. For a fixed \(K\), the cluster-position
and predictive directions are held fixed, while each seed independently
resamples the cluster-specific link-function assignments, training and test
examples,
label noise, and MLP initialization. No test labels are used to construct the
gates.

\paragraph{Source ReLU MLP.}
For each \(K\), sample size, and seed, we train a one-hidden-layer ReLU network
of width \(4{,}096\),
\[
    \widehat f(x)
    =
    \sum_{j=1}^{4096}
    a_j\operatorname{ReLU}(w_j^\top x+b_j)
    +b_{\mathrm{out}}.
\]
All weights and biases are trainable. Optimization uses Adam with minibatches
of size \(4{,}096\), zero weight decay, and cosine learning-rate decay from
\(10^{-2}\) to \(10^{-4}\). Training runs for at least \(20{,}000\) steps and
at most \(200{,}000\) steps. It terminates when the full training MSE reaches
\(4\times10^{-4}\), or when the training loss plateaus. We restore the
checkpoint with the smallest full training MSE. Neither test MSE nor cluster
identity is used for checkpoint selection.

\paragraph{Selecting active neurons.}
After training, we freeze the ReLU MLP and select the active-neuron set
\(\mathcal J\) using the shared importance-based criterion in
Section~\ref{app:shared-specialization-analysis}.
Here, these neurons are used to construct contribution profiles for
clustering inputs and gate-specific features for the local predictors,
rather than to measure alignment with the predictive directions.

\paragraph{Contribution profiles and learned gates.}
For every training input \(x\), we compute its normalized contribution profile
over the active neurons:
\begin{equation}
    p_j(x)
    =
    \frac{
        |a_j|\operatorname{ReLU}(w_j^\top x+b_j)
    }{
        \displaystyle
        \sum_{\ell\in\mathcal J}
        |a_\ell|\operatorname{ReLU}(w_\ell^\top x+b_\ell)
    },
    \qquad j\in\mathcal J.
    \label{eq:appendix-contribution-profile}
\end{equation}
Thus \(p(x)\) describes which active first-layer neurons contribute to the
MLP prediction on input \(x\), independent of the overall magnitude of the
prediction.

We apply \(K\)-means++ to the training contribution profiles and use
\(M=K\) learned gates. Knowledge of the number of clusters is therefore
provided, but the true cluster assignments are never observed. We use eight
initializations and at most 30 Lloyd iterations, retaining the solution with
the smallest training inertia. Every training example is assigned to exactly
one gate through its nearest centroid.

Let \(q_c\) denote the centroid of learned gate \(c\), with coordinate
\(q_{cj}\) corresponding to active neuron \(j\). Although examples receive
hard gate assignments, the centroid coordinates provide a soft association
between neurons and gates: the same neuron may contribute to several gates.

\paragraph{Constructing gate-specific features.}
For each learned gate \(c\), we form the positive-semidefinite metric
\begin{equation}
    G_c
    =
    \frac{
        \displaystyle
        \sum_{j\in\mathcal J}q_{cj}w_jw_j^\top
    }{
        \displaystyle
        \sum_{j\in\mathcal J}q_{cj}
    }.
    \label{eq:appendix-mlp-derived-metric}
\end{equation}
The corresponding gate-specific representation is
\begin{equation}
    z_c(x)=G_c^{1/2}x.
    \label{eq:appendix-mlp-derived-projection}
\end{equation}
The MLP biases affect the contribution profiles and hence the gate
assignments, while \(G_c\) itself is constructed from the incoming
first-layer weight directions.

For every learned gate, we project only the training examples assigned to
that gate and fit an independent local Laplace or RFM predictor in the
resulting representation.

\paragraph{Test-time prediction.}
For a new test input \(x\), we:

\begin{enumerate}
    \item pass \(x\) through the frozen MLP and compute its contribution
    profile \(p(x)\);
    \item assign \(x\) to the nearest fixed training centroid;
    \item transform it using the corresponding representation
    \(z_c(x)=G_c^{1/2}x\); and
    \item evaluate the local predictor fitted for that learned gate.
\end{enumerate}

The test input is therefore assigned using only its MLP contribution profile.
Neither its response nor its true cluster identity is used.

\paragraph{Local Laplace and RFM predictors.}
The local Laplace predictor for gate \(c\) uses
\[
    k_c(x,x')
    =
    \exp\!\left(
        -\frac{
            \|G_c^{1/2}(x-x')\|_2
        }{h_c}
    \right),
\]
where \(h_c\) is initialized using the median pairwise distance within that
gate. We use ridge parameter \(10^{-6}\).

The local RFM begins from this Laplace kernel and performs three metric
updates, producing iterations \(0,1,2,3\), where iteration \(0\) is the local
Laplace predictor. After each update, the bandwidth is recomputed using
distances under the updated metric. Kernel systems with fewer than
\(11{,}000\) local training examples are solved directly; larger systems use
EigenPro. The same procedure is used for the global and ground-truth local
RFM baselines.

\paragraph{Compared methods.}
Figure~\ref{fig:local-sample-complexity_kmeans} compares:

\begin{itemize}
    \item the source ReLU MLP;
    \item one global Laplace predictor fitted to all training examples;
    \item one global RFM fitted to all training examples;
    \item MLP-gated local Laplace, using the gates and projections extracted
    from the ReLU MLP;
    \item MLP-gated local RFM, using the same extracted gates and projections;
    \item local Laplace, which is given the true cluster identities at
    training and test time;
    \item local RFM, defined analogously; and
    \item the irreducible noise level, whose expected MSE is
    \(0.02^2=4\times10^{-4}\).
\end{itemize}

The local methods using the true clusters are diagnostic references rather
than fair deployable predictors because they receive the true cluster
identity. By contrast, the
MLP-gated methods receive only the training input--response pairs and the
known number \(K\) of gates.





\paragraph{Aggregation.}
Every curve in Figure~\ref{fig:local-sample-complexity_kmeans} is the
arithmetic mean over three independent runs. Each run uses a newly sampled
dataset from the same generative model and an independently initialized MLP.


\newpage

\subsection{Additional Details for Clustered Multi-Index Models}
\label{app:clustered-multi-index-details}

This appendix supplements Section~\ref{sec:clustered-multi-index}.
Section~\ref{app:multiindex-data} specifies the data distribution.
Section~\ref{app:multiindex-training} reports the models and training settings, and
Section~\ref{app:multiindex-specialization-null} extends the specialization
analysis to three-dimensional predictive spans.

\subsubsection{Data Generation}
\label{app:multiindex-data}

\paragraph{Cluster geometry.}
We follow the Gaussian-mixture construction in
Section~\ref{appendix:main_exp_detail}, increasing the input dimension to
\(d=50\), with coordinates \(1{:}20\) identifying the cluster,
\(21{:}40\) containing the predictive variables, and \(41{:}50\)
containing nuisance variables. We consider
\[
    K\in\{1,2,10,25\},
    \qquad
    n\in\{20{,}000,50{,}000,100{,}000,200{,}000,500{,}000\}.
\]
Clusters have equal sampling probabilities, and training sets contain
exactly \(n/K\) examples from each cluster.

For \(K>1\), cluster centers are
\[
    \mu_c=(s_c,0_{20},0_{10}),
    \qquad
    s_c\in\mathbb R^{20},
    \qquad
    \|s_c\|_2=1.
\]
The vectors \(s_c\) are selected from \(20{,}000\) random unit-vector
candidates in \(\mathbb R^{20}\) using the farthest-point procedure of
Section~\ref{appendix:main_exp_detail}. We use the same cluster-noise rule
in Eq.~\ref{eq:local-routing-noise}, with \(\gamma=4\), and the same
\(K=1\) convention. With the enlarged blocks and additional nuisance
variables, the conditional input distribution becomes
\[
    x\mid c
    \sim
    \mathcal N\!\left(
        \mu_c,\,
        \operatorname{diag}(\sigma_K^2I_{20},I_{20},I_{10})
    \right).
\]

\paragraph{Predictive spans and responses.}
Independently for each cluster, we generate a \(20\times3\) matrix with
independent standard Gaussian entries and compute its thin QR
decomposition. If \(q_{c,1},q_{c,2},q_{c,3}\) are the resulting orthonormal
columns, we set
\[
    v_{c,r}=(0_{20},q_{c,r},0_{10}),
    \qquad r=1,2,3.
\]
Directions are orthonormal within a cluster; frames from different
clusters are neither shared nor constrained to be orthogonal.

Each cluster independently receives one link function, sampled uniformly
with replacement from the family in Eq.~\ref{eq:local-link-function-family}
after excluding the third-order Hermite link. The remaining choices are
the normalized second-order Hermite polynomial, sine, and hyperbolic tangent.
Responses follow the three-index model in Section~\ref{sec:prelim}, with
the same observation-noise standard deviation \(0.02\) as in
Section~\ref{appendix:main_exp_detail}. Each predictive projection has
unit variance.

\paragraph{Shared data and repeated runs.}
For each \(K\), cluster centers and predictive frames are generated using
geometry seed \(271828+1000003K\) and held fixed across runs.
The three run seeds \(1000,1001,1002\) vary the link-function assignments,
sampled inputs, observation noise, network initialization, and
optimization minibatches.

All training sets are nested, cluster-balanced subsets of a
single \(500{,}000\)-example pool for each \((K,\text{seed})\).
Every run uses independently sampled validation and test sets of sizes
\(20{,}000\) and \(4{,}096\), respectively. All methods share the same
saved training, validation, and test data for each \((K,n,\text{seed})\).
Validation data are used for model selection, and test data only for final
evaluation. We report the mean test MSE over three runs without
standard-deviation bands.

\subsubsection{Models and Training}
\label{app:multiindex-training}

\paragraph{MLP architecture and initialization.}
We use the ReLU, GELU, ReGLU, and SwiGLU architectures described in
Section~\ref{appendix:main_exp_detail}, with width \(2048\) and input
dimension \(d=50\). Hidden and output biases are trainable.
Incoming weights are initialized with standard deviation
\(10^{-3}/\sqrt{50}\),
and biases are initialized to zero. Readout magnitudes follow the
norm-based formulas in Section~\ref{appendix:main_exp_detail}, with
proportionality constant one and independent random signs. The output
multiplier is \(1024/2048\).

\paragraph{Optimization and checkpoint selection.}
We use Adam, batch size \(4096\), and cosine decay from \(\eta_0\) to
\(\eta_0/100\), with a schedule horizon and maximum training duration of
\(100{,}000\) updates. Validation MSE is evaluated every \(1000\) updates, and the
checkpoint with the lowest validation MSE is retained.

Training may stop after \(30{,}000\) updates if validation MSE has not
improved by at least \(0.1\%\) relative to its last meaningful improvement
for twenty consecutive evaluations.

The following learning rates and Adam weight-decay coefficients are
used for all \(K\), sample sizes, and seeds:
\[
\begin{array}{lcc}
\text{Activation} & \eta_0 & \text{Weight decay}\\
\hline
\text{ReLU}   & 2\times10^{-3} & 10^{-5}\\
\text{GELU}   & 10^{-2}        & 0\\
\text{SwiGLU} & 5\times10^{-3} & 0
\end{array}
\]
ReGLU uses the following fixed condition-specific settings, shared
across the three seeds. Each entry is
\((\eta_0,\text{weight decay})\):
\begin{center}
\resizebox{\linewidth}{!}{\(
\begin{array}{c|cccc}
n & K=1 & K=2 & K=10 & K=25\\
\hline
20{,}000
& (10^{-3},10^{-5})
& (10^{-3},10^{-5})
& (9\times10^{-3},10^{-5})
& (9\times10^{-3},0)\\
50{,}000
& (10^{-3},10^{-5})
& (3\times10^{-3},10^{-5})
& (9\times10^{-3},10^{-6})
& (9\times10^{-3},10^{-6})\\
100{,}000
& (9\times10^{-3},10^{-6})
& (9\times10^{-3},10^{-6})
& (9\times10^{-3},10^{-6})
& (9\times10^{-3},10^{-6})\\
200{,}000,\ 500{,}000
& (10^{-3},10^{-6})
& (3\times10^{-3},10^{-6})
& (3\times10^{-3},10^{-6})
& (9\times10^{-3},10^{-6})
\end{array}
\)}
\end{center}
\paragraph{Kernel methods.}
We use the global and local Laplace and RFM predictors defined in
Section~\ref{appendix:main_exp_detail}, using the full fifty-dimensional
input, identity initialization, three RFM metric updates, and
iteration zero as the Laplace baseline. At each iteration, the bandwidth is the
median pairwise distance
among up to \(512\) uniformly sampled training inputs after applying
the current metric transform. Bandwidths are recomputed after metric
updates; neither bandwidth nor ridge regularization is tuned on validation.

Metric updates use the average gradient outer product from
Section~\ref{appendix:main_exp_detail}, estimated from up to \(20{,}000\)
training examples from the relevant global or local training set.
Here, the estimate is symmetrized and divided
by its largest matrix entry before constructing the next transform.

Kernel systems with fewer than \(11{,}000\) training examples are solved
directly; larger systems use EigenPro. Direct solves use regularization
\(K_{\mathrm{train}}+n_{\mathrm{local}}10^{-6}I\).
EigenPro uses a training-MSE target of \(4\times10^{-4}\) and a maximum
of \(5000\) epochs.
All completed EigenPro fits reached this training-error target.

The RFM iteration is selected by minimum validation MSE, including
iteration zero as a candidate. For the local method, we select one common
iteration using validation MSE pooled across clusters, weighted by their
validation sample counts. The selected predictor and the iteration-zero
Laplace predictor are then evaluated on the held-out test set.

\subsubsection{Specialization Analysis and Random-Subspace Control}
\label{app:multiindex-specialization-null}

\paragraph{Active neurons.}
We analyze the validation-selected models at \(K=10\) and
\(n=500{,}000\), using the active-neuron criterion in
Section~\ref{app:shared-specialization-analysis}, as in
Section~\ref{app:singleindex-specialization-null}. This includes the same
\(99.9\%\) importance threshold and separate gate/value analyses on the
same active set.

\paragraph{Span-specialization measure.}
We extend the direction-alignment measure in
Section~\ref{app:singleindex-specialization-null} to alignment with an
entire cluster-specific predictive span. The orthogonal projector onto
cluster \(c\)'s predictive span is
\[
    P_c=\sum_{r=1}^{3}v_{c,r}v_{c,r}^{\top}.
\]
For an incoming weight \(w_j\), define
\begin{equation}
    S_j=\max_{1\leq c\leq K}
    \frac{\|P_cw_j\|_2}{\|w_j\|_2}.
    \label{eq:multiindex-span-specialization}
12\end{equation}
The denominator includes all fifty coordinates, including
cluster-identifying and nuisance coordinates.
We report the percentage of active neurons satisfying
\(S_j\geq0.71\) or \(S_j\geq0.90\), using these literal thresholds.
The table reports mean and sample standard deviation over three
trained models. This measure detects alignment with a cluster's
predictive span, not necessarily with an individual direction within it.

\paragraph{Random-subspace control.}
We use the shared protocol in Section~\ref{app:specialization-null},
replacing the random directions of
Section~\ref{app:singleindex-specialization-null} with random
three-dimensional subspaces. For each of the \(B=5{,}000\) draws per
trained model, we independently generate ten such subspaces within the
twenty-dimensional predictive subspace, using the Gaussian QR construction
in Section~\ref{app:multiindex-data}. We then substitute their projectors
for the true projectors in Eq.~\ref{eq:multiindex-span-specialization} and
recompute the percentages at the same thresholds.

We aggregate the results as specified in
Section~\ref{app:specialization-null}. The same random subspace collections are used across activations
and weight types within each seed.

Table~\ref{tab:multiindex-mixed-k10-span-specialization} reports the
observed and random-baseline results. For every row and threshold, the
maximum per-seed Monte Carlo p-value is \(p_{\max}=1/5001\), using the shared definition in
Section~\ref{app:specialization-null}.

\newpage
\section{Proofs for Section~\ref{sec:specialization-theory}}
\label{app:specialization-theory-proofs}

\subsection{Proof of Theorem~\ref{thm:neurons-specialize}}
\label{app:proof-neurons-specialize}

We first prove Theorem~\ref{thm:neurons-specialize} invoking
\cref{lem:cubic-correlation-identity,lem:one-cluster-maxima,lem:specialized-maxima,lem:positive-signed-flow,lem:self-selected-specialization,lem:uniform-cluster-coverage,lem:early-time-tracking},
whose statements and proofs are provided after the main proof.
A roadmap of their dependencies is shown in
\cref{fig:theorem1-dependencies}.

{
Throughout this subsection, the normalized means and predictive directions
\[
    \frac{\mu_1}{R},\ldots,\frac{\mu_K}{R},
    v_1,\ldots,v_K
\]
form an orthonormal basis of $\mathbb R^d$, where $d=2K$. In the notation of
Section~\ref{sec:prelim}, we may write
\[
    c\sim\operatorname{Unif}([K]),
    \qquad
    x=\mu_c+z,
    \qquad
    z\sim\mathcal N(0,I_d),
\]
and
\[
    y=h_3\bigl(v_c^\top z\bigr),
    \qquad
    h_3(t)=\frac{t^3-3t}{\sqrt6}.
\]

For $\omega\in\mathbb S^{d-1}$, define the weight's correlations with the cluster means and predictive directions
\[
    b_c(\omega)
    =
    \langle\omega,\mu_c\rangle,
    \qquad
    \rho_c(\omega)
    =
    \langle\omega,v_c\rangle.
\]
}
Write the ReLU activation function and the Gaussian density as
\[
    \phi(t)=t_+,
    \qquad
    \varphi(t)
    =
    \frac{1}{\sqrt{2\pi}}e^{-t^2/2},
\]
and define the unsigned and signed population-correlation objectives, respectively, by
\[
    \Phi({\omega})
    =
    \mathbb E\left[{y}\,
      \phi({\omega}^\top{x})\right],
    \qquad
    \Psi_\zeta({\omega})
    =
    \zeta\Phi({\omega}),
    \qquad
    \zeta\in\{\pm1\}.
\]

\begin{proof}[Proof of Theorem~\ref{thm:neurons-specialize}]
For each initial direction ${\omega_j^0}$, let
$(\bar u_j,{\bar\omega_j})$ be the corresponding self-selected teacher-only
trajectory from Lemma~\ref{lem:self-selected-specialization}.

By Lemmas~\ref{lem:self-selected-specialization}
and~\ref{lem:specialized-maxima}, almost surely there are a cluster label
$J_j\in[K]$ and an orientation $\tau_j\in\{\pm1\}$ such that
\[
    \limsup_{t\to\infty}
    \left\|
        {\bar\omega_j(t)-\tau_j v_{J_j}}
    \right\|_2
    \leq
    \frac{C}{R}
\]
for a universal constant $C$.

Since $m<\infty$ and every teacher-only direction converges, for every
$\delta>0$ there is a finite time $T_\delta$ such that, simultaneously for
all $j\in[m]$,
\[
    \left\|
        {\bar\omega_j(T_\delta)-\tau_j v_{J_j}}
    \right\|_2
    \leq
    \frac{\delta}{2}+\frac{C}{R}.
\]

Apply Lemma~\ref{lem:early-time-tracking} on $[0,T_\delta]$. For sufficiently
small $\varepsilon$,
\[
    \left\|
        \frac{w_j(T_\delta)}{\|w_j(T_\delta)\|_2}
        -
        {\bar\omega_j(T_\delta)}
    \right\|_2
    \leq
    L_{T_\delta}m\varepsilon^2
    \leq
    \frac{\delta}{2}
\]
for every $j\in[m]$. The triangle inequality therefore gives
\[
    \left\|
        \frac{w_j(T_\delta)}{\|w_j(T_\delta)\|_2}
        -
        \tau_j {v_{J_j}}
    \right\|_2
    \leq
    \delta+\frac{C}{R}.
\]

The independence and uniformity of the labels $J_1,\ldots,J_m$, together
with the cluster-coverage probability, follow from
Lemma~\ref{lem:uniform-cluster-coverage}.
\end{proof}

{
\begin{figure}[ht!]
\centering
\resizebox{\linewidth}{!}{%
\begin{tikzpicture}[
    node distance=7mm and 12mm,
    >=Latex,
    every node/.style={
        draw,
        rounded corners,
        align=center,
        font=\small,
        inner sep=5pt,
        text width=3.4cm
    },
    theorem/.style={
        draw,
        very thick,
        text width=4.0cm
    },
    arrow/.style={
        ->,
        thick
    }
]

\node (corr)
{
Lemma~\ref{lem:cubic-correlation-identity}\\
\textbf{Cubic correlation identity}\\
$\Psi_\zeta(\omega)
 = \frac1K\sum_c
 \zeta\eta_3(b_c)\rho_c^3$
};

\node (onecluster) [below left=of corr]
{
Lemma~\ref{lem:one-cluster-maxima}\\
\textbf{One-cluster landscape}\\
positive local maxima cannot
mix predictive directions
};

\node (specialmax) [below right=of corr]
{
Lemma~\ref{lem:specialized-maxima}\\
\textbf{Specialized maxima}\\
$\omega^\star
 = \tau v_c
 + O(R^{-1})$
};

\node (flow) [below=14mm of $(onecluster)!0.5!(specialmax)$]
{
Lemma~\ref{lem:positive-signed-flow}\\
\textbf{Gradient-flow convergence}\\
unstable mixed critical points
are avoided a.s.
};

\node (selfsign) [below=of flow]
{
Lemma~\ref{lem:self-selected-specialization}\\
\textbf{Self-selected sign}\\
teacher-only dynamics select a sign
and specialize a.s.
};

\node (coverage) [below left=of selfsign]
{
Lemma~\ref{lem:uniform-cluster-coverage}\\
\textbf{Uniform labels and coverage}\\
$J_j \stackrel{\mathrm{iid}}{\sim}
\operatorname{Unif}([K])$
};

\node (tracking) [below right=of selfsign]
{
Lemma~\ref{lem:early-time-tracking}\\
\textbf{Full-network tracking}\\
coupled dynamics differ from
teacher-only dynamics by
$O(m\varepsilon^2)$
};

\node[theorem] (thm) [below=15mm of $(coverage)!0.5!(tracking)$]
{
Theorem~\ref{thm:neurons-specialize}\\
\textbf{Randomly initialized neurons specialize}\\[2pt]
$\displaystyle
\frac{w_j(T_\delta)}
     {\|w_j(T_\delta)\|_2}
\approx
\tau_j v_{J_j}$
};

\draw[arrow] (corr) -- (onecluster);
\draw[arrow] (corr) -- (specialmax);
\draw[arrow] (onecluster) -- (flow);
\draw[arrow] (specialmax) -- (flow);
\draw[arrow] (flow) -- (selfsign);
\draw[arrow] (selfsign) -- (coverage);
\draw[arrow] (selfsign) -- (tracking);
\draw[arrow, dashed]
    (corr.east)
    .. controls +(3.0,0) and +(2.5,1.0) ..
    (tracking.east);
\draw[arrow] (coverage) -- (thm);
\draw[arrow] (tracking) -- (thm);
\draw[arrow] (selfsign) -- (thm);
\draw[arrow] (specialmax)
    .. controls +(0,-2.0) and +(2.0,1.0) ..
    (thm.east);

\end{tikzpicture}%
}
\caption{
Dependency graph for the proof of
Theorem~\ref{thm:neurons-specialize}.
Solid arrows indicate the main logical dependencies.
The dashed arrow records that
Lemma~\ref{lem:cubic-correlation-identity}
is also used to establish smoothness of the teacher-only
vector field in the tracking argument.
}
\label{fig:theorem1-dependencies}
\end{figure}
}
\clearpage

\subsubsection{The Population-Correlation Landscape}

In this subsection, we study the critical points of the optimization problem
\begin{align}\label{eq:directional-alignment-optimization-problem}
\max_{\omega \in \mathbb{S}^{d-1}} \Psi_{\zeta}(\omega)\,.
\end{align}
The use of homogeneity to characterize parameter directions
is part of a broader literature on the implicit bias of gradient methods,
beginning with homogeneous linear predictors and extending to homogeneous
neural networks
\citep{soudry2018implicit,ji2019implicit,lyu2020gradient,ji2020directional}.
These results motivate the directional viewpoint used here, but do not
directly yield the squared-loss, small-initialization approximation in our
setting; we establish the required approximation directly in
Lemma~\ref{lem:early-time-tracking}. (In the subsequent section of the appendix, we will show how the sign $\zeta$ is determined for each neuron at very early times of training.) We show that the critical points of \eqref{eq:directional-alignment-optimization-problem} are specialized to clusters of the data distribution.

First, we provide an explicit formula for the signed population-correlation objective in terms of a weight's correlations with the cluster means and predictive directions. This lemma is stated for directions $\omega$ on the unit sphere, but is later applied to the directions of small weights using the homogeneity of ReLU networks.

\begin{lemma}[Cubic correlation identity]
\label{lem:cubic-correlation-identity}
For every $\omega\in\mathbb S^{d-1}$ and $\zeta\in\{\pm1\}$,
\[
    \Psi_\zeta(\omega)
    =
    -\frac{\zeta}{K\sqrt6}
    \sum_{c=1}^K
    b_c(\omega)\varphi\bigl(b_c(\omega)\bigr)\rho_c(\omega)^3.
\]
\end{lemma}

\begin{proof}
Fix $\omega\in\mathbb S^{d-1}$. Conditional on $c$,
\[
    \omega^\top x
    =b_c(\omega)
      +\omega^\top z,
    \qquad
    y=h_3\bigl(v_c^\top z\bigr).
\]
Hence the contribution of cluster $c$ to $\Phi(\omega) = \mathbb E\left[y\,
      \phi(\omega^\top x)\right]$ is
\[\mathbb E\left[
        h_3\bigl(v_c^\top z\bigr)
        \bigl(b_c(\omega)
          +\omega^\top z\bigr)_+
    \right].
\]

Set
\[
    U=v_c^\top z,
    \qquad
    V=\omega^\top z,
    \qquad
    \rho=\rho_c(\omega)
    =\langle v_c,\omega\rangle.
\]
Since $z\sim\mathcal N(0,I_d)$ and
$v_c,\omega$ are unit vectors,
$U$ and $V$ are standard Gaussian, and
\[
    \mathbb E[UV]
    =
    v_c^\top
      \mathbb E[zz^\top]\omega
    =
    \langle v_c,\omega\rangle
    =
    \rho.
\]
Thus $(U,V)$ is a jointly standard Gaussian pair with correlation $\rho$.

For $|\rho|<1$, define
\[
    Z:=\frac{U-\rho V}{\sqrt{1-\rho^2}}.
\]
Because $(U,V)$ is jointly Gaussian, $(Z,V)$ is also jointly Gaussian.
Moreover,
\[
    \mathbb E[Z]=0,
    \qquad
    \mathbb E[Z^2]=1,
\]
and
\[
    \mathbb E[ZV]
    =
    \frac{\mathbb E[UV]-\rho\mathbb E[V^2]}
         {\sqrt{1-\rho^2}}
    =0.
\]
Jointly Gaussian random variables with zero covariance are independent, so
$Z\sim\mathcal N(0,1)$ is independent of $V$. Therefore
\[
    U=\rho V+\sqrt{1-\rho^2}\,Z.
\]
When $|\rho|=1$, the same representation holds with the second term equal
to zero.

Conditioning on $V$ and using the independence of $Z$ gives
\[
    \mathbb E[U\mid V]=\rho V
\]
and
\begin{align*}
    \mathbb E[U^3\mid V]
    &=
    \mathbb E\left[
        \left(\rho V+\sqrt{1-\rho^2}\,Z\right)^3
        \,\middle|\,V
    \right]\\
    &=
    \rho^3V^3+3\rho(1-\rho^2)V,
\end{align*}
since $\mathbb E[Z]=\mathbb E[Z^3]=0$ and $\mathbb E[Z^2]=1$.
Recalling that
\[
    h_3(x)=\frac{x^3-3x}{\sqrt6},
\]
we obtain
\begin{align*}
    \mathbb E[h_3(U)\mid V]
    &=
    \frac{1}{\sqrt6}
    \left(
        \mathbb E[U^3\mid V]
        -3\mathbb E[U\mid V]
    \right)\\
    &=
    \rho^3 h_3(V).
\end{align*}

Since $(b_c(\omega)+V)_+$ depends only on $V$, the tower property now gives
\begin{align*}
    &\mathbb E\left[
        h_3(U)\bigl(b_c(\omega)+V\bigr)_+
    \right]\\
    &\qquad=
    \mathbb E\left[
        \bigl(b_c(\omega)+V\bigr)_+
        \mathbb E[h_3(U)\mid V]
    \right]\\
    &\qquad=
    \rho_c(\omega)^3
    \mathbb E_{G\sim\mathcal N(0,1)}
    \left[
        h_3(G)\bigl(b_c(\omega)+G\bigr)_+
    \right].
\end{align*}

It remains to compute the one-dimensional expectation. For any $b\in\mathbb R$,
\[
    \mathbb E\left[h_3(G)(b+G)_+\right]
    =
    \frac1{\sqrt6}
    \int_{-b}^{\infty}
        (x^3-3x)(b+x)\varphi(x)\,dx.
\]
Using
\[
    \frac{d}{dx}
    \left[(x^2-1)\varphi(x)\right]
    =
    -(x^3-3x)\varphi(x),
\]
integration by parts yields
\[
    \int_{-b}^{\infty}
        (x^3-3x)(b+x)\varphi(x)\,dx
    =
    \int_{-b}^{\infty}
        (x^2-1)\varphi(x)\,dx,
\]
where the boundary term vanishes because $b+x=0$ at $x=-b$ and
$\varphi(x)\to0$ as $x\to\infty$. Since
\[
    \frac{d}{dx}\bigl(x\varphi(x)\bigr)
    =
    (1-x^2)\varphi(x),
\]
we further obtain
\[
    \int_{-b}^{\infty}
        (x^2-1)\varphi(x)\,dx
    =
    -b\varphi(b).
\]
Therefore
\[
    \mathbb E\left[h_3(G)(b+G)_+\right]
    =
    -\frac{b\varphi(b)}{\sqrt6}.
\]

Applying this with $b=b_c(\omega)$, the contribution of cluster $c$ to
$\Phi(\omega)$ is
\[
    -\frac{1}{\sqrt6}
    b_c(\omega)\varphi\bigl(b_c(\omega)\bigr)\rho_c(\omega)^3.
\]
Since $c$ is uniform on $[K]$,
\[
    \Phi(\omega)
    =
    -\frac1{K\sqrt6}
    \sum_{c=1}^K
    b_c(\omega)\varphi\bigl(b_c(\omega)\bigr)\rho_c(\omega)^3.
\]
Finally, $\Psi_\zeta(\omega)=\zeta\Phi(\omega)$, and hence
\[
    \Psi_\zeta(\omega)
    =
    -\frac{\zeta}{K\sqrt6}
    \sum_{c=1}^K
    b_c(\omega)\varphi\bigl(b_c(\omega)\bigr)\rho_c(\omega)^3,
\]
as claimed.
\end{proof}


Because the $\mu_c/R$ and $v_c$ vectors
form an orthonormal basis, the
sphere constraint is
\[
    \sum_{c=1}^K\frac{b_c^2}{R^2}
    +
    \sum_{c=1}^K\rho_c^2
    =
    1.
\]
Next, we consider the optimization problem of signed population-correlation maximization between the value of the neuron and the ground truth:
$\max_{\omega\in\mathbb S^{d-1}}\Psi_\zeta(\omega)$. This is written and analyzed below under the linear change of
coordinates
\[
    b_c=\langle\omega,\mu_c\rangle,
    \qquad
    \rho_c=\langle\omega,v_c\rangle\,,
\]
where we show that the neuron maximizes correlation with the response by specializing to one cluster. Here, we use the fact that the link function is the cubic Hermite polynomial. (A similar argument would hold for all higher-degree Hermite link functions as well.)

\begin{lemma}[Positive local maxima use one cluster]
\label{lem:one-cluster-maxima}
Consider the constrained maximization problem
\[
    \max_{b_1,\ldots,b_K,\rho_1,\ldots,\rho_K}
    F(b_1,\ldots,b_K,\rho_1,\ldots,\rho_K),
\]
where
\[
    F(b_1,\ldots,b_K,\rho_1,\ldots,\rho_K)
    :=
    -\frac{\zeta}{K\sqrt6}
    \sum_{c=1}^K
    b_c\varphi(b_c)\rho_c^3,
\]
subject to
\[
    \sum_{c=1}^K\frac{b_c^2}{R^2}
    +
    \sum_{c=1}^K\rho_c^2
    =
    1.
\]

Every constrained local maximum of $F$ with positive objective value has
exactly one nonzero predictive coordinate $\rho_c$. Moreover, if
$\rho_c=0$, then the corresponding routing coordinate $b_c$ also vanishes.
\end{lemma}

\begin{proof}
At a constrained local maximum, first-order stationarity of the Lagrangian
\[
    F(b_1,\ldots,b_K,\rho_1,\ldots,\rho_K)
    -
    \lambda\left(
        \sum_{c=1}^K\frac{b_c^2}{R^2}
        +
        \sum_{c=1}^K\rho_c^2
        -1
    \right)
\]
with respect to $b_c$ and $\rho_c$ gives, for every $c\in[K]$,
\begin{align}
    -\frac{\zeta}{K\sqrt6}
    (1-b_c^2)\varphi(b_c)\rho_c^3
    &=
    \frac{2\lambda}{R^2}b_c,
    \label{eq:kkt-b}\tag{B.1}\\
    -\frac{3\zeta}{K\sqrt6}
    b_c\varphi(b_c)\rho_c^2
    &=
    2\lambda\rho_c.
    \label{eq:kkt-rho}\tag{B.2}
\end{align}
Here we used
\[
    \frac{d}{db}\bigl(b\varphi(b)\bigr)
    =
    (1-b^2)\varphi(b).
\]

Multiplying~\eqref{eq:kkt-rho} by $\rho_c$ and summing over $c$ gives
\[
    3F(b_1,\ldots,b_K,\rho_1,\ldots,\rho_K)
    =
    2\lambda\sum_{c=1}^K\rho_c^2.
\]
Since the objective value is positive, at least one predictive coordinate
is nonzero. Hence both the left-hand side and
$\sum_c\rho_c^2$ are positive, so
\[
    \lambda>0.
\]

If $\rho_c=0$, then~\eqref{eq:kkt-b} reduces to
\[
    0=\frac{2\lambda}{R^2}b_c.
\]
Since $\lambda>0$, this implies $b_c=0$.

It remains to rule out two nonzero predictive coordinates. Suppose that
$\rho_c\neq0$ and $\rho_r\neq0$ for two distinct clusters $c\neq r$.
Keep all routing coordinates and all other predictive coordinates fixed,
and consider
\[
    \rho_c(t)
    =
    \operatorname{sign}(\rho_c)
    \sqrt{\rho_c^2+t},
    \qquad
    \rho_r(t)
    =
    \operatorname{sign}(\rho_r)
    \sqrt{\rho_r^2-t}.
\]
For sufficiently small $|t|$, this perturbation is well defined, preserves
the signs of the two coordinates, and satisfies
\[
    \rho_c(t)^2+\rho_r(t)^2
    =
    \rho_c^2+\rho_r^2.
\]
Thus it preserves the constraint exactly.

Differentiating $F$ along this feasible curve at $t=0$ gives
\[
    F'(0)
    =
    -\frac{3\zeta}{2K\sqrt6}
    \left(
        b_c\varphi(b_c)\rho_c
        -
        b_r\varphi(b_r)\rho_r
    \right).
\]
Since $\rho_c,\rho_r\neq0$, dividing~\eqref{eq:kkt-rho} by the corresponding
predictive coordinate gives
\[
    -\frac{3\zeta}{K\sqrt6}
    b_c\varphi(b_c)\rho_c
    =
    2\lambda,
    \qquad
    -\frac{3\zeta}{K\sqrt6}
    b_r\varphi(b_r)\rho_r
    =
    2\lambda.
\]
Therefore
\[
    F'(0)=0.
\]

The second derivative along the same feasible curve is
\[
    F''(0)
    =
    -\frac{3\zeta}{4K\sqrt6}
    \left(
        \frac{b_c\varphi(b_c)}{\rho_c}
        +
        \frac{b_r\varphi(b_r)}{\rho_r}
    \right).
\]
Using the same KKT identities gives
\[
    F''(0)
    =
    \frac{\lambda}{2}
    \left(
        \frac{1}{\rho_c^2}
        +
        \frac{1}{\rho_r^2}
    \right)
    >0.
\]
Thus a critical point with two nonzero predictive coordinates has a feasible
direction of positive curvature and cannot be a constrained local maximum.

Hence at most one predictive coordinate is nonzero. Since the objective
value is positive, at least one is nonzero, so exactly one is active. As
shown above, every routing coordinate corresponding to an inactive predictive
coordinate also vanishes.
\end{proof}

Returning to our objective $\max_{\omega \in \mathbb{S}^{d-1}} \Psi_{\zeta}(\omega)$, let $c$ denote the unique active cluster and let
\[
    \tau=\operatorname{sign}(\rho_c)\in\{\pm1\}.
\]
Then
\[
    b_j=\rho_j=0
    \qquad
    \text{for every }j\neq c,
\]
and the constraint becomes
\[
    \frac{b_c^2}{R^2}+\rho_c^2=1.
\]
Hence
\[
    \rho_c
    =
    \tau\sqrt{1-\frac{b_c^2}{R^2}}.
\]
Substituting this into $F$, and writing $b=b_c$, reduces the problem to the
one-dimensional objective
\[
    -\frac{\tau\zeta}{K\sqrt6}
    b\varphi(b)
    \left(1-\frac{b^2}{R^2}\right)^{3/2},
    \qquad
    |b|<R.
\]


\begin{lemma}[Specialized maxima]
\label{lem:specialized-maxima}
Fix $c\in[K]$ and $\zeta,\tau\in\{\pm1\}$. On the one-cluster feasible branch,
write
\[
    f(b)
    :=
    -\frac{\tau\zeta}{K\sqrt6}
    b\varphi(b)
    \left(1-\frac{b^2}{R^2}\right)^{3/2},
    \qquad |b|<R,
\]
and let
\[
    I:=\{b\in(-R,R):f(b)>0\}.
\]
Then $I$ is a connected open interval. For all sufficiently large $R$, $f$
has a unique critical point $b^\star\in I$, which is a strict local maximum
and satisfies
\[
    b^\star=-\tau\zeta+O(R^{-2}).
\]
The corresponding direction
\[
    \omega^\star
    =
    \tau\sqrt{1-\frac{(b^\star)^2}{R^2}}\,v_c
    +
    \frac{b^\star}{R^2}\mu_c
\]
is a strict local maximum of $\Psi_\zeta$ and obeys
\[
    \omega^\star
    =
    \tau v_c
    -
    \frac{\tau\zeta}{R^2}\mu_c
    +
    O(R^{-2}).
\]
\end{lemma}

\begin{proof}
For every $b\in(-R,R)$,
\[
    \frac{1}{K\sqrt6}>0,
    \qquad
    \varphi(b)>0,
    \qquad
    \left(1-\frac{b^2}{R^2}\right)^{3/2}>0.
\]
Hence the sign of $f(b)$ is determined entirely by $-\tau\zeta b$:
\[
    f(b)>0
    \quad\Longleftrightarrow\quad
    -\tau\zeta b>0.
\]
Since $\tau\zeta\in\{\pm1\}$, it follows that
\[
    I=
    \begin{cases}
        (-R,0), & \tau\zeta=1,\\
        (0,R), & \tau\zeta=-1.
    \end{cases}
\]
Thus $I$ is connected.

For $b\in I$, set
\[
    t=-\tau\zeta b.
\]
The preceding characterization of $I$ shows that this is a bijective linear
change of variable from $I$ onto $(0,R)$. Since $\varphi$ is even and
$(\tau\zeta)^2=1$, we obtain
\[
    f(b)
    =
    \frac{1}{K\sqrt6}
    t\varphi(t)
    \left(1-\frac{t^2}{R^2}\right)^{3/2}.
\]
Therefore the critical points of $f$ on $I$ correspond exactly to the
critical points on $(0,R)$ of
\[
    t\longmapsto
    t\varphi(t)
    \left(1-\frac{t^2}{R^2}\right)^{3/2}.
\]

This function is strictly positive on $(0,R)$, so its critical points can be
found from its logarithmic derivative:
\[
    \frac{d}{dt}
    \log\left[
        t\varphi(t)
        \left(1-\frac{t^2}{R^2}\right)^{3/2}
    \right]
    =
    \frac1t-t-\frac{3t}{R^2-t^2}.
\]
Thus a critical point satisfies
\[
    \frac1t-t-\frac{3t}{R^2-t^2}=0.
\]
Multiplying by $t(R^2-t^2)$ gives
\[
    t^4-(R^2+4)t^2+R^2=0.
\]
Solving this quadratic equation in $t^2$ gives
\[
    t^2
    =
    \frac{R^2+4\pm\sqrt{R^4+4R^2+16}}{2}.
\]
The root with the plus sign is larger than $R^2$, whereas the root with the
minus sign lies in $(0,R^2)$. Hence there is exactly one critical point in
$(0,R)$, and it satisfies
\[
    t^2
    =
    \frac{
        R^2+4-\sqrt{R^4+4R^2+16}
    }{2}
    =
    1+O(R^{-2}).
\]
Since $t>0$,
\[
    t=1+O(R^{-2}).
\]
Returning to $b=-\tau\zeta t$ gives
\[
    b^\star
    =
    -\tau\zeta+O(R^{-2}).
\]



The function
\[
    t\varphi(t)
    \left(1-\frac{t^2}{R^2}\right)^{3/2}
\]
is positive on $(0,R)$ and tends to zero as $t\to0$ or $t\to R$.
Since it has exactly one critical point in $(0,R)$, this point is its
unique strict maximum. Hence $b^\star$ is the unique maximizer of $f$ on
$I$.

It remains to show that the corresponding point is a local maximum of the
full constrained objective $F$. Since the feasible set is compact, $F$
attains a global maximum. The point constructed above has positive objective
value, so the global maximum is positive. By
Lemma~\ref{lem:one-cluster-maxima}, every global maximizer must have exactly
one nonzero predictive coordinate and therefore lies on one of the
one-cluster feasible branches.

For every choice of the active cluster $c$ and orientation $\tau$, the
change of variable
\[
    t=-\tau\zeta b
\]
reduces the positive part of the corresponding one-cluster objective to
\[
    \frac{1}{K\sqrt6}
    t\varphi(t)
    \left(1-\frac{t^2}{R^2}\right)^{3/2}.
\]
Thus all one-cluster branches have the same maximal value, attained uniquely
at the point identified above. Consequently, every corresponding
$\omega^\star$ is a global maximizer of $\Psi_\zeta$.

There are only finitely many such maximizers, one for each
$c\in[K]$ and $\tau\in\{\pm1\}$. Hence each is isolated, and therefore each
$\omega^\star$ is a strict local maximum of $\Psi_\zeta$.

\end{proof}


\subsubsection{Random Initialization, Self-Selected Signs, and Cluster Coverage}

In the previous section, we considered the signed population-correlation objective $\Psi_{\zeta}$ for each neuron. In this section, we note that the sign $\zeta$ for each neuron is learned at very early times of training and does not need to be fixed. Consider
the teacher-only joint dynamics
\begin{equation}
    \dot u=\Phi({\omega}),
    \qquad
    {\dot\omega}
    =\tanh(u)\nabla_{\mathbb S^{d-1}}\Phi({\omega}).
\end{equation}
These equations arise from the two-layer parameterization when the
second-layer weight is initially zero. Crucially, they neglect interactions between neurons, which we will later show is a fine approximation because the network is small at initialization.

\begin{lemma}[Positive signed-correlation flows specialize]
\label{lem:positive-signed-flow}
Let ${\omega^0}$ have an absolutely continuous distribution on
$\mathbb S^{d-1}$, fix $\zeta\in\{\pm1\}$, and assume
$\Psi_\zeta({\omega^0})>0$. For all sufficiently large $R$, spherical gradient
ascent
\[
    {\dot\omega}
    =\nabla_{\mathbb S^{d-1}}\Psi_\zeta({\omega}),
    \qquad
    {\omega(0)}={\omega^0},
\]
converges almost surely to one of the specialized maxima in
Lemma~\ref{lem:specialized-maxima}.
\end{lemma}



\begin{proof}
Along spherical gradient ascent,
\[
    \frac{d}{dt}\Psi_\zeta({\omega(t)})
    =
    \left\|
        \nabla_{\mathbb S^{d-1}}\Psi_\zeta({\omega(t)})
    \right\|_2^2
    \geq 0.
\]
Since the sphere is compact, $\Psi_\zeta$ is bounded above. Hence
\[
    \int_0^\infty
    \left\|
        \nabla_{\mathbb S^{d-1}}\Psi_\zeta({\omega(t)})
    \right\|_2^2\,dt
    <\infty.
\]
Moreover, $\Psi_\zeta$ is smooth on the sphere, so its gradient and Hessian
are bounded. It follows that
$\|\nabla_{\mathbb S^{d-1}}\Psi_\zeta({\omega(t)})\|_2^2$ has bounded derivative.
A nonnegative integrable function with bounded derivative must converge to
zero, and therefore
\[
    \left\|
        \nabla_{\mathbb S^{d-1}}\Psi_\zeta({\omega(t)})
    \right\|_2
    \longrightarrow 0.
\]

We next note that $\Psi_\zeta$ has only finitely many positive critical
points. Indeed, at any positive critical point, the stationarity conditions
\eqref{eq:kkt-b}--\eqref{eq:kkt-rho} imply that $\lambda>0$, and hence
$b_c=0$ whenever $\rho_c=0$. For every active coordinate $\rho_c\neq0$,
eliminating $\lambda$ from the two stationarity equations gives
\[
    (1-b_c^2)\rho_c^2
    =
    \frac{3b_c^2}{R^2}.
\]
Thus $0<|b_c|<1$ and
\[
    \rho_c^2
    =
    \frac{3b_c^2}{R^2(1-b_c^2)}.
\]
Equation~\ref{eq:kkt-rho} further shows that
\[
    |b_c|\varphi(b_c)|\rho_c|
\]
has the same value for every active coordinate. Substituting the preceding
expression for $\rho_c^2$, this quantity is proportional to
\[
    \frac{|b_c|^2\varphi(|b_c|)}
         {\sqrt{1-|b_c|^2}},
\]
which is strictly increasing for $0<|b_c|<1$, since
\[
    \frac{d}{dt}
    \log\left(
        \frac{t^2\varphi(t)}{\sqrt{1-t^2}}
    \right)
    =
    \frac{2}{t}+\frac{t^3}{1-t^2}
    >0.
\]
Hence all active $|b_c|$ are equal. If there are $s$ active coordinates,
the sphere constraint then gives
\[
    s\,\frac{b_c^2(4-b_c^2)}
             {R^2(1-b_c^2)}
    =1,
\]
which has a unique solution for $b_c^2\in(0,1)$. Thus, for each choice of
the active coordinates and their signs, there is at most one positive
critical point. Since there are only finitely many such choices, the set of
positive critical points is finite.

Because $\Psi_\zeta({\omega(0)})>0$ and the objective is nondecreasing, every
accumulation point of the trajectory has positive objective value. By
compactness, accumulation points exist, and the convergence of the gradient
to zero implies that every accumulation point is a positive critical point.
Since there are only finitely many such points, the trajectory must converge
to one of them.

Finally, the argument in Lemma~\ref{lem:one-cluster-maxima} gives an
unstable tangent direction at every positive critical point with more than
one active predictive coordinate. Since there are only finitely many such
critical points, the center-stable manifold theorem implies that the set of
initial conditions converging to any of them has measure zero.
Lemma~\ref{lem:specialized-maxima} shows that the remaining positive critical
points are precisely the specialized maxima. Hence an absolutely continuous
initialization converges almost surely to a specialized maximum.
\end{proof}

For a single neuron, write
\[
    \omega(t)=\frac{w(t)}{\|w(t)\|_2}.
\]
ReLU homogeneity implies that gradient flow preserves
\[
    \|w(t)\|_2^2-a(t)^2
    =
    \|w(0)\|_2^2-a(0)^2
    =
    \varepsilon^2.
\]
We therefore define the scalar amplitude coordinate
\[
    u(t)
    :=
    \operatorname{arsinh}\!\left(\frac{a(t)}{\varepsilon}\right),
\]
so that, equivalently,
\[
    a(t)=\varepsilon\sinh u(t),
    \qquad
    \|w(t)\|_2=\varepsilon\cosh u(t).
\]
Under this parametrization, the target-only part of the directional dynamics
is
\begin{equation}
\label{eq:self-selected-teacher-flow}
    \dot u=\Phi(\omega),
    \qquad
    \dot\omega
    =
    \tanh(u)\nabla_{\mathbb S^{d-1}}\Phi(\omega).
\end{equation}

\begin{lemma}[Self-selected sign and almost-sure specialization]
\label{lem:self-selected-specialization}
Let $\omega^0$ have an absolutely continuous distribution on
$\mathbb S^{d-1}$, and let $(u(t),\omega(t))$ solve
\eqref{eq:self-selected-teacher-flow} with
\[
    u(0)=0,
    \qquad
    \omega(0)=\omega^0.
\]
For almost every $\omega^0$, define
\[
    \zeta
    =
    \operatorname{sign}(\Phi(\omega^0))
    \in\{\pm1\}.
\]
Then, for every $t>0$,
\[
    \zeta u(t)>0,
    \qquad
    \zeta\Phi(\omega(t))
    \geq
    |\Phi(\omega^0)|.
\]
Moreover, $\omega(t)$ is a positive time reparameterization of spherical
gradient ascent on $\Psi_\zeta=\zeta\Phi$, and, for all sufficiently large
$R$, it converges almost surely to one of the specialized maxima in
Lemma~\ref{lem:specialized-maxima}.
\end{lemma}


\begin{proof}
Since $\Phi$ is real analytic and not identically zero, its zero set has
spherical measure zero. Thus, for almost every ${\omega^0}$, $\Phi({\omega^0})\neq0$.
Fix such a ${\omega^0}$ and let
\[
    \zeta=\operatorname{sign}(\Phi({\omega^0})).
\]

At $t=0$,
\[
    \frac{d}{dt}\bigl(\zeta u(t)\bigr)\bigg|_{t=0}
    =
    \zeta\Phi({\omega^0})
    =
    |\Phi({\omega^0})|
    >0,
\]
so $\zeta u(t)>0$ for all sufficiently small $t>0$.

As long as $\zeta u(t)>0$, $\tanh(u(t))$ has sign $\zeta$, and hence
\[
    \zeta\tanh(u(t))=|\tanh(u(t))|.
\]
Therefore
\[
    \frac{d}{dt}\bigl(\zeta\Phi({\omega(t)})\bigr)
    =
    |\tanh(u(t))|
    \left\|
        \nabla_{\mathbb S^{d-1}}\Phi({\omega(t)})
    \right\|_2^2
    \geq0.
\]
It follows that
\[
    \zeta\Phi({\omega(t)})
    \geq
    \zeta\Phi({\omega^0})
    =
    |\Phi({\omega^0})|.
\]
Consequently,
\[
    \frac{d}{dt}\bigl(\zeta u(t)\bigr)
    =
    \zeta\Phi({\omega(t)})
    \geq
    |\Phi({\omega^0})|,
\]
and therefore
\[
    \zeta u(t)\geq |\Phi({\omega^0})|t.
\]
In particular, $\zeta u(t)$ cannot return to zero, so the preceding
inequalities hold for every $t>0$. This proves
\[
    \zeta u(t)>0,
    \qquad
    \zeta\Phi({\omega(t)})
    \geq
    |\Phi({\omega^0})|.
\]

Since $\Psi_\zeta=\zeta\Phi$ and $\tanh(u(t))$ has sign $\zeta$, the
directional dynamics satisfy
\[
    {\dot\omega(t)}
    =
    |\tanh(u(t))|
    \nabla_{\mathbb S^{d-1}}\Psi_\zeta({\omega(t)}).
\]
Thus ${\omega(t)}$ follows the same orbit as spherical gradient ascent on
$\Psi_\zeta$, up to a positive reparameterization of time. Moreover,
$\zeta u(t)\geq|\Phi({\omega^0})|t$ implies $|u(t)|\to\infty$, and hence
$|\tanh(u(t))|\to1$. In particular,
\[
    \int_0^\infty |\tanh(u(t))|\,dt=\infty,
\]
so this reparameterization covers the entire forward gradient-ascent
trajectory.

Finally,
\[
    \Psi_\zeta({\omega^0})
    =
    \zeta\Phi({\omega^0})
    =
    |\Phi({\omega^0})|
    >0.
\]
Lemma~\ref{lem:positive-signed-flow} therefore implies that ${\omega(t)}$
converges almost surely to one of the specialized maxima in
Lemma~\ref{lem:specialized-maxima}.
\end{proof}



\begin{lemma}[Uniform selected labels and coverage]
\label{lem:uniform-cluster-coverage}
Let ${\omega_1^0,\ldots,\omega_m^0}$ be independent and uniform on
$\mathbb S^{d-1}$, and let $J_j$ be the cluster selected by the self-selected
flow in Lemma~\ref{lem:self-selected-specialization}. Then
$J_1,\ldots,J_m$ are independent and uniform on $[K]$. Consequently,
\[
    \mathbb P\bigl(\{J_1,\ldots,J_m\}\neq[K]\bigr)
    \leq
    K\left(1-\frac1K\right)^m
    \leq
    Ke^{-m/K}.
\]
\end{lemma}

\begin{proof}
For a permutation $\pi$ of $[K]$, let $P_\pi$ be the orthogonal map satisfying
\[
    P_\pi{\frac{\mu_c}{R}}
    =
    {\frac{\mu_{\pi(c)}}{R}},
    \qquad
    P_\pi{v_c}
    =
    {v_{\pi(c)}}.
\]
The data distribution, target, and self-selected vector field are equivariant
under $P_\pi$. Hence, if ${\omega^0}$ selects cluster $c$, then $P_\pi {\omega^0}$ selects
cluster $\pi(c)$. Uniform spherical measure is invariant under $P_\pi$, so
all selection probabilities are equal. By
Lemma~\ref{lem:self-selected-specialization}, they sum to one, and each is
therefore $1/K$. Independence follows because each teacher-only label is a
deterministic function of an independent initial direction.

A fixed cluster is missed with probability $(1-1/K)^m$. A union bound yields
the coverage estimate.
\end{proof}

\subsubsection{Tracking by the Full Small-Initialization Dynamics}

Finally, we put the above ingredients analyzing the trajectories of individual neurons together, and show that they describe the trajectory of the neurons in a neural network (up to rescaling) when the network is initialized small.
The proof strategy of studying feature learning through
effectively independent neuron dynamics in early-time or
small-initialization regimes has been put forward in prior work
\citep{abbe2022merged,min2024early,glasgow2024sgd}.
The population loss is
\[
    \mathcal L(\theta)
    =
    \frac12
    \mathbb E\left[
        \bigl(f_\theta({x})-{y}\bigr)^2
    \right],
\]
and its gradient-flow equations are
\begin{align*}
    \dot a_j
    &=
    \mathbb E\left[
        ({y}-f_\theta({x}))
        \phi(w_j^\top{x})
    \right],\\
    \dot w_j
    &=
    a_j
    \mathbb E\left[
        ({y}-f_\theta({x}))
        \mathbf 1_{\{w_j^\top{x}>0\}}{x}
    \right].
\end{align*}

\begin{lemma}[Early-time tracking from a zero output layer]
\label{lem:early-time-tracking}
Initialize
\[
    a_j(0)=0,
    \qquad
    w_j(0)=\varepsilon {\omega_j^0},
    \qquad
    {\omega_j^0}\in\mathbb S^{d-1}.
\]
Let $(\bar u_j,{\bar\omega_j})$ solve the self-selected teacher-only system
\[
    \dot{\bar u}_j
    =
    \Phi({\bar\omega_j}),
    \qquad
    {\dot{\bar\omega}_j}
    =
    \tanh(\bar u_j)
    \nabla_{\mathbb S^{d-1}}\Phi({\bar\omega_j}),
\]
with
\[
    (\bar u_j(0),{\bar\omega_j(0)})
    =
    (0,{\omega_j^0}).
\]
For every fixed $T<\infty$, there are constants $L_T<\infty$ and
$\varepsilon_T>0$ such that, for $0<\varepsilon\leq\varepsilon_T$,
\[
    \sup_{0\leq t\leq T}
    \left(
        |u_j(t)-\bar u_j(t)|
        +
        \|{\omega_j(t)-\bar\omega_j(t)}\|_2
    \right)
    \leq
    L_Tm\varepsilon^2,
\]
where
\[
    a_j(t)=\varepsilon\sinh(u_j(t)),
    \qquad
    {\omega_j(t)}=\frac{w_j(t)}{\|w_j(t)\|_2}.
\]
\end{lemma}

\begin{proof}
ReLU homogeneity gives the invariant
\[
    \frac{d}{dt}
    \left(
        \|w_j\|_2^2-a_j^2
    \right)
    =
    0.
\]
Under the stated initialization,
\[
    \|w_j(t)\|_2^2-a_j(t)^2
    =
    \varepsilon^2.
\]
Since $a_j=\varepsilon\sinh u_j$, it follows that
\[
    \|w_j\|_2
    =
    \varepsilon\cosh u_j.
\]

Substituting
\[
    w_j
    =
    \varepsilon\cosh(u_j){\omega_j}
\]
into the gradient-flow equations and using ReLU homogeneity gives
\begin{align*}
    \dot u_j
    ={}&
    \Phi({\omega_j})
    -
    \mathbb E\left[
        f_\theta({x})
        \phi({\omega_j}^\top{x})
    \right],\\
    {\dot\omega_j}
    ={}&
    \tanh(u_j)
    \Bigg(
        \nabla_{\mathbb S^{d-1}}\Phi({\omega_j})\\
    &\hspace{25mm}
        -
        \bigl(I-{\omega_j\omega_j^\top}\bigr)
        \mathbb E\left[
            f_\theta({x})
            \mathbf 1_{\{{\omega_j}^\top
              {x}>0\}}{x}
        \right]
    \Bigg).
\end{align*}

Lemma~\ref{lem:cubic-correlation-identity} shows that the teacher-only
vector field is smooth on bounded $u$-intervals and on the sphere. On every
fixed interval $[0,T]$, the quantities $u_j(t)$ remain uniformly bounded for
sufficiently small $\varepsilon$. Consequently, there is a constant
$C_T<\infty$ such that
\[
    |a_j(t)|+\|w_j(t)\|_2
    \leq
    C_T\varepsilon,
    \qquad
    0\leq t\leq T.
\]

Since ${x}$ has finite moments of every order,
\[
    \|f_{\theta(t)}\|_{L^2(P_{{x}})}
    \leq
    C_T
    \sum_{j=1}^m
    |a_j(t)|\|w_j(t)\|_2
    \leq
    C_Tm\varepsilon^2.
\]
Cauchy--Schwarz therefore bounds both interaction terms in the
$(u_j,{\omega_j})$ equations by $C_Tm\varepsilon^2$.

The full and teacher-only systems have the same initial conditions, and
their vector fields are Lipschitz on the relevant compact set. Gr\"onwall's
inequality yields
\[
    \sup_{0\leq t\leq T}
    \left(
        |u_j(t)-\bar u_j(t)|
        +
        \|{\omega_j(t)-\bar\omega_j(t)}\|_2
    \right)
    \leq
    L_Tm\varepsilon^2.
\]
\end{proof}

\paragraph{Perturbations of the cubic link function.}
The Hermite link function isolates the specialization mechanism and makes the
correlation landscape exact. More generally, one could perturb the induced
correlation objective and expect specialization to persist.
In particular, the specialized maxima and their attraction basins should be
stable under sufficiently small perturbations near both the maxima and the
strict-saddle regions.




\subsection{Proof of Theorem~\ref{thm:rfm-mlp-gap}}
\label{app:proof-mlp-rfm-separation}


We first prove the positive MLP result in a more general fixed-dimensional
multi-index setting. We then recover the single-index MLP result in
Theorem~\ref{thm:rfm-mlp-gap} as a corollary and prove the kernel and RFM lower bounds in that
single-index setting.

\subsubsection{A General Multi-Index Specialized MLP Construction}

Fix an integer $r_0\geq1$ and a constant $R\geq1$. For each $K$, let
$d=K+Kr_0$. Let $E_c\in\mathbb R^{Kr_0\times r_0}$ embed
$\mathbb R^{r_0}$ into the $c$th block of $\mathbb R^{Kr_0}$, so that
\[
    E_c^\top E_{c'}=\mathbf 1_{\{c=c'\}}I_{r_0}.
\]
Define
\[
    \mu_c=
    \begin{bmatrix}
        Re_c\\
        0
    \end{bmatrix}
    \in\mathbb R^d,
    \qquad
    V_c=
    \begin{bmatrix}
        0\\
        E_c
    \end{bmatrix}
    \in\mathbb R^{d\times r_0}.
\]
The data distribution is
\[
    C\sim\operatorname{Unif}([K]),
    \qquad
    X=\mu_C+
    \begin{bmatrix}
        0\\
        Z
    \end{bmatrix},
    \qquad
    Z\sim\mathcal N(0,I_{Kr_0}).
\]
Let $g:\mathbb R^{r_0}\to\mathbb R$ be fixed, bounded, and Lipschitz, with
\[
    \|g\|_\infty\leq G_0,
    \qquad
    |g(u)-g(u')|\leq L\|u-u'\|_2,
\]
and define
\[
    f_{r_0}^\star(X)=g(V_C^\top X).
\]
Thus, within cluster $c$, the target may depend on all $r_0$ coordinates in
the cluster-specific predictive block $V_c^\top X$.

For a two-layer ReLU network, define the path norm
\[
    \|f\|_{\mathcal P}
    =\inf\left\{
        \sum_j|a_j|\|w_j\|_2:
        f(x)=\sum_j a_j\phi(w_j^\top x)
    \right\}.
\]
For a width-$m$ parameterization
$\theta=((a_j,w_j))_{j=1}^m$, define
\[
    \|\theta\|_{\mathrm F}^2
    :=\frac12\sum_{j=1}^m
        \left(a_j^2+\|w_j\|_2^2\right),
\]
and let
\[
    \mathcal F_{B,m}
    :=
    \left\{
        f_\theta(x)=\sum_{j=1}^m a_j\phi(w_j^\top x):
        \|\theta\|_{\mathrm F}^2\leq B
    \right\}.
\]

Define the ramp function
\[
    \psi(t):=\phi(t)-\phi(t-1).
\]
Every occurrence of $\psi$ below is implemented by two ordinary ReLU
neurons; the network itself has neither hidden biases nor output clipping.

\begin{lemma}[Ridge-ramp approximation]
\label{lem:ridge-ramp-density}
Let $\gamma_{r_0}=\mathcal N(0,I_{r_0})$. Finite ridge-ramp expansions of the form
\[
    p(u)=b+\sum_{j=1}^J a_j\psi(q_j^\top u-t_j)
\]
are dense in $L^2(\gamma_{r_0})$. Consequently, there is a sequence
$(p_s)_{s\geq1}$ of such functions satisfying
\[
    \|p_s-g\|_{L^2(\gamma_{r_0})}\longrightarrow0.
\]
\end{lemma}

\begin{proof}
Suppose $h\in L^2(\gamma_{r_0})$ is orthogonal to every ridge-ramp function,
and define the finite signed measure
\[
    \nu(A):=\int_A h(u)\,d\gamma_{r_0}(u).
\]
For every $q\neq0$, $t\in\mathbb R$, and $\lambda>0$,
\[
    \int h(u)\psi\!\left(\lambda(q^\top u-t)\right)
        \,d\gamma_{r_0}(u)=0.
\]
Because $0\leq\psi\leq1$, dominated convergence gives, as
$\lambda\to\infty$,
\[
    \nu\bigl(\{u:q^\top u>t\}\bigr)=0.
\]
Halfspaces determine finite Borel measures on $\mathbb R^{r_0}$, so
$\nu=0$ and therefore $h=0$ in $L^2(\gamma_{r_0})$. This proves density.
\end{proof}

\begin{lemma}[Routing a multi-index approximant]
\label{lem:routed-multi-index-approximant}
\label{lem:exact-specialized-mlp}
Let
\[
    p(u)=b+\sum_{j=1}^J a_j\psi(q_j^\top u-t_j).
\]
There is a two-layer ReLU network
\[
    F_{K,p}=\sum_{c=1}^K F_{K,p,c}
\]
with no hidden biases such that, on an input from cluster $c_0$,
\[
    F_{K,p,c}(x)
    =\mathbf 1_{\{c=c_0\}}p(V_c^\top x).
\]
In particular, each cluster-specific group cancels on all other clusters.
Every hidden weight in $F_{K,p,c}$ lies in
\[
    \operatorname{span}\bigl(\{\mu_c\}\cup\operatorname{range}(V_c)\bigr).
\]
The network has width at most $K(2J+1)$ and path norm at most
\[
    K\,\mathcal C_R(p),
\]
where
\begin{align*}
    \mathcal C_R(p)
    &:={|b|\over R}\\
    &\quad+
    \sum_{j=1}^J |a_j|
    \left[
        \sqrt{\|q_j\|_2^2+{t_j^2\over R^2}}
        +
        \sqrt{\|q_j\|_2^2+{(t_j+1)^2\over R^2}}
    \right].
\end{align*}
It also admits a parameterization satisfying
\[
    \|\theta\|_{\mathrm F}^2\leq K\,\mathcal C_R(p).
\]
\end{lemma}

\begin{proof}
Define
\begin{align*}
    F_{K,p,c}(x)
    :=\;&
        b\,\phi\!\left({\mu_c^\top x\over R^2}\right)\\
    &+\sum_{j=1}^J a_j
        \phi\!\left(
            \left(V_cq_j-{t_j\over R^2}\mu_c\right)^\top x
        \right)\\
    &-\sum_{j=1}^J a_j
        \phi\!\left(
            \left(V_cq_j-{t_j+1\over R^2}\mu_c\right)^\top x
        \right).
\end{align*}
On the support of the data distribution,
\[
    {\mu_c^\top x\over R^2}=\mathbf1_{\{C=c\}}.
\]
Suppose $x$ belongs to cluster $c_0$. If $c\neq c_0$, the two ReLUs in
each pair have the same argument and cancel, while the routing neuron is
zero. If $c=c_0$, then
\[
    F_{K,p,c_0}(x)
    =b+\sum_{j=1}^J a_j
      \left[
        \phi(q_j^\top V_{c_0}^\top x-t_j)
        -\phi(q_j^\top V_{c_0}^\top x-t_j-1)
      \right]
    =p(V_{c_0}^\top x).
\]
This proves the specialization claim.

The routing weight has norm $1/R$. The two hidden weights associated with
$(a_j,q_j,t_j)$ have norms
\[
    \sqrt{\|q_j\|_2^2+{t_j^2\over R^2}}
    \quad\text{and}\quad
    \sqrt{\|q_j\|_2^2+{(t_j+1)^2\over R^2}}.
\]
Summing the corresponding path-norm contributions over all clusters gives
the stated bound. Finally, positive homogeneity allows each neuron to be
balanced so that $|a_j|=\|w_j\|_2$. For this parameterization, its
contribution to $\|\theta\|_{\mathrm F}^2$ equals its contribution to the
path norm.
\end{proof}

The finite-network path norm is the discrete analogue of the Barron or
variation norm for two-layer networks
\citep{bach2017breaking,e2022barron}.
The following sample-dependent estimate is a version of standard
path-norm Rademacher-complexity bounds
\citep{e2022barron,golowich2018size}.

\begin{lemma}[Path-norm generalization without clipping]
\label{lem:path-rademacher}
For any deterministic sample $x_1,\ldots,x_n\in\mathbb R^d$,
\[
    \widehat{\mathfrak R}_n(\mathcal F_{B,m})
    \leq
    {2B\over n}
    \left(\sum_{i=1}^n\|x_i\|_2^2\right)^{1/2}.
\]
For the multi-index distribution above, let
$\widehat f_{\mathrm{MLP}}^{(r_0)}$ be an empirical squared-loss minimizer
over $\mathcal F_{B,m}$. There is a universal constant $C<\infty$ such that
\begin{align*}
    \mathbb E\left[
        \|\widehat f_{\mathrm{MLP}}^{(r_0)}-f_{r_0}^\star
        \|_{L^2(P_x)}^2
    \right]
    &\leq
    \inf_{f\in\mathcal F_{B,m}}
    \|f-f_{r_0}^\star\|_{L^2(P_x)}^2\\
    &\quad+
    CB^2\sqrt{
        {(Kr_0+R^2)(Kr_0+R^2+\log(en))\over n}
    }\\
    &\quad+
    CBG_0\sqrt{{Kr_0+R^2\over n}}.
\end{align*}
\end{lemma}

\begin{proof}
By the arithmetic--geometric mean inequality, every
$f\in\mathcal F_{B,m}$ has path norm at most $B$. By positive homogeneity,
every such function can be represented as
\[
    f(x)=\sum_j c_j\phi(\omega_j^\top x),
    \qquad
    \|\omega_j\|_2=1,
    \qquad
    \sum_j|c_j|\leq B.
\]
Thus $\mathcal F_{B,m}$ is contained in $B$ times the absolutely convex hull
of unit-norm ReLU atoms. Contraction and Cauchy--Schwarz give
\begin{align*}
    \widehat{\mathfrak R}_n(\mathcal F_{B,m})
    &\leq
    {B\over n}\mathbb E_\epsilon
    \sup_{\|\omega\|_2\leq1}
    \left|
        \sum_{i=1}^n\epsilon_i\phi(\omega^\top x_i)
    \right|\\
    &\leq
    {2B\over n}\mathbb E_\epsilon
    \left\|\sum_{i=1}^n\epsilon_i x_i\right\|_2\\
    &\leq
    {2B\over n}
    \left(\sum_{i=1}^n\|x_i\|_2^2\right)^{1/2}.
\end{align*}

For the squared-loss bound, let
\[
    M_n:=\max_{1\leq i\leq n}\|x_i\|_2.
\]
The path-norm bound implies
\[
    |f(x_i)|\leq BM_n
    \qquad\text{for every }f\in\mathcal F_{B,m}.
\]
Since $|y_i|\leq G_0$, the squared loss is
$2(BM_n+G_0)$-Lipschitz in the prediction over the range attained on the
sample. The standard ERM symmetrization and contraction argument therefore
gives
\begin{align*}
    \mathbb E\left[
        \|\widehat f_{\mathrm{MLP}}^{(r_0)}-f_{r_0}^\star
        \|_{L^2(P_x)}^2
    \right]
    &\leq
    \inf_{f\in\mathcal F_{B,m}}
    \|f-f_{r_0}^\star\|_{L^2(P_x)}^2\\
    &\quad+
    C\,\mathbb E\left[
        {B(BM_n+G_0)\over n}
        \left(\sum_{i=1}^n\|x_i\|_2^2\right)^{1/2}
    \right].
\end{align*}
Under the multi-index distribution,
\[
    \|x_i\|_2^2=R^2+\|Z_i\|_2^2,
    \qquad
    Z_i\sim\mathcal N(0,I_{Kr_0}).
\]
Hence
\[
    \mathbb E\sum_{i=1}^n\|x_i\|_2^2=n(Kr_0+R^2),
\]
and Gaussian norm concentration gives
\[
    \mathbb E[M_n^2]
    \leq C\bigl(Kr_0+R^2+\log(en)\bigr).
\]
The result follows from Cauchy--Schwarz.
\end{proof}

\begin{theorem}[Specialized multi-index MLP approximation and consistency]
\label{thm:multi-index-mlp-consistency}
Fix $r_0$, $R$, and a bounded Lipschitz function
$g:\mathbb R^{r_0}\to\mathbb R$ as above. There is a sequence of
cluster-specialized two-layer ReLU networks $(F_K)$ satisfying
\[
    m(F_K)\leq \left\lceil K^{9/8}\right\rceil,
    \qquad
    \|\theta(F_K)\|_{\mathrm F}^2\leq K^{9/8},
\]
and
\[
    \|F_K-f_{r_0}^\star\|_{L^2(P_x)}^2\longrightarrow0.
\]
Moreover, if
\[
    m_K=\left\lceil K^{9/8}\right\rceil,
    \qquad
    B_K=K^{9/8},
    \qquad
    n_K=\left\lceil K^7\right\rceil,
\]
then any empirical squared-loss minimizer over
$\mathcal F_{B_K,m_K}$ satisfies
\[
    \lim_{K\to\infty}
    \mathbb E\left[
        \|\widehat f_{\mathrm{MLP}}^{(r_0)}-f_{r_0}^\star
        \|_{L^2(P_x)}^2
    \right]
    =0.
\]
\end{theorem}

\begin{proof}
By Lemma~\ref{lem:ridge-ramp-density}, choose ridge-ramp networks
$(p_s)_{s\geq1}$ such that
\[
    \|p_s-g\|_{L^2(\gamma_{r_0})}\longrightarrow0.
\]
Write $J_s$ for the number of ridge-ramp units in $p_s$, and define
\[
    D_s:=\max\left\{2J_s+1,\mathcal C_R(p_s),1\right\},
    \qquad
    A_s:=\max_{1\leq j\leq s}D_j.
\]
Let $p_0=0$ and set
\[
    s_K
    :=\max\left(
        \{0\}\cup
        \{s\leq K:A_s\leq K^{1/8}\}
      \right).
\]
Every $A_s$ is finite, so $s_K\to\infty$.

Apply Lemma~\ref{lem:routed-multi-index-approximant} to $p_{s_K}$. The
resulting network $F_K$ has width at most
\[
    K(2J_{s_K}+1)\leq K^{9/8}
\]
and admits a parameterization satisfying
\[
    \|\theta(F_K)\|_{\mathrm F}^2
    \leq K\mathcal C_R(p_{s_K})
    \leq K^{9/8}.
\]
It is cluster-specialized in the sense of
Lemma~\ref{lem:routed-multi-index-approximant}, and
\[
    \|F_K-f_{r_0}^\star\|_{L^2(P_x)}^2
    =\|p_{s_K}-g\|_{L^2(\gamma_{r_0})}^2
    \longrightarrow0.
\]

Applying Lemma~\ref{lem:path-rademacher} with
$B_K=K^{9/8}$ and $n_K=\lceil K^7\rceil$ gives
\begin{align*}
    \mathbb E\left[
        \|\widehat f_{\mathrm{MLP}}^{(r_0)}-f_{r_0}^\star
        \|_{L^2(P_x)}^2
    \right]
    &\leq
    \|F_K-f_{r_0}^\star\|_{L^2(P_x)}^2\\
    &\quad+
    CB_K^2\sqrt{
        {(Kr_0+R^2)(Kr_0+R^2+\log(en_K))\over n_K}
    }\\
    &\quad+
    CB_KG_0\sqrt{{Kr_0+R^2\over n_K}}.
\end{align*}
Because $r_0$ and $R$ are fixed, the two estimation terms are respectively
$O(K^{-1/4})$ and $O(K^{-15/8})$. Both vanish, proving consistency.
\end{proof}

\begin{corollary}[Single-index specialized MLP upper bound]
\label{cor:single-index-mlp-consistency}
Under \cref{data:rfm-mlp-gap}, the target $f^\star$ admits approximation by
cluster-specialized two-layer ReLU networks with polynomially growing width
and squared Frobenius norm and with approximation error tending to zero.
Moreover, the MLP estimator with
\[
    m_K=\left\lceil K^{9/8}\right\rceil,
    \qquad
    B_K=K^{9/8},
    \qquad
    n_K=\left\lceil K^7\right\rceil
\]
is consistent.
\end{corollary}

\begin{proof}
Take $r_0=1$ and $R=1$ in
Theorem~\ref{thm:multi-index-mlp-consistency}. Then $V_c$ reduces to the
single predictive direction $v_c$, and the multi-index distribution reduces
to \cref{data:rfm-mlp-gap}.
\end{proof}

\subsubsection{The Single-Index Kernel and RFM Lower Bounds}

For the remainder of the proof, we return to the single-index setting.
Throughout, $d=2K$ and $R\geq1$ is fixed. We use
\[
    \mu_c=[Re_c;0],
    \qquad
    v_c=[0;e_c],
    \qquad
    \Sigma_c=\operatorname{diag}(0_{K\times K},I_K).
\]
Equivalently,
\[
    c\sim\operatorname{Unif}([K]),
    \qquad
    x=\mu_c+
    \begin{bmatrix}
        0\\
        z
    \end{bmatrix},
    \qquad
    z\sim\mathcal N(0,I_K),
\]
and
\[
    f^\star(x)=g(\langle x,v_c\rangle),
\]
where $g:\mathbb R\to\mathbb R$ is fixed, nonconstant, bounded, and
Lipschitz.



\subsubsection{The Ground-Truth AGOP Removes the Routing Block}


Define the ground-truth population AGOP across all input
coordinates by
\[
    M
    :=\mathbb E\left[
        \nabla_x f^\star(x)\nabla_x f^\star(x)^\top
      \right].
\]

\begin{lemma}[The AGOP kernel loses the cluster identity]
\label{lem:ground-truth-agop-loses-routing}
Let $P_V:=\sum_{c=1}^K v_cv_c^\top$ be the orthogonal projector onto the
predictive subspace, let $G\sim\mathcal N(0,1)$, and set
\[
    \alpha:=\mathbb E[g'(G)^2].
\]
Then $0<\alpha\leq L^2$, and
\[
    M=\frac{\alpha}{K}P_V.
\]
Moreover, every measurable function $h$ of $\sqrt Mx$ satisfies
\[
    \mathbb E\left[\bigl(h(\sqrt Mx)-f^\star(x)\bigr)^2\right]
    \geq
    \left(1-\frac1K\right)\operatorname{Var}(g(G)).
\]
\end{lemma}

\begin{proof}
On the support of cluster $c$, the regression function is
$f^\star(x)=g(\langle x,v_c\rangle)$ and hence, almost everywhere,
\[
    \nabla_x f^\star(x)
    =g'(\langle x,v_c\rangle)v_c.
\]

Since $g$ is Lipschitz, $|g'|\leq L$ almost everywhere, so
$\alpha\leq L^2$. If $\alpha=0$, then $g'=0$ almost everywhere with respect
to Gaussian measure and therefore almost everywhere with respect to Lebesgue
measure. Since a Lipschitz function is absolutely continuous, this would make
$g$ constant. Hence $\alpha>0$.

Averaging the corresponding outer product over the input and the uniform
cluster index gives
\[
    M
    =\frac1K\sum_{c=1}^K
      \mathbb E\left[g'(\langle x,v_c\rangle)^2\mid c\right]v_cv_c^\top
    =\frac{\alpha}{K}\sum_{c=1}^K v_cv_c^\top
    =\frac{\alpha}{K}P_V.
\]
Therefore $\sqrt M=\sqrt{\alpha/K}\,P_V$. In particular, $\sqrt Mx$
contains the predictive coordinates
$S_j:=\langle x,v_j\rangle$, $j\in[K]$, but none of the routing coordinates.
Under the data distribution, $S=(S_1,\ldots,S_K)\sim\mathcal N(0,I_K)$ and is
independent of the uniform cluster index $c$, while
$f^\star(x)=g(S_c)$. Thus, conditional on $\sqrt Mx$,
\[
    \mathbb E[f^\star(x)\mid\sqrt Mx]
    =\frac1K\sum_{j=1}^K g(S_j).
\]
The conditional-expectation characterization of squared-loss regression
implies that every $h(\sqrt Mx)$ has risk at least
\begin{align*}
    \mathbb E\!\left[\operatorname{Var}
        \bigl(f^\star(x)\mid\sqrt Mx\bigr)\right]
    &=\mathbb E\!\left[
        \frac1K\sum_{j=1}^K g(S_j)^2
        -\left(\frac1K\sum_{j=1}^K g(S_j)\right)^2
      \right]\\
    &=\left(1-\frac1K\right)\operatorname{Var}(g(G)).
\end{align*}

Because $g$ is continuous and nonconstant and the Gaussian distribution has
full support, $\operatorname{Var}(g(G))>0$.

\end{proof}

\subsubsection{An All-Orders Lower Bound for Kernel Methods and Regularized RFM}

We next study the full family in \cref{train:rfm-gap}. Let
$P_R:=I_d-P_V$ denote the projector onto the routing subspace. By
Lemma~\ref{lem:ground-truth-agop-loses-routing},
\[
    M_\rho
    =\rho P_R+\left(\rho+\frac{\alpha}{K}\right)P_V.
\]
Thus both the standard metric $I_d$ and the regularized RFM metric $M_\rho$
commute with every rotation of the predictive subspace. The lower bound below
is uniform over $\rho\geq0$, applies to any rotationally invariant base kernel
$\mathcal K$, and is independent of the empirical ridge parameter. 

This degree-by-degree rotational-invariance obstruction is closely related to
high-dimensional results showing that rotationally invariant kernel methods
with polynomially many samples recover only bounded-degree components
\citep{ghorbani2021linearized,ghorbani2020outperform}.
Related group-based analyses of invariant kernel and random-feature models
appear in \citet{mei2021invariances}. More broadly,
\citet{abbe2022nonuniversality} develop symmetry-based limitations for
equivariant learning algorithms.


Let $\operatorname{He}_r$ denote the probabilists' Hermite polynomial and
$h_r=\operatorname{He}_r/\sqrt{r!}$ its normalized version. Write
\[
    g(t)=\sum_{r=0}^\infty\widehat g_r h_r(t),
    \qquad
    \widehat g_r=\mathbb E[g(G)h_r(G)].
\]


\begin{lemma}[Bounded nonconstant links have infinite Hermite support]
\label{lem:infinite-hermite-support}
\label{lem:clipped-ramp-hermite}
The Hermite coefficients of $g$ are not eventually zero. Equivalently,
\[
    \sum_{r>r_0}\widehat g_r^{\,2}>0
    \qquad\text{for every finite }r_0.
\]
\end{lemma}

\begin{proof}
Suppose instead that, for some finite $r_0$,
\[
    \widehat g_r=0
    \qquad\text{for all }r>r_0.
\]
Then, in $L^2(\mathcal N(0,1))$,
\[
    g(t)=\sum_{r=0}^{r_0}\widehat g_rh_r(t)=:p(t),
\]
where $p$ is a polynomial. Both $g$ and $p$ are continuous, and the Gaussian
distribution has full support. Hence equality almost everywhere implies
$g(t)=p(t)$ for every $t\in\mathbb R$. But $g$ is bounded, so $p$ is a
bounded polynomial and must be constant, contradicting the assumption that
$g$ is nonconstant. Parseval's identity then gives the stated positive-tail
property.
\end{proof}


For $r\geq2$, let $\mathcal V_{r,K}$ be the traceless order-$r$ Gaussian-chaos
space. Equivalently, under the standard isometry between the $r$th Gaussian
chaos and symmetric order-$r$ tensors, $\mathcal V_{r,K}$ corresponds to the
kernel of the tensor trace map. It is the irreducible $O(K)$ representation
of harmonic degree-$r$ polynomials.

\begin{lemma}[Harmonic dimension and directional energy]
\label{lem:harmonic-order-r}
The dimension of $\mathcal V_{r,K}$ is
\[
    H_{r,K}
    =\binom{K+r-1}{r}-\binom{K+r-3}{r-2}.
\]
Let $P_{r,K}$ be orthogonal projection onto $\mathcal V_{r,K}$. For every
unit vector $q\in\mathbb R^K$,
\[
    \left\|P_{r,K}h_r(q^\top z)\right\|_{L^2}^2
    =\alpha_{r,K},
\]
where, writing $r=2m$ or $r=2m+1$,
\[
    \alpha_{2m,K}
    =\prod_{j=0}^{m-1}
      \frac{K+2j-1}{K+2m+2j-2},
\]
and
\[
    \alpha_{2m+1,K}
    =\prod_{j=0}^{m-1}
      \frac{K+2j-1}{K+2m+2j}.
\]
For every fixed $r$, $H_{r,K}=\Theta_r(K^r)$ and
$\alpha_{r,K}\to1$ as $K\to\infty$.
\end{lemma}

\begin{proof}
The dimension of symmetric order-$r$ tensors is
$\binom{K+r-1}{r}$. The trace map onto symmetric order-$(r-2)$ tensors is
surjective, and its kernel is the traceless subspace. This gives the stated
dimension.

Under the Gaussian-chaos/tensor isometry,
$h_r(q^\top z)$ corresponds to $q^{\otimes r}$. By rotational
invariance, take $q=e_1$. The orthogonal projection of the homogeneous
polynomial $x_1^r$ onto harmonic degree $r$ is
\[
    \mathcal H_r[x_1^r]
    =\sum_{j=0}^{\lfloor r/2\rfloor}
      \frac{(-1)^j\Gamma(r-j+K/2-1)}
           {4^j j!\,\Gamma(r+K/2-1)}
      \|x\|^{2j}\Delta^j x_1^r.
\]
Since
\[
    \Delta^j x_1^r=\frac{r!}{(r-2j)!}x_1^{r-2j},
\]
the coefficient of $x_1^r$ in this harmonic projection is
\[
    \sum_{j=0}^{\lfloor r/2\rfloor}
      \frac{(-1)^j r!\,\Gamma(r-j+K/2-1)}
           {4^j j!(r-2j)!\,\Gamma(r+K/2-1)}.
\]
In the symmetric-tensor inner product, this coefficient equals
$\langle e_1^{\otimes r},P_{r,K}e_1^{\otimes r}\rangle$, which equals the
squared norm of the projection. Simplifying the finite sum gives the two
product formulas displayed above. The asymptotic statements follow
immediately from those formulas and the dimension expression.
\end{proof}

\begin{lemma}[Invariant random-subspace bound]
\label{lem:random-subspace-bound}
Let $V$ be a $D$-dimensional irreducible orthogonal representation of a compact
group. Let $S\subseteq V$ be a random subspace with invariant law and
$\dim S\leq N$ almost surely. Then, for every fixed $v\in V$,
\[
    \mathbb E\operatorname{dist}(v,S)^2
    \geq\left(1-\frac ND\right)_+\|v\|_2^2.
\]
\end{lemma}

\begin{proof}
Let $\Pi_S$ be orthogonal projection onto $S$. Invariance implies that
$\mathbb E\Pi_S$ commutes with the group action. Schur's lemma gives
$\mathbb E\Pi_S=\beta I_V$. Taking traces yields
$\beta D=\mathbb E\dim S\leq N$. Therefore
\[
    \mathbb E\operatorname{dist}(v,S)^2
    =\|v\|_2^2-\mathbb E\|\Pi_Sv\|_2^2
    =(1-\beta)\|v\|_2^2
    \geq\left(1-\frac ND\right)_+\|v\|_2^2.
\]
\end{proof}

\begin{lemma}[All-orders lower bound for regularized AGOP kernel methods]
\label{lem:all-orders-krr-lower-bound}
For every rotationally invariant kernel $\mathcal K$, every sample size $n$,
every $\rho\geq0$, and every empirical ridge parameter $\lambda_n\geq0$,
both estimators from \cref{train:rfm-gap} satisfy
\[
    \mathbb E\left[
        \|\widehat f-f^\star\|_{L^2(P_x)}^2
    \right]
    \geq
    \sum_{r=2}^\infty
    \widehat g_r^{\,2}\alpha_{r,K}
    \left(1-\frac{n}{H_{r,K}}\right)_+,
    \qquad
    \widehat f\in
    \left\{\widehat f_{\mathrm{Kernel}},\widehat f_{\mathrm{RFM}}\right\}.
\]
\end{lemma}

\begin{proof}
Let $A=I_d$ for $\widehat f_{\mathrm{Kernel}}$ and $A=M_\rho$ for
$\widehat f_{\mathrm{RFM}}$.
By the representer theorem, for a fixed training sample the predictor belongs
to the span of the $n$ transformed-kernel sections centered at the training
inputs. For a test point in cluster $c$, restrict these sections to that
cluster's support and write them as functions of its predictive argument
$\xi\in\mathbb R^K$:
\[
    \psi_{c,i}^{A}(\xi)
    :=\mathcal K\!\left(
        \sqrt A x_i,
        \sqrt A
        \left(\mu_c+\begin{bmatrix}0\\\xi\end{bmatrix}\right)
      \right),
    \qquad 1\leq i\leq n.
\]
At Hermite order $r\geq2$, define
\[
    S_{c,r}
    :=\operatorname{span}\left\{
        P_{r,K}\psi_{c,i}^{A}:1\leq i\leq n
      \right\}
    \subseteq\mathcal V_{r,K}.
\]
Then $\dim S_{c,r}\leq n$. The predictive coordinates of the training inputs
are standard Gaussian, $A$ commutes with predictive rotations, and
$\mathcal K$ is rotationally invariant. Hence the law of $S_{c,r}$ is
invariant under the $O(K)$ action on $\mathcal V_{r,K}$ for both estimators
and every $\rho\geq0$.

The order-$r$ target component in cluster $c$ is
\[
    \gamma_{c,r}
    =\widehat g_rP_{r,K}h_r(e_c^\top\xi),
\]
whose squared norm is $\widehat g_r^{\,2}\alpha_{r,K}$ by
Lemma~\ref{lem:harmonic-order-r}. The projected predictor lies in $S_{c,r}$.
Since distinct Gaussian-chaos orders are orthogonal, for every finite $L_0$,
\[
    \|\widehat f_c-g(e_c^\top\!\cdot)\|_{L^2}^2
    \geq
    \sum_{r=2}^{L_0}\operatorname{dist}(\gamma_{c,r},S_{c,r})^2.
\]
Apply Lemma~\ref{lem:random-subspace-bound} and then let $L_0\to\infty$ by
monotone convergence. This gives, for every cluster $c$,
\[
    \mathbb E\left[
        \|\widehat f_c-g(e_c^\top\!\cdot)\|_{L^2}^2
    \right]
    \geq
    \sum_{r=2}^\infty
    \widehat g_r^{\,2}\alpha_{r,K}
    \left(1-\frac{n}{H_{r,K}}\right)_+.
\]
Averaging over the uniform test cluster proves the result.
\end{proof}

\begin{lemma}[Polynomial sample budgets leave a nonzero kernel/RFM error]
\label{lem:polynomial-kernel-rfm-failure}
Fix $A<\infty$, let $n_K=O(K^A)$, and let $\rho_K\geq0$ be any sequence.
Then, for each
\[
    \widehat f\in
    \left\{\widehat f_{\mathrm{Kernel}},\widehat f_{\mathrm{RFM}}\right\},
\]
\[
    \liminf_{K\to\infty}
    \mathbb E\left[
        \|\widehat f-f^\star\|_{L^2(P_x)}^2
    \right]
    \geq
    \sum_{\substack{r\geq2\\r>A}}\widehat g_r^{\,2}>0.
\]
\end{lemma}

\begin{proof}
Fix an order $r>A$. By Lemma~\ref{lem:harmonic-order-r},
$H_{r,K}=\Theta_r(K^r)$ and $\alpha_{r,K}\to1$, so
\[
    \frac{n_K}{H_{r,K}}\longrightarrow0.
\]
Apply Lemma~\ref{lem:all-orders-krr-lower-bound}. For any finite collection
of orders satisfying $r>A$, each corresponding summand converges to
$\widehat g_r^{\,2}$. Taking the lower limit and then increasing the finite
collection gives, by monotone convergence,
\[
    \liminf_{K\to\infty}
    \mathbb E\left[
        \|\widehat f-f^\star\|_{L^2(P_x)}^2
    \right]
    \geq
    \sum_{\substack{r\geq2\\r>A}}\widehat g_r^{\,2}.
\]

The right-hand side is strictly positive by
Lemma~\ref{lem:infinite-hermite-support}.

\end{proof}


\begin{proof}[Proof of Theorem~\ref{thm:rfm-mlp-gap}]
The specialized MLP approximation and consistency conclusions follow from
Corollary~\ref{cor:single-index-mlp-consistency}. In particular, the choices
\[
    m_K=\left\lceil K^{9/8}\right\rceil,
    \qquad
    B_K=K^{9/8},
    \qquad
    n_K=\left\lceil K^7\right\rceil
\]
are polynomial in $K$ and give
\[
    \mathbb E\left[
        \|\widehat f_{\mathrm{MLP}}-f^\star\|_{L^2(P_x)}^2
    \right]
    \longrightarrow0.
\]

For the lower bound, fix any $A<\infty$, any sequence $n_K=O(K^A)$, and any
sequence $\rho_K\geq0$. Lemma~\ref{lem:polynomial-kernel-rfm-failure} gives
\[
    \liminf_{K\to\infty}
    \mathbb E\left[
        \|\widehat f-f^\star\|_{L^2(P_x)}^2
    \right]
    \geq
    \sum_{\substack{r\geq2\\r>A}}\widehat g_r^{\,2}>0,
    \qquad
    \widehat f\in
    \left\{\widehat f_{\mathrm{Kernel}},\widehat f_{\mathrm{RFM}}\right\}.
\]
The lower bound is uniform over $\rho_K$ and over the empirical ridge
parameters allowed in \cref{train:rfm-gap}. This proves the theorem.
\end{proof}


\end{document}
